%% Generated by Sphinx.
\def\sphinxdocclass{report}
\documentclass[letterpaper,10pt,english]{sphinxmanual}
\ifdefined\pdfpxdimen
   \let\sphinxpxdimen\pdfpxdimen\else\newdimen\sphinxpxdimen
\fi \sphinxpxdimen=.75bp\relax
\ifdefined\pdfimageresolution
    \pdfimageresolution= \numexpr \dimexpr1in\relax/\sphinxpxdimen\relax
\fi
\newdimen\sphinxremdimen\sphinxremdimen = 10pt
%% let collapsible pdf bookmarks panel have high depth per default
\PassOptionsToPackage{bookmarksdepth=5}{hyperref}

\PassOptionsToPackage{booktabs}{sphinx}
\PassOptionsToPackage{colorrows}{sphinx}

\PassOptionsToPackage{warn}{textcomp}
\usepackage[utf8]{inputenc}
\ifdefined\DeclareUnicodeCharacter
% support both utf8 and utf8x syntaxes
  \ifdefined\DeclareUnicodeCharacterAsOptional
    \def\sphinxDUC#1{\DeclareUnicodeCharacter{"#1}}
  \else
    \let\sphinxDUC\DeclareUnicodeCharacter
  \fi
  \sphinxDUC{00A0}{\nobreakspace}
  \sphinxDUC{2500}{\sphinxunichar{2500}}
  \sphinxDUC{2502}{\sphinxunichar{2502}}
  \sphinxDUC{2514}{\sphinxunichar{2514}}
  \sphinxDUC{251C}{\sphinxunichar{251C}}
  \sphinxDUC{2572}{\textbackslash}
\fi
\usepackage{cmap}
\usepackage[T1]{fontenc}
\usepackage{amsmath,amssymb,amstext}
\usepackage{babel}



\usepackage{tgtermes}
\usepackage{tgheros}
\renewcommand{\ttdefault}{txtt}



\usepackage[Bjarne]{fncychap}
\usepackage{sphinx}

\fvset{fontsize=auto}
\usepackage{geometry}


% Include hyperref last.
\usepackage{hyperref}
% Fix anchor placement for figures with captions.
\usepackage{hypcap}% it must be loaded after hyperref.
% Set up styles of URL: it should be placed after hyperref.
\urlstyle{same}

\addto\captionsenglish{\renewcommand{\contentsname}{Contents:}}

\usepackage{sphinxmessages}
\setcounter{tocdepth}{1}



\title{PreencheFacil}
\date{Mar 13, 2026}
\release{v0.0.3}
\author{Avance Telecom}
\newcommand{\sphinxlogo}{\vbox{}}
\renewcommand{\releasename}{Release}
\makeindex
\begin{document}

\ifdefined\shorthandoff
  \ifnum\catcode`\=\string=\active\shorthandoff{=}\fi
  \ifnum\catcode`\"=\active\shorthandoff{"}\fi
\fi

\pagestyle{empty}
\sphinxmaketitle
\pagestyle{plain}
\sphinxtableofcontents
\pagestyle{normal}
\phantomsection\label{\detokenize{index::doc}}


\sphinxAtStartPar
Add your content using \sphinxcode{\sphinxupquote{reStructuredText}} syntax. See the
\sphinxhref{https://www.sphinx-doc.org/en/master/usage/restructuredtext/index.html}{reStructuredText}
documentation for details.

\sphinxstepscope


\chapter{text\_extractor}
\label{\detokenize{api/modules:text-extractor}}\label{\detokenize{api/modules::doc}}
\sphinxstepscope


\section{app module}
\label{\detokenize{api/app:module-app}}\label{\detokenize{api/app:app-module}}\label{\detokenize{api/app::doc}}\index{module@\spxentry{module}!app@\spxentry{app}}\index{app@\spxentry{app}!module@\spxentry{module}}\index{abs\_path() (in module app)@\spxentry{abs\_path()}\spxextra{in module app}}

\begin{fulllineitems}
\phantomsection\label{\detokenize{api/app:app.abs_path}}
\pysigstartsignatures
\pysiglinewithargsret
{\sphinxcode{\sphinxupquote{app.}}\sphinxbfcode{\sphinxupquote{abs\_path}}}
{\sphinxparam{\DUrole{n}{p}}}
{}
\pysigstopsignatures
\sphinxAtStartPar
Resolve um caminho de arquivo para um caminho absoluto.

\sphinxAtStartPar
Esta função é crucial para lidar com recursos em um aplicativo que pode ser
executado tanto em um ambiente de desenvolvimento quanto como um executável
“congelado” (por exemplo, com PyInstaller).
\begin{itemize}
\item {} 
\sphinxAtStartPar
Se o caminho já for absoluto, ele é retornado sem modificação.

\item {} \begin{description}
\sphinxlineitem{Se o caminho for relativo, a função determina o diretório base:}\begin{itemize}
\item {} 
\sphinxAtStartPar
Em um executável PyInstaller, usa o diretório temporário \sphinxtitleref{\_MEIPASS}.

\item {} 
\sphinxAtStartPar
Em um script Python normal, usa o diretório do próprio script.

\end{itemize}

\end{description}

\item {} 
\sphinxAtStartPar
O caminho relativo é então combinado com este diretório base.

\end{itemize}
\begin{quote}\begin{description}
\sphinxlineitem{Parameters}
\sphinxAtStartPar
\sphinxstyleliteralstrong{\sphinxupquote{p}} (\DUrole{sphinx_autodoc_typehints-type}{\sphinxcode{\sphinxupquote{str}}}) \textendash{} O caminho do arquivo a ser resolvido.

\sphinxlineitem{Returns}
\sphinxAtStartPar
O caminho absoluto correspondente, ou uma string vazia se a entrada for nula/vazia.

\sphinxlineitem{Return type}
\sphinxAtStartPar
\DUrole{sphinx_autodoc_typehints-type}{\sphinxcode{\sphinxupquote{str}}}

\end{description}\end{quote}

\end{fulllineitems}

\index{file\_exists() (in module app)@\spxentry{file\_exists()}\spxextra{in module app}}

\begin{fulllineitems}
\phantomsection\label{\detokenize{api/app:app.file_exists}}
\pysigstartsignatures
\pysiglinewithargsret
{\sphinxcode{\sphinxupquote{app.}}\sphinxbfcode{\sphinxupquote{file\_exists}}}
{\sphinxparam{\DUrole{n}{p}}}
{}
\pysigstopsignatures
\sphinxAtStartPar
Verifica se um arquivo existe e pode ser lido.
\begin{quote}\begin{description}
\sphinxlineitem{Parameters}
\sphinxAtStartPar
\sphinxstyleliteralstrong{\sphinxupquote{p}} (\DUrole{sphinx_autodoc_typehints-type}{\sphinxcode{\sphinxupquote{str}}}) \textendash{} String que representa o caminho para o arquivo a ser verificado.

\sphinxlineitem{Returns}
\sphinxAtStartPar
True se o arquivo existe e pode ser lido, False caso contrário.

\sphinxlineitem{Return type}
\sphinxAtStartPar
\DUrole{sphinx_autodoc_typehints-type}{\sphinxcode{\sphinxupquote{bool}}}

\end{description}\end{quote}

\end{fulllineitems}

\index{SettingsDialog (class in app)@\spxentry{SettingsDialog}\spxextra{class in app}}

\begin{fulllineitems}
\phantomsection\label{\detokenize{api/app:app.SettingsDialog}}
\pysigstartsignatures
\pysiglinewithargsret
{\sphinxbfcode{\sphinxupquote{\DUrole{k}{class}\DUrole{w}{ }}}\sphinxcode{\sphinxupquote{app.}}\sphinxbfcode{\sphinxupquote{SettingsDialog}}}
{\sphinxparam{\DUrole{n}{parent}}}
{}
\pysigstopsignatures
\sphinxAtStartPar
Bases: \sphinxcode{\sphinxupquote{QDialog}}

\sphinxAtStartPar
Classe de diálogo para gerenciar as configurações globais do aplicativo, permitindo ao usuário acessar os ajustes de tema visual e a lista de domínios de e\sphinxhyphen{}mail sugeridos.
\begin{quote}\begin{description}
\sphinxlineitem{Parameters}
\sphinxAtStartPar
\sphinxstyleliteralstrong{\sphinxupquote{QDialog}} (\sphinxstyleliteralemphasis{\sphinxupquote{\_type\_}}) \textendash{} Objeto do tipo QDialog.

\end{description}\end{quote}

\end{fulllineitems}

\index{App (class in app)@\spxentry{App}\spxextra{class in app}}

\begin{fulllineitems}
\phantomsection\label{\detokenize{api/app:app.App}}
\pysigstartsignatures
\pysiglinewithargsret
{\sphinxbfcode{\sphinxupquote{\DUrole{k}{class}\DUrole{w}{ }}}\sphinxcode{\sphinxupquote{app.}}\sphinxbfcode{\sphinxupquote{App}}}
{\sphinxparam{\DUrole{o}{*}\DUrole{n}{args}\DUrole{p}{:}\DUrole{w}{ }\DUrole{n}{Any}}\sphinxparamcomma \sphinxparam{\DUrole{o}{**}\DUrole{n}{kwargs}\DUrole{p}{:}\DUrole{w}{ }\DUrole{n}{Any}}}
{}
\pysigstopsignatures
\sphinxAtStartPar
Bases: \sphinxcode{\sphinxupquote{QWidget}}

\sphinxAtStartPar
Interface principal do aplicativo “Preenche Fácil”.
Gerencia a configuração de campos, a interface de usuário dinâmica,
a integração com planilhas Excel e a persistência de dados.
\begin{quote}\begin{description}
\sphinxlineitem{Parameters}
\sphinxAtStartPar
\sphinxstyleliteralstrong{\sphinxupquote{QWidget}} (\sphinxstyleliteralemphasis{\sphinxupquote{\_type\_}}) \textendash{} Objeto do tipo QWidget.

\end{description}\end{quote}
\index{open\_license() (app.App method)@\spxentry{open\_license()}\spxextra{app.App method}}

\begin{fulllineitems}
\phantomsection\label{\detokenize{api/app:app.App.open_license}}
\pysigstartsignatures
\pysiglinewithargsret
{\sphinxbfcode{\sphinxupquote{open\_license}}}
{}
{}
\pysigstopsignatures
\sphinxAtStartPar
Abre o diálogo de informações da licença, exibindo o texto da licença
e aplicando o tema visual atual do aplicativo.
\begin{quote}\begin{description}
\sphinxlineitem{Return type}
\sphinxAtStartPar
\DUrole{sphinx_autodoc_typehints-type}{\sphinxcode{\sphinxupquote{None}}}

\end{description}\end{quote}

\end{fulllineitems}

\index{open\_email\_domains() (app.App method)@\spxentry{open\_email\_domains()}\spxextra{app.App method}}

\begin{fulllineitems}
\phantomsection\label{\detokenize{api/app:app.App.open_email_domains}}
\pysigstartsignatures
\pysiglinewithargsret
{\sphinxbfcode{\sphinxupquote{open\_email\_domains}}}
{}
{}
\pysigstopsignatures
\sphinxAtStartPar
Abre o diálogo de gerenciamento de domínios de e\sphinxhyphen{}mail.
Permite adicionar, editar ou remover domínios da lista de sugestões.
Ao aceitar, salva as alterações e atualiza todos os widgets de e\sphinxhyphen{}mail ativos.
\begin{quote}\begin{description}
\sphinxlineitem{Return type}
\sphinxAtStartPar
\DUrole{sphinx_autodoc_typehints-type}{\sphinxcode{\sphinxupquote{None}}}

\end{description}\end{quote}

\end{fulllineitems}

\index{clear\_selected\_sheet() (app.App method)@\spxentry{clear\_selected\_sheet()}\spxextra{app.App method}}

\begin{fulllineitems}
\phantomsection\label{\detokenize{api/app:app.App.clear_selected_sheet}}
\pysigstartsignatures
\pysiglinewithargsret
{\sphinxbfcode{\sphinxupquote{clear\_selected\_sheet}}}
{}
{}
\pysigstopsignatures
\sphinxAtStartPar
Limpa a seleção da planilha atual sem apagar o arquivo do disco.

\sphinxAtStartPar
Remove o vínculo da interface com a planilha, limpa a lista de abas,
atualiza a configuração persistida e mantém os campos do controle na UI.
\begin{quote}\begin{description}
\sphinxlineitem{Return type}
\sphinxAtStartPar
\DUrole{sphinx_autodoc_typehints-type}{\sphinxcode{\sphinxupquote{None}}}

\end{description}\end{quote}

\end{fulllineitems}

\index{on\_sheet\_changed() (app.App method)@\spxentry{on\_sheet\_changed()}\spxextra{app.App method}}

\begin{fulllineitems}
\phantomsection\label{\detokenize{api/app:app.App.on_sheet_changed}}
\pysigstartsignatures
\pysiglinewithargsret
{\sphinxbfcode{\sphinxupquote{on\_sheet\_changed}}}
{\sphinxparam{\DUrole{n}{new\_sheet}}}
{}
\pysigstopsignatures
\sphinxAtStartPar
Atualiza a aba ativa e recarrega os campos correspondentes. Se a aba for nova,
tenta sincronizar os cabeçalhos a partir do arquivo Excel.
\begin{quote}\begin{description}
\sphinxlineitem{Parameters}
\sphinxAtStartPar
\sphinxstyleliteralstrong{\sphinxupquote{new\_sheet}} (\DUrole{sphinx_autodoc_typehints-type}{\sphinxcode{\sphinxupquote{str}}}) \textendash{} String representando a nova aba

\sphinxlineitem{Return type}
\sphinxAtStartPar
\DUrole{sphinx_autodoc_typehints-type}{\sphinxcode{\sphinxupquote{None}}}

\end{description}\end{quote}

\end{fulllineitems}

\index{create\_new\_sheet() (app.App method)@\spxentry{create\_new\_sheet()}\spxextra{app.App method}}

\begin{fulllineitems}
\phantomsection\label{\detokenize{api/app:app.App.create_new_sheet}}
\pysigstartsignatures
\pysiglinewithargsret
{\sphinxbfcode{\sphinxupquote{create\_new\_sheet}}}
{}
{}
\pysigstopsignatures
\sphinxAtStartPar
Função responsável por criar uma nova aba na planilha
\begin{quote}\begin{description}
\sphinxlineitem{Return type}
\sphinxAtStartPar
\DUrole{sphinx_autodoc_typehints-type}{\sphinxcode{\sphinxupquote{None}}}

\end{description}\end{quote}

\end{fulllineitems}

\index{refresh\_sheets() (app.App method)@\spxentry{refresh\_sheets()}\spxextra{app.App method}}

\begin{fulllineitems}
\phantomsection\label{\detokenize{api/app:app.App.refresh_sheets}}
\pysigstartsignatures
\pysiglinewithargsret
{\sphinxbfcode{\sphinxupquote{refresh\_sheets}}}
{}
{}
\pysigstopsignatures
\sphinxAtStartPar
Atualiza a lista de abas (sheets) disponíveis no arquivo Excel selecionado e
mantém o estado da aba atual consistente com o conteúdo do arquivo.
\begin{quote}\begin{description}
\sphinxlineitem{Return type}
\sphinxAtStartPar
\DUrole{sphinx_autodoc_typehints-type}{\sphinxcode{\sphinxupquote{None}}}

\end{description}\end{quote}

\end{fulllineitems}

\index{open\_settings() (app.App method)@\spxentry{open\_settings()}\spxextra{app.App method}}

\begin{fulllineitems}
\phantomsection\label{\detokenize{api/app:app.App.open_settings}}
\pysigstartsignatures
\pysiglinewithargsret
{\sphinxbfcode{\sphinxupquote{open\_settings}}}
{}
{}
\pysigstopsignatures
\sphinxAtStartPar
Função responsável por abrir o diálogo de configurações
\begin{quote}\begin{description}
\sphinxlineitem{Return type}
\sphinxAtStartPar
\DUrole{sphinx_autodoc_typehints-type}{\sphinxcode{\sphinxupquote{None}}}

\end{description}\end{quote}

\end{fulllineitems}

\index{open\_help() (app.App method)@\spxentry{open\_help()}\spxextra{app.App method}}

\begin{fulllineitems}
\phantomsection\label{\detokenize{api/app:app.App.open_help}}
\pysigstartsignatures
\pysiglinewithargsret
{\sphinxbfcode{\sphinxupquote{open\_help}}}
{}
{}
\pysigstopsignatures
\sphinxAtStartPar
Acessa a página web com informações sobre o software
\begin{quote}\begin{description}
\sphinxlineitem{Return type}
\sphinxAtStartPar
\DUrole{sphinx_autodoc_typehints-type}{\sphinxcode{\sphinxupquote{None}}}

\end{description}\end{quote}

\end{fulllineitems}

\index{render\_fields() (app.App method)@\spxentry{render\_fields()}\spxextra{app.App method}}

\begin{fulllineitems}
\phantomsection\label{\detokenize{api/app:app.App.render_fields}}
\pysigstartsignatures
\pysiglinewithargsret
{\sphinxbfcode{\sphinxupquote{render\_fields}}}
{}
{}
\pysigstopsignatures
\sphinxAtStartPar
Renderiza a lista de campos na interface, recriando as linhas de entrada.

\sphinxAtStartPar
Preserva os valores já digitados (incluindo partes do e\sphinxhyphen{}mail), limpa o layout
atual e monta novamente os widgets conforme o tipo de cada campo (email,
booleano ou texto), conectando sinais de editar/excluir e aplicando ícones e
estado de bloqueio.
\begin{quote}\begin{description}
\sphinxlineitem{Return type}
\sphinxAtStartPar
\DUrole{sphinx_autodoc_typehints-type}{\sphinxcode{\sphinxupquote{None}}}

\sphinxlineitem{Returns}
\sphinxAtStartPar
None

\end{description}\end{quote}

\end{fulllineitems}

\index{sync\_fields\_from\_existing\_excel() (app.App method)@\spxentry{sync\_fields\_from\_existing\_excel()}\spxextra{app.App method}}

\begin{fulllineitems}
\phantomsection\label{\detokenize{api/app:app.App.sync_fields_from_existing_excel}}
\pysigstartsignatures
\pysiglinewithargsret
{\sphinxbfcode{\sphinxupquote{sync\_fields\_from\_existing\_excel}}}
{\sphinxparam{\DUrole{n}{path}}}
{}
\pysigstopsignatures
\sphinxAtStartPar
Sincroniza os campos do controle com o cabeçalho existente na planilha.

\sphinxAtStartPar
Lê os cabeçalhos da aba, cria cabeçalho se estiver ausente, ou então mapeia
cada coluna para um Campo (reaproveitando id/tipo quando possível e inferindo
para novos títulos). Ao final, persiste a configuração, aplica regras de
tipos de coluna e re\sphinxhyphen{}renderiza os campos na UI.
\begin{quote}\begin{description}
\sphinxlineitem{Parameters}
\sphinxAtStartPar
\sphinxstyleliteralstrong{\sphinxupquote{path}} (\DUrole{sphinx_autodoc_typehints-type}{\sphinxcode{\sphinxupquote{str}}}) \textendash{} Caminho do arquivo Excel.

\sphinxlineitem{Returns}
\sphinxAtStartPar
True se a sincronização foi concluída com sucesso; False em caso de erro
ou ausência de cabeçalho válido.

\sphinxlineitem{Return type}
\sphinxAtStartPar
\DUrole{sphinx_autodoc_typehints-type}{\sphinxcode{\sphinxupquote{bool}}}

\end{description}\end{quote}

\end{fulllineitems}

\index{choose\_file() (app.App method)@\spxentry{choose\_file()}\spxextra{app.App method}}

\begin{fulllineitems}
\phantomsection\label{\detokenize{api/app:app.App.choose_file}}
\pysigstartsignatures
\pysiglinewithargsret
{\sphinxbfcode{\sphinxupquote{choose\_file}}}
{}
{}
\pysigstopsignatures
\sphinxAtStartPar
Abre o diálogo para escolher/criar uma planilha Excel e aplica a seleção.

\sphinxAtStartPar
Se o arquivo já existir, pergunta se o usuário deseja importar os campos a
partir do cabeçalho (linha 1). Caso aceite, sincroniza os campos do app com
a planilha; caso contrário, mantém os campos atuais e prepara a planilha.
\begin{quote}\begin{description}
\sphinxlineitem{Return type}
\sphinxAtStartPar
\DUrole{sphinx_autodoc_typehints-type}{\sphinxcode{\sphinxupquote{None}}}

\sphinxlineitem{Returns}
\sphinxAtStartPar
None

\end{description}\end{quote}

\end{fulllineitems}

\index{on\_field\_lock\_toggled() (app.App method)@\spxentry{on\_field\_lock\_toggled()}\spxextra{app.App method}}

\begin{fulllineitems}
\phantomsection\label{\detokenize{api/app:app.App.on_field_lock_toggled}}
\pysigstartsignatures
\pysiglinewithargsret
{\sphinxbfcode{\sphinxupquote{on\_field\_lock\_toggled}}}
{\sphinxparam{\DUrole{n}{field\_id}}\sphinxparamcomma \sphinxparam{\DUrole{n}{locked}}}
{}
\pysigstopsignatures
\sphinxAtStartPar
Atualiza o estado de bloqueio de um campo e persiste a alteração.

\sphinxAtStartPar
Localiza o Campo pelo \sphinxtitleref{field\_id}, ajusta o atributo \sphinxtitleref{locked}, salva a
configuração e re\sphinxhyphen{}renderiza a lista de campos para refletir o novo estado.
\begin{quote}\begin{description}
\sphinxlineitem{Parameters}\begin{itemize}
\item {} 
\sphinxAtStartPar
\sphinxstyleliteralstrong{\sphinxupquote{field\_id}} (\DUrole{sphinx_autodoc_typehints-type}{\sphinxcode{\sphinxupquote{str}}}) \textendash{} ID do campo cujo bloqueio foi alterado.

\item {} 
\sphinxAtStartPar
\sphinxstyleliteralstrong{\sphinxupquote{locked}} (\DUrole{sphinx_autodoc_typehints-type}{\sphinxcode{\sphinxupquote{bool}}}) \textendash{} Novo estado de bloqueio (True = bloqueado).

\end{itemize}

\sphinxlineitem{Return type}
\sphinxAtStartPar
\DUrole{sphinx_autodoc_typehints-type}{\sphinxcode{\sphinxupquote{None}}}

\sphinxlineitem{Returns}
\sphinxAtStartPar
None

\end{description}\end{quote}

\end{fulllineitems}

\index{add\_field() (app.App method)@\spxentry{add\_field()}\spxextra{app.App method}}

\begin{fulllineitems}
\phantomsection\label{\detokenize{api/app:app.App.add_field}}
\pysigstartsignatures
\pysiglinewithargsret
{\sphinxbfcode{\sphinxupquote{add\_field}}}
{}
{}
\pysigstopsignatures
\sphinxAtStartPar
Adiciona um novo campo de preenchimento.

\sphinxAtStartPar
Solicita ao usuário o título e o tipo do campo, valida duplicidade de título,
gera um ID único, salva no controle, re\sphinxhyphen{}renderiza a UI e (se houver planilha
selecionada) tenta garantir/atualizar as colunas no Excel.
\begin{quote}\begin{description}
\sphinxlineitem{Return type}
\sphinxAtStartPar
\DUrole{sphinx_autodoc_typehints-type}{\sphinxcode{\sphinxupquote{None}}}

\sphinxlineitem{Returns}
\sphinxAtStartPar
None

\end{description}\end{quote}

\end{fulllineitems}

\index{edit\_field() (app.App method)@\spxentry{edit\_field()}\spxextra{app.App method}}

\begin{fulllineitems}
\phantomsection\label{\detokenize{api/app:app.App.edit_field}}
\pysigstartsignatures
\pysiglinewithargsret
{\sphinxbfcode{\sphinxupquote{edit\_field}}}
{\sphinxparam{\DUrole{n}{field\_id}}}
{}
\pysigstopsignatures
\sphinxAtStartPar
Edita o título e/ou o tipo de um campo existente.

\sphinxAtStartPar
Impede edição de campos bloqueados, abre diálogos para alterar título e tipo,
valida duplicidade, persiste a mudança e tenta refletir no Excel (renomeando
o cabeçalho e reaplicando regras de tipos). Após re\sphinxhyphen{}renderizar, restaura o
valor previamente digitado, ajustando\sphinxhyphen{}o quando há troca entre tipos com/sem
máscara (telefone/data) e e\sphinxhyphen{}mail.
\begin{quote}\begin{description}
\sphinxlineitem{Parameters}
\sphinxAtStartPar
\sphinxstyleliteralstrong{\sphinxupquote{field\_id}} (\DUrole{sphinx_autodoc_typehints-type}{\sphinxcode{\sphinxupquote{str}}}) \textendash{} ID do campo a ser editado.

\sphinxlineitem{Return type}
\sphinxAtStartPar
\DUrole{sphinx_autodoc_typehints-type}{\sphinxcode{\sphinxupquote{None}}}

\sphinxlineitem{Returns}
\sphinxAtStartPar
None

\end{description}\end{quote}

\end{fulllineitems}

\index{delete\_field() (app.App method)@\spxentry{delete\_field()}\spxextra{app.App method}}

\begin{fulllineitems}
\phantomsection\label{\detokenize{api/app:app.App.delete_field}}
\pysigstartsignatures
\pysiglinewithargsret
{\sphinxbfcode{\sphinxupquote{delete\_field}}}
{\sphinxparam{\DUrole{n}{field\_id}}}
{}
\pysigstopsignatures
\sphinxAtStartPar
Exclui um campo do controle (e a coluna correspondente no Excel, quando possível).

\sphinxAtStartPar
Bloqueia a exclusão de campos protegidos, solicita confirmação do usuário,
remove o campo da lista, persiste a alteração e tenta remover a coluna na
planilha e reaplicar as regras de tipos. Em seguida, re\sphinxhyphen{}renderiza a UI e
atualiza o status.
\begin{quote}\begin{description}
\sphinxlineitem{Parameters}
\sphinxAtStartPar
\sphinxstyleliteralstrong{\sphinxupquote{field\_id}} (\DUrole{sphinx_autodoc_typehints-type}{\sphinxcode{\sphinxupquote{str}}}) \textendash{} ID do campo a ser excluído.

\sphinxlineitem{Return type}
\sphinxAtStartPar
\DUrole{sphinx_autodoc_typehints-type}{\sphinxcode{\sphinxupquote{None}}}

\sphinxlineitem{Returns}
\sphinxAtStartPar
None

\end{description}\end{quote}

\end{fulllineitems}

\index{generate\_protocol() (app.App method)@\spxentry{generate\_protocol()}\spxextra{app.App method}}

\begin{fulllineitems}
\phantomsection\label{\detokenize{api/app:app.App.generate_protocol}}
\pysigstartsignatures
\pysiglinewithargsret
{\sphinxbfcode{\sphinxupquote{generate\_protocol}}}
{}
{}
\pysigstopsignatures
\sphinxAtStartPar
Gera um número de protocolo no formato ddMMyyHHmmss.
Se existir algum campo do tipo ‘protocolo’, preenche automaticamente
o primeiro campo encontrado com o valor gerado.
Não cria campo novo automaticamente.
\begin{quote}\begin{description}
\sphinxlineitem{Return type}
\sphinxAtStartPar
\DUrole{sphinx_autodoc_typehints-type}{\sphinxcode{\sphinxupquote{None}}}

\end{description}\end{quote}

\end{fulllineitems}

\index{delete\_all\_fields() (app.App method)@\spxentry{delete\_all\_fields()}\spxextra{app.App method}}

\begin{fulllineitems}
\phantomsection\label{\detokenize{api/app:app.App.delete_all_fields}}
\pysigstartsignatures
\pysiglinewithargsret
{\sphinxbfcode{\sphinxupquote{delete\_all\_fields}}}
{}
{}
\pysigstopsignatures
\sphinxAtStartPar
Exclui todos os campos não protegidos do controle e, quando possível,
remove também as colunas correspondentes na planilha Excel.

\sphinxAtStartPar
A função:
\sphinxhyphen{} verifica se há campos cadastrados;
\sphinxhyphen{} impede a exclusão de campos protegidos;
\sphinxhyphen{} solicita confirmação do usuário;
\sphinxhyphen{} remove do controle apenas os campos permitidos;
\sphinxhyphen{} persiste a alteração;
\sphinxhyphen{} tenta remover as colunas correspondentes no Excel;
\sphinxhyphen{} reaplica as regras de tipo nas colunas restantes;
\sphinxhyphen{} re\sphinxhyphen{}renderiza a interface e atualiza o status.
\begin{quote}\begin{description}
\sphinxlineitem{Return type}
\sphinxAtStartPar
\DUrole{sphinx_autodoc_typehints-type}{\sphinxcode{\sphinxupquote{None}}}

\sphinxlineitem{Returns}
\sphinxAtStartPar
None

\end{description}\end{quote}

\end{fulllineitems}

\index{clear\_fields() (app.App method)@\spxentry{clear\_fields()}\spxextra{app.App method}}

\begin{fulllineitems}
\phantomsection\label{\detokenize{api/app:app.App.clear_fields}}
\pysigstartsignatures
\pysiglinewithargsret
{\sphinxbfcode{\sphinxupquote{clear\_fields}}}
{}
{}
\pysigstopsignatures
\sphinxAtStartPar
Limpa os valores preenchidos nos campos da interface.

\sphinxAtStartPar
Zera os QLineEdit e, no caso de e\sphinxhyphen{}mail, mantém o domínio e limpa apenas a
parte local. Também atualiza a mensagem de status.
\begin{quote}\begin{description}
\sphinxlineitem{Return type}
\sphinxAtStartPar
\DUrole{sphinx_autodoc_typehints-type}{\sphinxcode{\sphinxupquote{None}}}

\sphinxlineitem{Returns}
\sphinxAtStartPar
None

\end{description}\end{quote}

\end{fulllineitems}

\index{save\_lead() (app.App method)@\spxentry{save\_lead()}\spxextra{app.App method}}

\begin{fulllineitems}
\phantomsection\label{\detokenize{api/app:app.App.save_lead}}
\pysigstartsignatures
\pysiglinewithargsret
{\sphinxbfcode{\sphinxupquote{save\_lead}}}
{}
{}
\pysigstopsignatures
\sphinxAtStartPar
Salva os dados preenchidos como uma nova linha na planilha.

\sphinxAtStartPar
Valida se há arquivo selecionado, verifica possível bloqueio do Excel,
coleta os valores dos campos (tratando e\sphinxhyphen{}mail e booleano), normaliza texto
com máscara e evita salvar quando tudo estiver vazio. Em seguida, tenta
gravar a linha com novas tentativas em caso de PermissionError e, ao
concluir, limpa os campos e informa o usuário.
\begin{quote}\begin{description}
\sphinxlineitem{Return type}
\sphinxAtStartPar
\DUrole{sphinx_autodoc_typehints-type}{\sphinxcode{\sphinxupquote{None}}}

\sphinxlineitem{Returns}
\sphinxAtStartPar
None

\end{description}\end{quote}

\end{fulllineitems}

\index{open\_theme\_settings() (app.App method)@\spxentry{open\_theme\_settings()}\spxextra{app.App method}}

\begin{fulllineitems}
\phantomsection\label{\detokenize{api/app:app.App.open_theme_settings}}
\pysigstartsignatures
\pysiglinewithargsret
{\sphinxbfcode{\sphinxupquote{open\_theme\_settings}}}
{}
{}
\pysigstopsignatures
\sphinxAtStartPar
Abre o diálogo de tema e permite ajustar a cor de fundo (HSL) em tempo real.

\sphinxAtStartPar
Aplica uma prévia do tema conforme os sliders mudam. Se o usuário confirmar,
salva as novas configurações (background\_hsl e theme) e aplica o tema
derivado. Se cancelar, restaura o tema/configuração anterior.
\begin{quote}\begin{description}
\sphinxlineitem{Return type}
\sphinxAtStartPar
\DUrole{sphinx_autodoc_typehints-type}{\sphinxcode{\sphinxupquote{None}}}

\sphinxlineitem{Returns}
\sphinxAtStartPar
None

\end{description}\end{quote}

\end{fulllineitems}


\end{fulllineitems}

\index{main() (in module app)@\spxentry{main()}\spxextra{in module app}}

\begin{fulllineitems}
\phantomsection\label{\detokenize{api/app:app.main}}
\pysigstartsignatures
\pysiglinewithargsret
{\sphinxcode{\sphinxupquote{app.}}\sphinxbfcode{\sphinxupquote{main}}}
{}
{}
\pysigstopsignatures
\sphinxAtStartPar
Ponto de entrada da aplicação.

\sphinxAtStartPar
Inicializa o QApplication e uma tela de carregamento, solicita o e\sphinxhyphen{}mail do
usuário e valida a licença online (com tentativas quando aplicável). Em
seguida, valida/marca o dispositivo localmente e, se tudo estiver ok,
instancia e exibe a janela principal.
\begin{quote}\begin{description}
\sphinxlineitem{Return type}
\sphinxAtStartPar
\DUrole{sphinx_autodoc_typehints-type}{\sphinxcode{\sphinxupquote{None}}}

\sphinxlineitem{Returns}
\sphinxAtStartPar
None

\end{description}\end{quote}

\end{fulllineitems}


\sphinxstepscope


\section{colors module}
\label{\detokenize{api/colors:module-colors}}\label{\detokenize{api/colors:colors-module}}\label{\detokenize{api/colors::doc}}\index{module@\spxentry{module}!colors@\spxentry{colors}}\index{colors@\spxentry{colors}!module@\spxentry{module}}\index{hex\_to\_rgb() (in module colors)@\spxentry{hex\_to\_rgb()}\spxextra{in module colors}}

\begin{fulllineitems}
\phantomsection\label{\detokenize{api/colors:colors.hex_to_rgb}}
\pysigstartsignatures
\pysiglinewithargsret
{\sphinxcode{\sphinxupquote{colors.}}\sphinxbfcode{\sphinxupquote{hex\_to\_rgb}}}
{\sphinxparam{\DUrole{n}{hex\_color}}}
{}
\pysigstopsignatures
\sphinxAtStartPar
Convert a hexadecimal color string to an RGB tuple.

\sphinxAtStartPar
Accepts strings in the form \sphinxcode{\sphinxupquote{\#RRGGBB}} or \sphinxcode{\sphinxupquote{RRGGBB}}. If the input is not a
6\sphinxhyphen{}digit hex color, returns \sphinxcode{\sphinxupquote{(0, 0, 0)}}.
\begin{quote}\begin{description}
\sphinxlineitem{Parameters}
\sphinxAtStartPar
\sphinxstyleliteralstrong{\sphinxupquote{hex\_color}} (\DUrole{sphinx_autodoc_typehints-type}{\sphinxcode{\sphinxupquote{str}}}) \textendash{} Hex color string.

\sphinxlineitem{Return type}
\sphinxAtStartPar
\DUrole{sphinx_autodoc_typehints-type}{\sphinxcode{\sphinxupquote{Tuple}}{[}\sphinxcode{\sphinxupquote{int}}, \sphinxcode{\sphinxupquote{int}}, \sphinxcode{\sphinxupquote{int}}{]}}

\sphinxlineitem{Returns}
\sphinxAtStartPar
A tuple \sphinxcode{\sphinxupquote{(r, g, b)}} with values in the range 0..255.

\end{description}\end{quote}

\end{fulllineitems}

\index{rgb\_to\_hex() (in module colors)@\spxentry{rgb\_to\_hex()}\spxextra{in module colors}}

\begin{fulllineitems}
\phantomsection\label{\detokenize{api/colors:colors.rgb_to_hex}}
\pysigstartsignatures
\pysiglinewithargsret
{\sphinxcode{\sphinxupquote{colors.}}\sphinxbfcode{\sphinxupquote{rgb\_to\_hex}}}
{\sphinxparam{\DUrole{n}{r}}\sphinxparamcomma \sphinxparam{\DUrole{n}{g}}\sphinxparamcomma \sphinxparam{\DUrole{n}{b}}}
{}
\pysigstopsignatures
\sphinxAtStartPar
Convert RGB components to a hexadecimal color string.

\sphinxAtStartPar
Components are clamped to the range 0..255.
\begin{quote}\begin{description}
\sphinxlineitem{Parameters}\begin{itemize}
\item {} 
\sphinxAtStartPar
\sphinxstyleliteralstrong{\sphinxupquote{r}} (\DUrole{sphinx_autodoc_typehints-type}{\sphinxcode{\sphinxupquote{int}}}) \textendash{} Red component (0..255).

\item {} 
\sphinxAtStartPar
\sphinxstyleliteralstrong{\sphinxupquote{g}} (\DUrole{sphinx_autodoc_typehints-type}{\sphinxcode{\sphinxupquote{int}}}) \textendash{} Green component (0..255).

\item {} 
\sphinxAtStartPar
\sphinxstyleliteralstrong{\sphinxupquote{b}} (\DUrole{sphinx_autodoc_typehints-type}{\sphinxcode{\sphinxupquote{int}}}) \textendash{} Blue component (0..255).

\end{itemize}

\sphinxlineitem{Return type}
\sphinxAtStartPar
\DUrole{sphinx_autodoc_typehints-type}{\sphinxcode{\sphinxupquote{str}}}

\sphinxlineitem{Returns}
\sphinxAtStartPar
Uppercase hex string in the form \sphinxcode{\sphinxupquote{\#RRGGBB}}.

\end{description}\end{quote}

\end{fulllineitems}

\index{rgb\_to\_hsl() (in module colors)@\spxentry{rgb\_to\_hsl()}\spxextra{in module colors}}

\begin{fulllineitems}
\phantomsection\label{\detokenize{api/colors:colors.rgb_to_hsl}}
\pysigstartsignatures
\pysiglinewithargsret
{\sphinxcode{\sphinxupquote{colors.}}\sphinxbfcode{\sphinxupquote{rgb\_to\_hsl}}}
{\sphinxparam{\DUrole{n}{r}}\sphinxparamcomma \sphinxparam{\DUrole{n}{g}}\sphinxparamcomma \sphinxparam{\DUrole{n}{b}}}
{}
\pysigstopsignatures
\sphinxAtStartPar
Convert RGB (0..255) to HSL\sphinxhyphen{}like integers.

\sphinxAtStartPar
Internally uses Python’s \sphinxcode{\sphinxupquote{colorsys.rgb\_to\_hls}} which returns (H, L, S).
This function returns an integer tuple that matches your UI needs:
\begin{itemize}
\item {} 
\sphinxAtStartPar
H: 0..359 (degrees)

\item {} 
\sphinxAtStartPar
S: 0..100 (percent)

\item {} 
\sphinxAtStartPar
L: 0..100 (percent)

\end{itemize}
\begin{quote}\begin{description}
\sphinxlineitem{Parameters}\begin{itemize}
\item {} 
\sphinxAtStartPar
\sphinxstyleliteralstrong{\sphinxupquote{r}} (\DUrole{sphinx_autodoc_typehints-type}{\sphinxcode{\sphinxupquote{int}}}) \textendash{} Red component (0..255).

\item {} 
\sphinxAtStartPar
\sphinxstyleliteralstrong{\sphinxupquote{g}} (\DUrole{sphinx_autodoc_typehints-type}{\sphinxcode{\sphinxupquote{int}}}) \textendash{} Green component (0..255).

\item {} 
\sphinxAtStartPar
\sphinxstyleliteralstrong{\sphinxupquote{b}} (\DUrole{sphinx_autodoc_typehints-type}{\sphinxcode{\sphinxupquote{int}}}) \textendash{} Blue component (0..255).

\end{itemize}

\sphinxlineitem{Return type}
\sphinxAtStartPar
\DUrole{sphinx_autodoc_typehints-type}{\sphinxcode{\sphinxupquote{Tuple}}{[}\sphinxcode{\sphinxupquote{int}}, \sphinxcode{\sphinxupquote{int}}, \sphinxcode{\sphinxupquote{int}}{]}}

\sphinxlineitem{Returns}
\sphinxAtStartPar
Tuple \sphinxcode{\sphinxupquote{(h, s, l)}} with integer values.

\end{description}\end{quote}

\end{fulllineitems}

\index{hsl\_to\_rgb() (in module colors)@\spxentry{hsl\_to\_rgb()}\spxextra{in module colors}}

\begin{fulllineitems}
\phantomsection\label{\detokenize{api/colors:colors.hsl_to_rgb}}
\pysigstartsignatures
\pysiglinewithargsret
{\sphinxcode{\sphinxupquote{colors.}}\sphinxbfcode{\sphinxupquote{hsl\_to\_rgb}}}
{\sphinxparam{\DUrole{n}{h}}\sphinxparamcomma \sphinxparam{\DUrole{n}{s}}\sphinxparamcomma \sphinxparam{\DUrole{n}{l}}}
{}
\pysigstopsignatures
\sphinxAtStartPar
Convert HSL\sphinxhyphen{}like integers to RGB (0..255).

\sphinxAtStartPar
This function expects H in degrees and S/L in percent. Values are clamped
where appropriate, and H is normalized with modulo 360.

\sphinxAtStartPar
Note: \sphinxcode{\sphinxupquote{colorsys.hls\_to\_rgb}} expects (H, L, S), hence the parameter order.
\begin{quote}\begin{description}
\sphinxlineitem{Parameters}\begin{itemize}
\item {} 
\sphinxAtStartPar
\sphinxstyleliteralstrong{\sphinxupquote{h}} (\DUrole{sphinx_autodoc_typehints-type}{\sphinxcode{\sphinxupquote{int}}}) \textendash{} Hue in degrees (any int; normalized to 0..359).

\item {} 
\sphinxAtStartPar
\sphinxstyleliteralstrong{\sphinxupquote{s}} (\DUrole{sphinx_autodoc_typehints-type}{\sphinxcode{\sphinxupquote{int}}}) \textendash{} Saturation in percent (0..100).

\item {} 
\sphinxAtStartPar
\sphinxstyleliteralstrong{\sphinxupquote{l}} (\DUrole{sphinx_autodoc_typehints-type}{\sphinxcode{\sphinxupquote{int}}}) \textendash{} Lightness in percent (0..100).

\end{itemize}

\sphinxlineitem{Return type}
\sphinxAtStartPar
\DUrole{sphinx_autodoc_typehints-type}{\sphinxcode{\sphinxupquote{Tuple}}{[}\sphinxcode{\sphinxupquote{int}}, \sphinxcode{\sphinxupquote{int}}, \sphinxcode{\sphinxupquote{int}}{]}}

\sphinxlineitem{Returns}
\sphinxAtStartPar
Tuple \sphinxcode{\sphinxupquote{(r, g, b)}} with integer values 0..255.

\end{description}\end{quote}

\end{fulllineitems}

\index{luminance() (in module colors)@\spxentry{luminance()}\spxextra{in module colors}}

\begin{fulllineitems}
\phantomsection\label{\detokenize{api/colors:colors.luminance}}
\pysigstartsignatures
\pysiglinewithargsret
{\sphinxcode{\sphinxupquote{colors.}}\sphinxbfcode{\sphinxupquote{luminance}}}
{\sphinxparam{\DUrole{n}{rgb}}}
{}
\pysigstopsignatures
\sphinxAtStartPar
Compute a simple perceived luminance from an RGB tuple.

\sphinxAtStartPar
Uses weighted coefficients (Rec. 709\sphinxhyphen{}like) and normalizes the result to 0..1.
\begin{quote}\begin{description}
\sphinxlineitem{Parameters}
\sphinxAtStartPar
\sphinxstyleliteralstrong{\sphinxupquote{rgb}} (\DUrole{sphinx_autodoc_typehints-type}{\sphinxcode{\sphinxupquote{Tuple}}{[}\sphinxcode{\sphinxupquote{int}}, \sphinxcode{\sphinxupquote{int}}, \sphinxcode{\sphinxupquote{int}}{]}}) \textendash{} Tuple \sphinxcode{\sphinxupquote{(r, g, b)}} with values 0..255.

\sphinxlineitem{Return type}
\sphinxAtStartPar
\DUrole{sphinx_autodoc_typehints-type}{\sphinxcode{\sphinxupquote{float}}}

\sphinxlineitem{Returns}
\sphinxAtStartPar
Luminance value in the range 0..1.

\end{description}\end{quote}

\end{fulllineitems}

\index{blend() (in module colors)@\spxentry{blend()}\spxextra{in module colors}}

\begin{fulllineitems}
\phantomsection\label{\detokenize{api/colors:colors.blend}}
\pysigstartsignatures
\pysiglinewithargsret
{\sphinxcode{\sphinxupquote{colors.}}\sphinxbfcode{\sphinxupquote{blend}}}
{\sphinxparam{\DUrole{n}{a}}\sphinxparamcomma \sphinxparam{\DUrole{n}{b}}\sphinxparamcomma \sphinxparam{\DUrole{n}{t}}}
{}
\pysigstopsignatures
\sphinxAtStartPar
Blend two RGB colors.

\sphinxAtStartPar
Linear interpolation between color \sphinxcode{\sphinxupquote{a}} and \sphinxcode{\sphinxupquote{b}} by factor \sphinxcode{\sphinxupquote{t}}.
\begin{itemize}
\item {} 
\sphinxAtStartPar
\sphinxcode{\sphinxupquote{t = 0}} returns \sphinxcode{\sphinxupquote{a}}

\item {} 
\sphinxAtStartPar
\sphinxcode{\sphinxupquote{t = 1}} returns \sphinxcode{\sphinxupquote{b}}

\end{itemize}
\begin{quote}\begin{description}
\sphinxlineitem{Parameters}\begin{itemize}
\item {} 
\sphinxAtStartPar
\sphinxstyleliteralstrong{\sphinxupquote{a}} (\DUrole{sphinx_autodoc_typehints-type}{\sphinxcode{\sphinxupquote{Tuple}}{[}\sphinxcode{\sphinxupquote{int}}, \sphinxcode{\sphinxupquote{int}}, \sphinxcode{\sphinxupquote{int}}{]}}) \textendash{} First color as \sphinxcode{\sphinxupquote{(r, g, b)}}.

\item {} 
\sphinxAtStartPar
\sphinxstyleliteralstrong{\sphinxupquote{b}} (\DUrole{sphinx_autodoc_typehints-type}{\sphinxcode{\sphinxupquote{Tuple}}{[}\sphinxcode{\sphinxupquote{int}}, \sphinxcode{\sphinxupquote{int}}, \sphinxcode{\sphinxupquote{int}}{]}}) \textendash{} Second color as \sphinxcode{\sphinxupquote{(r, g, b)}}.

\item {} 
\sphinxAtStartPar
\sphinxstyleliteralstrong{\sphinxupquote{t}} (\DUrole{sphinx_autodoc_typehints-type}{\sphinxcode{\sphinxupquote{float}}}) \textendash{} Blend factor in the range 0..1 (clamped).

\end{itemize}

\sphinxlineitem{Return type}
\sphinxAtStartPar
\DUrole{sphinx_autodoc_typehints-type}{\sphinxcode{\sphinxupquote{Tuple}}{[}\sphinxcode{\sphinxupquote{int}}, \sphinxcode{\sphinxupquote{int}}, \sphinxcode{\sphinxupquote{int}}{]}}

\sphinxlineitem{Returns}
\sphinxAtStartPar
Blended color as \sphinxcode{\sphinxupquote{(r, g, b)}}.

\end{description}\end{quote}

\end{fulllineitems}

\index{derive\_theme\_from\_background() (in module colors)@\spxentry{derive\_theme\_from\_background()}\spxextra{in module colors}}

\begin{fulllineitems}
\phantomsection\label{\detokenize{api/colors:colors.derive_theme_from_background}}
\pysigstartsignatures
\pysiglinewithargsret
{\sphinxcode{\sphinxupquote{colors.}}\sphinxbfcode{\sphinxupquote{derive\_theme\_from\_background}}}
{\sphinxparam{\DUrole{n}{bg\_hex}}\sphinxparamcomma \sphinxparam{\DUrole{n}{base\_theme}}}
{}
\pysigstopsignatures
\sphinxAtStartPar
Derive a full theme palette from a background color.

\sphinxAtStartPar
Given a background color, this function computes:
\sphinxhyphen{} surface colors (slightly blended toward white/black)
\sphinxhyphen{} border color
\sphinxhyphen{} text color (pulled toward white/black but influenced by the background)
\sphinxhyphen{} muted text color (between text and border)

\sphinxAtStartPar
The algorithm adapts depending on whether the background is considered dark
(luminance \textless{} 0.5) or light.
\begin{quote}\begin{description}
\sphinxlineitem{Parameters}\begin{itemize}
\item {} 
\sphinxAtStartPar
\sphinxstyleliteralstrong{\sphinxupquote{bg\_hex}} (\DUrole{sphinx_autodoc_typehints-type}{\sphinxcode{\sphinxupquote{str}}}) \textendash{} Background color as \sphinxcode{\sphinxupquote{\#RRGGBB}}.

\item {} 
\sphinxAtStartPar
\sphinxstyleliteralstrong{\sphinxupquote{base\_theme}} (\DUrole{sphinx_autodoc_typehints-type}{\sphinxcode{\sphinxupquote{dict}}}) \textendash{} Base theme dictionary to copy and override. Must contain keys
like \sphinxcode{\sphinxupquote{text}}, \sphinxcode{\sphinxupquote{primary}}, \sphinxcode{\sphinxupquote{danger}} etc. Missing keys are kept as\sphinxhyphen{}is.

\end{itemize}

\sphinxlineitem{Returns}
\sphinxAtStartPar
\sphinxcode{\sphinxupquote{background}}, \sphinxcode{\sphinxupquote{surface}}, \sphinxcode{\sphinxupquote{surface\_alt}}, \sphinxcode{\sphinxupquote{border}}, \sphinxcode{\sphinxupquote{text}},
and \sphinxcode{\sphinxupquote{muted\_text}}.

\sphinxlineitem{Return type}
\sphinxAtStartPar
\DUrole{sphinx_autodoc_typehints-type}{\sphinxcode{\sphinxupquote{dict}}}

\end{description}\end{quote}

\end{fulllineitems}


\sphinxstepscope


\section{config module}
\label{\detokenize{api/config:module-config}}\label{\detokenize{api/config:config-module}}\label{\detokenize{api/config::doc}}\index{module@\spxentry{module}!config@\spxentry{config}}\index{config@\spxentry{config}!module@\spxentry{module}}\index{default\_fields\_config() (in module config)@\spxentry{default\_fields\_config()}\spxextra{in module config}}

\begin{fulllineitems}
\phantomsection\label{\detokenize{api/config:config.default_fields_config}}
\pysigstartsignatures
\pysiglinewithargsret
{\sphinxcode{\sphinxupquote{config.}}\sphinxbfcode{\sphinxupquote{default\_fields\_config}}}
{}
{}
\pysigstopsignatures
\sphinxAtStartPar
Configurações padrão para campos e abas.
\begin{quote}\begin{description}
\sphinxlineitem{Returns}
\sphinxAtStartPar
Configurações iniciais com uma aba “Preenche Fácil” contendo campos básicos.

\sphinxlineitem{Return type}
\sphinxAtStartPar
\DUrole{sphinx_autodoc_typehints-type}{\sphinxcode{\sphinxupquote{dict}}}

\end{description}\end{quote}

\end{fulllineitems}

\index{rename\_sheet\_key() (in module config)@\spxentry{rename\_sheet\_key()}\spxextra{in module config}}

\begin{fulllineitems}
\phantomsection\label{\detokenize{api/config:config.rename_sheet_key}}
\pysigstartsignatures
\pysiglinewithargsret
{\sphinxcode{\sphinxupquote{config.}}\sphinxbfcode{\sphinxupquote{rename\_sheet\_key}}}
{\sphinxparam{\DUrole{n}{cfg}}\sphinxparamcomma \sphinxparam{\DUrole{n}{old\_name}}\sphinxparamcomma \sphinxparam{\DUrole{n}{new\_name}}}
{}
\pysigstopsignatures
\sphinxAtStartPar
Renomeia uma aba (sheet) dentro da configuração, movendo os campos associados para a nova chave.
\begin{quote}\begin{description}
\sphinxlineitem{Parameters}\begin{itemize}
\item {} 
\sphinxAtStartPar
\sphinxstyleliteralstrong{\sphinxupquote{cfg}} (\DUrole{sphinx_autodoc_typehints-type}{\sphinxcode{\sphinxupquote{dict}}}) \textendash{} Configuração completa onde as abas estão definidas. Deve conter a chave “abas” que é um dicionário de abas.

\item {} 
\sphinxAtStartPar
\sphinxstyleliteralstrong{\sphinxupquote{old\_name}} (\DUrole{sphinx_autodoc_typehints-type}{\sphinxcode{\sphinxupquote{str}}}) \textendash{} Nome atual da aba que se deseja renomear. Se não existir, a função apenas garante que a nova aba exista sem campos.

\item {} 
\sphinxAtStartPar
\sphinxstyleliteralstrong{\sphinxupquote{new\_name}} (\DUrole{sphinx_autodoc_typehints-type}{\sphinxcode{\sphinxupquote{str}}}) \textendash{} Novo nome para a aba. Se já existir, os campos só serão movidos se a aba de destino estiver vazia. Se for igual ao old\_name, não faz nada.

\end{itemize}

\sphinxlineitem{Return type}
\sphinxAtStartPar
\DUrole{sphinx_autodoc_typehints-type}{\sphinxcode{\sphinxupquote{None}}}

\end{description}\end{quote}

\end{fulllineitems}

\index{get\_sheet\_campos() (in module config)@\spxentry{get\_sheet\_campos()}\spxextra{in module config}}

\begin{fulllineitems}
\phantomsection\label{\detokenize{api/config:config.get_sheet_campos}}
\pysigstartsignatures
\pysiglinewithargsret
{\sphinxcode{\sphinxupquote{config.}}\sphinxbfcode{\sphinxupquote{get\_sheet\_campos}}}
{\sphinxparam{\DUrole{n}{cfg}}\sphinxparamcomma \sphinxparam{\DUrole{n}{sheet\_name}}}
{}
\pysigstopsignatures
\sphinxAtStartPar
Obtém os campos configurados para uma aba específica dentro da configuração.
\begin{quote}\begin{description}
\sphinxlineitem{Parameters}\begin{itemize}
\item {} 
\sphinxAtStartPar
\sphinxstyleliteralstrong{\sphinxupquote{cfg}} (\DUrole{sphinx_autodoc_typehints-type}{\sphinxcode{\sphinxupquote{dict}}}) \textendash{} Configuração completa onde as abas estão definidas. Deve conter a chave “abas” que é um dicionário de abas, cada uma com uma lista de campos.

\item {} 
\sphinxAtStartPar
\sphinxstyleliteralstrong{\sphinxupquote{sheet\_name}} (\DUrole{sphinx_autodoc_typehints-type}{\sphinxcode{\sphinxupquote{str}}}) \textendash{} Nome da aba para a qual se deseja obter os campos. Se a aba não existir, retorna uma lista vazia.

\end{itemize}

\sphinxlineitem{Returns}
\sphinxAtStartPar
Lista de campos configurados para a aba especificada. Cada campo é um dicionário com chaves como “id”, “titulo”, “tipo” e “fixo”. Se a aba ou os campos não estiverem definidos, retorna uma lista vazia.

\sphinxlineitem{Return type}
\sphinxAtStartPar
\DUrole{sphinx_autodoc_typehints-type}{\sphinxcode{\sphinxupquote{list}}}

\end{description}\end{quote}

\end{fulllineitems}

\index{set\_sheet\_campos() (in module config)@\spxentry{set\_sheet\_campos()}\spxextra{in module config}}

\begin{fulllineitems}
\phantomsection\label{\detokenize{api/config:config.set_sheet_campos}}
\pysigstartsignatures
\pysiglinewithargsret
{\sphinxcode{\sphinxupquote{config.}}\sphinxbfcode{\sphinxupquote{set\_sheet\_campos}}}
{\sphinxparam{\DUrole{n}{cfg}}\sphinxparamcomma \sphinxparam{\DUrole{n}{sheet\_name}}\sphinxparamcomma \sphinxparam{\DUrole{n}{campos}}}
{}
\pysigstopsignatures
\sphinxAtStartPar
Define os campos para uma aba específica dentro da configuração, normalizando a lista de campos e garantindo a estrutura correta.
\begin{quote}\begin{description}
\sphinxlineitem{Parameters}\begin{itemize}
\item {} 
\sphinxAtStartPar
\sphinxstyleliteralstrong{\sphinxupquote{cfg}} (\DUrole{sphinx_autodoc_typehints-type}{\sphinxcode{\sphinxupquote{dict}}}) \textendash{} Configuração completa onde as abas estão definidas. Deve conter a chave “abas” que é um dicionário de abas, cada uma com uma lista de campos.

\item {} 
\sphinxAtStartPar
\sphinxstyleliteralstrong{\sphinxupquote{sheet\_name}} (\DUrole{sphinx_autodoc_typehints-type}{\sphinxcode{\sphinxupquote{str}}}) \textendash{} Nome da aba para a qual se deseja definir os campos. Se a aba não existir, será criada automaticamente.

\item {} 
\sphinxAtStartPar
\sphinxstyleliteralstrong{\sphinxupquote{campos}} (\DUrole{sphinx_autodoc_typehints-type}{\sphinxcode{\sphinxupquote{list}}}) \textendash{} Lista de campos a ser definida para a aba especificada. Cada campo deve ser um dicionário, mas o formato é flexível. A função irá normalizar os campos, garantindo que cada um tenha as chaves mínimas e tipos válidos, conforme definido na função \_normalize\_campos.

\end{itemize}

\sphinxlineitem{Return type}
\sphinxAtStartPar
\DUrole{sphinx_autodoc_typehints-type}{\sphinxcode{\sphinxupquote{None}}}

\end{description}\end{quote}

\end{fulllineitems}

\index{load\_fields\_config() (in module config)@\spxentry{load\_fields\_config()}\spxextra{in module config}}

\begin{fulllineitems}
\phantomsection\label{\detokenize{api/config:config.load_fields_config}}
\pysigstartsignatures
\pysiglinewithargsret
{\sphinxcode{\sphinxupquote{config.}}\sphinxbfcode{\sphinxupquote{load\_fields\_config}}}
{\sphinxparam{\DUrole{n}{field\_types}}}
{}
\pysigstopsignatures
\sphinxAtStartPar
Carrega a configuração de campos e abas a partir do arquivo JSON, aplicando normalização e migração de formato se necessário.
\begin{quote}\begin{description}
\sphinxlineitem{Parameters}
\sphinxAtStartPar
\sphinxstyleliteralstrong{\sphinxupquote{field\_types}} (\DUrole{sphinx_autodoc_typehints-type}{\sphinxcode{\sphinxupquote{List}}{[}\sphinxcode{\sphinxupquote{str}}{]}}) \textendash{} Lista de tipos de campo válidos para normalização. Usada para garantir que os campos carregados tenham tipos reconhecidos, definindo como “texto” se o tipo for inválido ou ausente.

\sphinxlineitem{Returns}
\sphinxAtStartPar
Configuração de campos carregada e normalizada.

\sphinxlineitem{Return type}
\sphinxAtStartPar
\DUrole{sphinx_autodoc_typehints-type}{\sphinxcode{\sphinxupquote{dict}}}

\end{description}\end{quote}

\end{fulllineitems}

\index{save\_fields\_config() (in module config)@\spxentry{save\_fields\_config()}\spxextra{in module config}}

\begin{fulllineitems}
\phantomsection\label{\detokenize{api/config:config.save_fields_config}}
\pysigstartsignatures
\pysiglinewithargsret
{\sphinxcode{\sphinxupquote{config.}}\sphinxbfcode{\sphinxupquote{save\_fields\_config}}}
{\sphinxparam{\DUrole{n}{cfg}}}
{}
\pysigstopsignatures
\sphinxAtStartPar
Salva a configuração de campos e abas em um arquivo JSON, mantendo a estrutura esperada.
\begin{quote}\begin{description}
\sphinxlineitem{Parameters}
\sphinxAtStartPar
\sphinxstyleliteralstrong{\sphinxupquote{cfg}} (\DUrole{sphinx_autodoc_typehints-type}{\sphinxcode{\sphinxupquote{dict}}}) \textendash{} Configuração de campos a ser salva. Deve conter a estrutura esperada, incluindo “abas” como um dicionário de abas, cada uma com uma lista de campos. A função não realiza validação ou normalização, portanto, é responsabilidade do chamador garantir que a configuração esteja no formato correto antes de chamar esta função.

\sphinxlineitem{Return type}
\sphinxAtStartPar
\DUrole{sphinx_autodoc_typehints-type}{\sphinxcode{\sphinxupquote{None}}}

\end{description}\end{quote}

\end{fulllineitems}

\index{load\_sheet\_config() (in module config)@\spxentry{load\_sheet\_config()}\spxextra{in module config}}

\begin{fulllineitems}
\phantomsection\label{\detokenize{api/config:config.load_sheet_config}}
\pysigstartsignatures
\pysiglinewithargsret
{\sphinxcode{\sphinxupquote{config.}}\sphinxbfcode{\sphinxupquote{load\_sheet\_config}}}
{}
{}
\pysigstopsignatures
\sphinxAtStartPar
Carrega a configuração de planilha a partir do arquivo JSON. Se o arquivo não existir ou ocorrer um erro durante a leitura, retorna um dicionário vazio. A função não realiza validação ou normalização da configuração, portanto, é responsabilidade do chamador garantir que a configuração esteja no formato correto após o carregamento.
\begin{quote}\begin{description}
\sphinxlineitem{Returns}
\sphinxAtStartPar
Configuração de planilha carregada a partir do arquivo JSON. Se o arquivo não existir ou ocorrer um erro durante a leitura, retorna um dicionário vazio. A função não realiza validação ou normalização da configuração, portanto, é responsabilidade do chamador garantir que a configuração esteja no formato correto após o carregamento.

\sphinxlineitem{Return type}
\sphinxAtStartPar
\DUrole{sphinx_autodoc_typehints-type}{\sphinxcode{\sphinxupquote{dict}}}

\end{description}\end{quote}

\end{fulllineitems}

\index{save\_sheet\_config() (in module config)@\spxentry{save\_sheet\_config()}\spxextra{in module config}}

\begin{fulllineitems}
\phantomsection\label{\detokenize{api/config:config.save_sheet_config}}
\pysigstartsignatures
\pysiglinewithargsret
{\sphinxcode{\sphinxupquote{config.}}\sphinxbfcode{\sphinxupquote{save\_sheet\_config}}}
{\sphinxparam{\DUrole{n}{path}}}
{}
\pysigstopsignatures
\sphinxAtStartPar
Salva o caminho da planilha na configuração, escrevendo em um arquivo JSON. A função irá criar ou sobrescrever o arquivo de configuração com um dicionário contendo a chave “last\_sheet\_path” e o valor do caminho fornecido. A função não realiza validação do caminho fornecido, portanto, é responsabilidade do chamador garantir que o caminho seja válido e acessível antes de chamar esta função.
\begin{quote}\begin{description}
\sphinxlineitem{Parameters}
\sphinxAtStartPar
\sphinxstyleliteralstrong{\sphinxupquote{path}} (\DUrole{sphinx_autodoc_typehints-type}{\sphinxcode{\sphinxupquote{str}}}) \textendash{} Caminho da planilha a ser salvo na configuração. A função irá salvar este caminho no arquivo de configuração JSON, substituindo o valor existente. A função não realiza validação do caminho fornecido, portanto, é responsabilidade do chamador garantir que o caminho seja válido e acessível antes de chamar esta função.

\sphinxlineitem{Return type}
\sphinxAtStartPar
\DUrole{sphinx_autodoc_typehints-type}{\sphinxcode{\sphinxupquote{None}}}

\end{description}\end{quote}

\end{fulllineitems}

\index{get\_last\_sheet\_path() (in module config)@\spxentry{get\_last\_sheet\_path()}\spxextra{in module config}}

\begin{fulllineitems}
\phantomsection\label{\detokenize{api/config:config.get_last_sheet_path}}
\pysigstartsignatures
\pysiglinewithargsret
{\sphinxcode{\sphinxupquote{config.}}\sphinxbfcode{\sphinxupquote{get\_last\_sheet\_path}}}
{}
{}
\pysigstopsignatures
\sphinxAtStartPar
Obtém o caminho da última planilha salvo na configuração. A função lê o arquivo de configuração JSON para obter o valor associado à chave “last\_sheet\_path”. Se o arquivo de configuração não existir, se a chave “last\_sheet\_path” não estiver presente ou se o caminho armazenado não corresponder a um arquivo existente no sistema de arquivos, a função retorna None. Caso contrário, retorna o caminho da última planilha como uma string.
\begin{quote}\begin{description}
\sphinxlineitem{Returns}
\sphinxAtStartPar
Caminho da última planilha salvo na configuração. Se o caminho não estiver presente na configuração ou se o arquivo de configuração não existir, retorna None. Se o caminho estiver presente mas o arquivo correspondente não existir no sistema de arquivos, também retorna None. Caso contrário, retorna o caminho da última planilha como uma string.

\sphinxlineitem{Return type}
\sphinxAtStartPar
\DUrole{sphinx_autodoc_typehints-type}{\sphinxcode{\sphinxupquote{Optional}}{[}\sphinxcode{\sphinxupquote{str}}{]}}

\end{description}\end{quote}

\end{fulllineitems}

\index{default\_ui\_config() (in module config)@\spxentry{default\_ui\_config()}\spxextra{in module config}}

\begin{fulllineitems}
\phantomsection\label{\detokenize{api/config:config.default_ui_config}}
\pysigstartsignatures
\pysiglinewithargsret
{\sphinxcode{\sphinxupquote{config.}}\sphinxbfcode{\sphinxupquote{default\_ui\_config}}}
{}
{}
\pysigstopsignatures
\sphinxAtStartPar
Configurações padrão para a interface do usuário, incluindo ícones, tema de cores e HTML para o rodapé. Esta configuração é usada como base para a personalização da aparência da aplicação. A função retorna um dicionário contendo as seguintes chaves:
\begin{quote}\begin{description}
\sphinxlineitem{Returns}
\sphinxAtStartPar
Configurações padrão para a interface do usuário, incluindo ícones, tema de cores e HTML para o rodapé. Esta configuração é usada como base para a personalização da aparência da aplicação. A função retorna um dicionário contendo as seguintes chaves:

\sphinxlineitem{Return type}
\sphinxAtStartPar
\DUrole{sphinx_autodoc_typehints-type}{\sphinxcode{\sphinxupquote{dict}}}

\end{description}\end{quote}
\begin{itemize}
\item {} 
\sphinxAtStartPar
“window\_icon”: Caminho para o ícone da janela principal.

\item {} 
\sphinxAtStartPar
“button\_icons”: Dicionário de caminhos para ícones de botões específicos, como “choose\_file”, “add\_field”, “save\_lead”, “clear”, “edit\_title”, “delete\_field” e “settings”.

\item {} 
\sphinxAtStartPar
“theme”: Dicionário de cores para a interface, incluindo “background”, “surface

\end{itemize}

\sphinxAtStartPar
“, “surface\_alt”, “text”, “muted\_text”, “primary”, “danger” e “border”.
\sphinxhyphen{} “background\_hsl”: Dicionário com valores de matiz (h), saturação (s) e luminosidade (l) para a cor de fundo, usado para ajustes dinâmicos de tema.
\sphinxhyphen{} “footer\_left\_html”: String contendo HTML para o conteúdo do rodapé à esquerda, permitindo personalização de mensagens ou links.

\end{fulllineitems}

\index{load\_ui\_config() (in module config)@\spxentry{load\_ui\_config()}\spxextra{in module config}}

\begin{fulllineitems}
\phantomsection\label{\detokenize{api/config:config.load_ui_config}}
\pysigstartsignatures
\pysiglinewithargsret
{\sphinxcode{\sphinxupquote{config.}}\sphinxbfcode{\sphinxupquote{load\_ui\_config}}}
{}
{}
\pysigstopsignatures
\sphinxAtStartPar
Carrega as configurações da interface do usuário a partir de um arquivo JSON. Se o arquivo não existir ou ocorrer um erro durante a leitura, retorna as configurações padrão definidas pela função default\_ui\_config(). A função garante que as chaves principais estejam presentes na configuração carregada, preenchendo com valores padrão para quaisquer chaves ausentes. A configuração inclui temas de cores, ícones e HTML para o rodapé, permitindo personalização da aparência da aplicação.
:returns: Configurações da interface do usuário carregadas a partir do arquivo JSON. Se o arquivo não existir ou ocorrer um erro durante a leitura, retorna as configurações padrão definidas pela função default\_ui\_config(). A função garante que as chaves principais estejam presentes na configuração carregada, preenchendo com valores padrão para quaisquer chaves ausentes. A configuração inclui temas de cores, ícones e HTML para o rodapé, permitindo personalização da aparência da aplicação.
\begin{quote}\begin{description}
\sphinxlineitem{Return type}
\sphinxAtStartPar
\DUrole{sphinx_autodoc_typehints-type}{\sphinxcode{\sphinxupquote{dict}}}

\end{description}\end{quote}

\end{fulllineitems}

\index{save\_ui\_config() (in module config)@\spxentry{save\_ui\_config()}\spxextra{in module config}}

\begin{fulllineitems}
\phantomsection\label{\detokenize{api/config:config.save_ui_config}}
\pysigstartsignatures
\pysiglinewithargsret
{\sphinxcode{\sphinxupquote{config.}}\sphinxbfcode{\sphinxupquote{save\_ui\_config}}}
{\sphinxparam{\DUrole{n}{cfg}}}
{}
\pysigstopsignatures
\sphinxAtStartPar
Salva as configurações da interface do usuário em um arquivo JSON. A função recebe um dicionário de configurações e escreve seu conteúdo em um arquivo de configuração JSON, criando ou sobrescrevendo o arquivo existente. A configuração deve conter chaves principais como “window\_icon”, “button\_icons”, “theme”, “background\_hsl” e “footer\_left\_html”. A função não realiza validação ou normalização da configuração, portanto, é responsabilidade do chamador garantir que a configuração esteja no formato correto antes de chamar esta função. O arquivo JSON resultante terá uma estrutura legível, com indentação para facilitar a edição manual se necessário.
\begin{quote}\begin{description}
\sphinxlineitem{Parameters}
\sphinxAtStartPar
\sphinxstyleliteralstrong{\sphinxupquote{cfg}} (\DUrole{sphinx_autodoc_typehints-type}{\sphinxcode{\sphinxupquote{dict}}}) \textendash{} Configurações da interface do usuário a ser salva. Deve conter as chaves principais como “window\_icon”, “button\_icons”, “theme”, “background\_hsl” e “footer\_left\_html”. A função não realiza validação ou normalização da configuração, portanto, é responsabilidade do chamador garantir que a configuração esteja no formato correto antes de chamar esta função. A função irá criar ou sobrescrever o arquivo de configuração JSON com o conteúdo fornecido.

\sphinxlineitem{Return type}
\sphinxAtStartPar
\DUrole{sphinx_autodoc_typehints-type}{\sphinxcode{\sphinxupquote{None}}}

\end{description}\end{quote}

\end{fulllineitems}

\index{default\_email\_domains\_config() (in module config)@\spxentry{default\_email\_domains\_config()}\spxextra{in module config}}

\begin{fulllineitems}
\phantomsection\label{\detokenize{api/config:config.default_email_domains_config}}
\pysigstartsignatures
\pysiglinewithargsret
{\sphinxcode{\sphinxupquote{config.}}\sphinxbfcode{\sphinxupquote{default\_email\_domains\_config}}}
{}
{}
\pysigstopsignatures
\sphinxAtStartPar
Configurações padrão para domínios de email, contendo uma lista de domínios comuns usados para preenchimento automático em campos de email. Esta configuração é usada como base para a personalização dos domínios de email disponíveis na aplicação. A função retorna um dicionário com a chave “dominios”, que é uma lista de strings representando os domínios de email padrão, como “@gmail.com”, “@hotmail.com”, “@yahoo.com” e “@outlook.com”.
\begin{quote}\begin{description}
\sphinxlineitem{Returns}
\sphinxAtStartPar
Configurações padrão para domínios de email, contendo uma lista de domínios comuns usados para preenchimento automático em campos de email. Esta configuração é usada como base para a personalização dos domínios de email disponíveis na aplicação. A função retorna um dicionário com a chave “dominios”, que é uma lista de strings representando os domínios de email padrão, como “@gmail.com”, “@hotmail.com”, “@yahoo.com” e “@outlook.com”.

\sphinxlineitem{Return type}
\sphinxAtStartPar
\DUrole{sphinx_autodoc_typehints-type}{\sphinxcode{\sphinxupquote{dict}}}

\end{description}\end{quote}

\end{fulllineitems}

\index{load\_email\_domains\_config() (in module config)@\spxentry{load\_email\_domains\_config()}\spxextra{in module config}}

\begin{fulllineitems}
\phantomsection\label{\detokenize{api/config:config.load_email_domains_config}}
\pysigstartsignatures
\pysiglinewithargsret
{\sphinxcode{\sphinxupquote{config.}}\sphinxbfcode{\sphinxupquote{load\_email\_domains\_config}}}
{}
{}
\pysigstopsignatures
\sphinxAtStartPar
Carrega as configurações de domínios de email a partir de um arquivo JSON. Se o arquivo não existir ou ocorrer um erro durante a leitura, retorna as configurações padrão definidas pela função default\_email\_domains\_config(). A função garante que a chave “dominios” esteja presente na configuração carregada, preenchendo com valores padrão para quaisquer chaves ausentes. A configuração inclui uma lista de domínios de email usados para preenchimento automático em campos de email, permitindo personalização dos domínios disponíveis na aplicação.
\begin{quote}\begin{description}
\sphinxlineitem{Returns}
\sphinxAtStartPar
Configurações de domínios de email carregadas do arquivo JSON ou as configurações padrão caso o arquivo não exista ou ocorra um erro durante a leitura.

\sphinxlineitem{Return type}
\sphinxAtStartPar
\DUrole{sphinx_autodoc_typehints-type}{\sphinxcode{\sphinxupquote{dict}}}

\end{description}\end{quote}

\end{fulllineitems}

\index{save\_email\_domains\_config() (in module config)@\spxentry{save\_email\_domains\_config()}\spxextra{in module config}}

\begin{fulllineitems}
\phantomsection\label{\detokenize{api/config:config.save_email_domains_config}}
\pysigstartsignatures
\pysiglinewithargsret
{\sphinxcode{\sphinxupquote{config.}}\sphinxbfcode{\sphinxupquote{save\_email\_domains\_config}}}
{\sphinxparam{\DUrole{n}{cfg}}}
{}
\pysigstopsignatures
\sphinxAtStartPar
Salva as configurações de domínios de email em um arquivo JSON. A função recebe um dicionário de configurações contendo a chave “dominios”, que é uma lista de strings representando os domínios de email. O conteúdo do dicionário é escrito em um arquivo de configuração JSON, criando ou sobrescrevendo o arquivo existente. A função não realiza validação ou normalização da configuração, portanto, é responsabilidade do chamador garantir que a configuração esteja no formato correto antes de chamar esta função. O arquivo JSON resultante terá uma estrutura legível, com indentação para facilitar a edição manual se necessário.
\begin{quote}\begin{description}
\sphinxlineitem{Parameters}
\sphinxAtStartPar
\sphinxstyleliteralstrong{\sphinxupquote{cfg}} (\DUrole{sphinx_autodoc_typehints-type}{\sphinxcode{\sphinxupquote{dict}}}) \textendash{} Configurações de domínios de email a ser salva. Deve conter a chave “dominios”, que é uma lista de strings representando os domínios de email. A função não realiza validação ou normalização da configuração, portanto, é responsabilidade do chamador garantir que a configuração esteja no formato correto antes de chamar esta função. A função irá criar ou sobrescrever o arquivo de configuração JSON com o conteúdo fornecido, mantendo uma estrutura legível com indentação.

\sphinxlineitem{Return type}
\sphinxAtStartPar
\DUrole{sphinx_autodoc_typehints-type}{\sphinxcode{\sphinxupquote{None}}}

\end{description}\end{quote}

\end{fulllineitems}

\index{delete\_sheet\_key() (in module config)@\spxentry{delete\_sheet\_key()}\spxextra{in module config}}

\begin{fulllineitems}
\phantomsection\label{\detokenize{api/config:config.delete_sheet_key}}
\pysigstartsignatures
\pysiglinewithargsret
{\sphinxcode{\sphinxupquote{config.}}\sphinxbfcode{\sphinxupquote{delete\_sheet\_key}}}
{\sphinxparam{\DUrole{n}{cfg\_fields}}\sphinxparamcomma \sphinxparam{\DUrole{n}{sheet\_name}}}
{}
\pysigstopsignatures
\sphinxAtStartPar
Remove do controle JSON a configuração referente a uma aba específica.
\begin{quote}\begin{description}
\sphinxlineitem{Parameters}\begin{itemize}
\item {} 
\sphinxAtStartPar
\sphinxstyleliteralstrong{\sphinxupquote{cfg\_fields}} (\DUrole{sphinx_autodoc_typehints-type}{\sphinxcode{\sphinxupquote{dict}}}) \textendash{} Configuração carregada do controle de campos.

\item {} 
\sphinxAtStartPar
\sphinxstyleliteralstrong{\sphinxupquote{sheet\_name}} (\DUrole{sphinx_autodoc_typehints-type}{\sphinxcode{\sphinxupquote{str}}}) \textendash{} Nome da aba a ser removida.

\end{itemize}

\sphinxlineitem{Returns}
\sphinxAtStartPar
True se removeu algo, False caso contrário.

\sphinxlineitem{Return type}
\sphinxAtStartPar
\DUrole{sphinx_autodoc_typehints-type}{\sphinxcode{\sphinxupquote{bool}}}

\end{description}\end{quote}

\end{fulllineitems}


\sphinxstepscope


\section{dialogs module}
\label{\detokenize{api/dialogs:module-dialogs}}\label{\detokenize{api/dialogs:dialogs-module}}\label{\detokenize{api/dialogs::doc}}\index{module@\spxentry{module}!dialogs@\spxentry{dialogs}}\index{dialogs@\spxentry{dialogs}!module@\spxentry{module}}\index{ThemeDialog (class in dialogs)@\spxentry{ThemeDialog}\spxextra{class in dialogs}}

\begin{fulllineitems}
\phantomsection\label{\detokenize{api/dialogs:dialogs.ThemeDialog}}
\pysigstartsignatures
\pysiglinewithargsret
{\sphinxbfcode{\sphinxupquote{\DUrole{k}{class}\DUrole{w}{ }}}\sphinxcode{\sphinxupquote{dialogs.}}\sphinxbfcode{\sphinxupquote{ThemeDialog}}}
{\sphinxparam{\DUrole{n}{parent}}\sphinxparamcomma \sphinxparam{\DUrole{n}{h}}\sphinxparamcomma \sphinxparam{\DUrole{n}{s}}\sphinxparamcomma \sphinxparam{\DUrole{n}{l}}}
{}
\pysigstopsignatures
\sphinxAtStartPar
Bases: \sphinxcode{\sphinxupquote{QDialog}}
\index{restore\_defaults() (dialogs.ThemeDialog method)@\spxentry{restore\_defaults()}\spxextra{dialogs.ThemeDialog method}}

\begin{fulllineitems}
\phantomsection\label{\detokenize{api/dialogs:dialogs.ThemeDialog.restore_defaults}}
\pysigstartsignatures
\pysiglinewithargsret
{\sphinxbfcode{\sphinxupquote{restore\_defaults}}}
{}
{}
\pysigstopsignatures
\sphinxAtStartPar
Restaura os sliders para os valores padrão do tema.
\begin{quote}\begin{description}
\sphinxlineitem{Return type}
\sphinxAtStartPar
\DUrole{sphinx_autodoc_typehints-type}{\sphinxcode{\sphinxupquote{None}}}

\end{description}\end{quote}

\end{fulllineitems}

\index{values() (dialogs.ThemeDialog method)@\spxentry{values()}\spxextra{dialogs.ThemeDialog method}}

\begin{fulllineitems}
\phantomsection\label{\detokenize{api/dialogs:dialogs.ThemeDialog.values}}
\pysigstartsignatures
\pysiglinewithargsret
{\sphinxbfcode{\sphinxupquote{values}}}
{}
{}
\pysigstopsignatures
\sphinxAtStartPar
Retorna os valores atuais dos sliders de cor.
\begin{quote}\begin{description}
\sphinxlineitem{Returns}
\sphinxAtStartPar
Valores do tom, saturação e brilho.

\sphinxlineitem{Return type}
\sphinxAtStartPar
\DUrole{sphinx_autodoc_typehints-type}{\sphinxcode{\sphinxupquote{Tuple}}{[}\sphinxcode{\sphinxupquote{int}}, \sphinxcode{\sphinxupquote{int}}, \sphinxcode{\sphinxupquote{int}}{]}}

\end{description}\end{quote}

\end{fulllineitems}


\end{fulllineitems}

\index{CursorStartLineEdit (class in dialogs)@\spxentry{CursorStartLineEdit}\spxextra{class in dialogs}}

\begin{fulllineitems}
\phantomsection\label{\detokenize{api/dialogs:dialogs.CursorStartLineEdit}}
\pysigstartsignatures
\pysiglinewithargsret
{\sphinxbfcode{\sphinxupquote{\DUrole{k}{class}\DUrole{w}{ }}}\sphinxcode{\sphinxupquote{dialogs.}}\sphinxbfcode{\sphinxupquote{CursorStartLineEdit}}}
{\sphinxparam{\DUrole{o}{*}\DUrole{n}{args}}\sphinxparamcomma \sphinxparam{\DUrole{n}{force\_cursor\_start}\DUrole{o}{=}\DUrole{default_value}{False}}\sphinxparamcomma \sphinxparam{\DUrole{o}{**}\DUrole{n}{kwargs}}}
{}
\pysigstopsignatures
\sphinxAtStartPar
Bases: \sphinxcode{\sphinxupquote{QLineEdit}}

\sphinxAtStartPar
QLineEdit que força o cursor para o início ao receber foco/clique.
Útil para campos com inputMask (telefone, data) para facilitar colar.
\index{set\_force\_cursor\_start() (dialogs.CursorStartLineEdit method)@\spxentry{set\_force\_cursor\_start()}\spxextra{dialogs.CursorStartLineEdit method}}

\begin{fulllineitems}
\phantomsection\label{\detokenize{api/dialogs:dialogs.CursorStartLineEdit.set_force_cursor_start}}
\pysigstartsignatures
\pysiglinewithargsret
{\sphinxbfcode{\sphinxupquote{set\_force\_cursor\_start}}}
{\sphinxparam{\DUrole{n}{enabled}}}
{}
\pysigstopsignatures
\sphinxAtStartPar
Define se o cursor será forçado para o início ao receber foco ou clique.
\begin{quote}\begin{description}
\sphinxlineitem{Parameters}
\sphinxAtStartPar
\sphinxstyleliteralstrong{\sphinxupquote{enabled}} (\DUrole{sphinx_autodoc_typehints-type}{\sphinxcode{\sphinxupquote{bool}}}) \textendash{} Se True, o cursor será forçado para o início ao receber foco ou clique.

\sphinxlineitem{Return type}
\sphinxAtStartPar
\DUrole{sphinx_autodoc_typehints-type}{\sphinxcode{\sphinxupquote{None}}}

\end{description}\end{quote}

\end{fulllineitems}

\index{focusInEvent() (dialogs.CursorStartLineEdit method)@\spxentry{focusInEvent()}\spxextra{dialogs.CursorStartLineEdit method}}

\begin{fulllineitems}
\phantomsection\label{\detokenize{api/dialogs:dialogs.CursorStartLineEdit.focusInEvent}}
\pysigstartsignatures
\pysiglinewithargsret
{\sphinxbfcode{\sphinxupquote{focusInEvent}}}
{\sphinxparam{\DUrole{n}{event}}}
{}
\pysigstopsignatures
\sphinxAtStartPar
Sobrescreve o evento de foco para mover o cursor para o início se a opção estiver habilitada.
\begin{quote}\begin{description}
\sphinxlineitem{Parameters}
\sphinxAtStartPar
\sphinxstyleliteralstrong{\sphinxupquote{event}} (\sphinxstyleliteralemphasis{\sphinxupquote{\_type\_}}) \textendash{} Evento de foco recebido.

\end{description}\end{quote}

\end{fulllineitems}

\index{mousePressEvent() (dialogs.CursorStartLineEdit method)@\spxentry{mousePressEvent()}\spxextra{dialogs.CursorStartLineEdit method}}

\begin{fulllineitems}
\phantomsection\label{\detokenize{api/dialogs:dialogs.CursorStartLineEdit.mousePressEvent}}
\pysigstartsignatures
\pysiglinewithargsret
{\sphinxbfcode{\sphinxupquote{mousePressEvent}}}
{\sphinxparam{\DUrole{n}{event}}}
{}
\pysigstopsignatures
\sphinxAtStartPar
Sobrescreve o evento de clique do mouse para mover o cursor para o início se a opção estiver habilitada.
:type event: \_type\_
:param event: Evento de clique do mouse recebido.
:type event: \_type\_

\end{fulllineitems}


\end{fulllineitems}

\index{EmailDomainsDialog (class in dialogs)@\spxentry{EmailDomainsDialog}\spxextra{class in dialogs}}

\begin{fulllineitems}
\phantomsection\label{\detokenize{api/dialogs:dialogs.EmailDomainsDialog}}
\pysigstartsignatures
\pysiglinewithargsret
{\sphinxbfcode{\sphinxupquote{\DUrole{k}{class}\DUrole{w}{ }}}\sphinxcode{\sphinxupquote{dialogs.}}\sphinxbfcode{\sphinxupquote{EmailDomainsDialog}}}
{\sphinxparam{\DUrole{n}{parent}}\sphinxparamcomma \sphinxparam{\DUrole{n}{domains}}}
{}
\pysigstopsignatures
\sphinxAtStartPar
Bases: \sphinxcode{\sphinxupquote{QDialog}}
\index{domains() (dialogs.EmailDomainsDialog method)@\spxentry{domains()}\spxextra{dialogs.EmailDomainsDialog method}}

\begin{fulllineitems}
\phantomsection\label{\detokenize{api/dialogs:dialogs.EmailDomainsDialog.domains}}
\pysigstartsignatures
\pysiglinewithargsret
{\sphinxbfcode{\sphinxupquote{domains}}}
{}
{}
\pysigstopsignatures
\sphinxAtStartPar
Retorna a lista de domínios de e\sphinxhyphen{}mail configurados no diálogo, garantindo que cada domínio comece com ‘@’ e seja único.
\begin{quote}\begin{description}
\sphinxlineitem{Returns}
\sphinxAtStartPar
Lista de domínios de e\sphinxhyphen{}mail configurados no diálogo.

\sphinxlineitem{Return type}
\sphinxAtStartPar
\DUrole{sphinx_autodoc_typehints-type}{\sphinxcode{\sphinxupquote{List}}{[}\sphinxcode{\sphinxupquote{str}}{]}}

\end{description}\end{quote}

\end{fulllineitems}


\end{fulllineitems}


\sphinxstepscope


\section{excel\_io module}
\label{\detokenize{api/excel_io:module-excel_io}}\label{\detokenize{api/excel_io:excel-io-module}}\label{\detokenize{api/excel_io::doc}}\index{module@\spxentry{module}!excel\_io@\spxentry{excel\_io}}\index{excel\_io@\spxentry{excel\_io}!module@\spxentry{module}}\index{list\_sheets\_with\_ids() (in module excel\_io)@\spxentry{list\_sheets\_with\_ids()}\spxextra{in module excel\_io}}

\begin{fulllineitems}
\phantomsection\label{\detokenize{api/excel_io:excel_io.list_sheets_with_ids}}
\pysigstartsignatures
\pysiglinewithargsret
{\sphinxcode{\sphinxupquote{excel\_io.}}\sphinxbfcode{\sphinxupquote{list\_sheets\_with\_ids}}}
{\sphinxparam{\DUrole{n}{path}}}
{}
\pysigstopsignatures
\sphinxAtStartPar
Lê um arquivo Excel e retorna uma lista de tuplas contendo o ID da planilha e o nome de cada aba presente no arquivo. O método utiliza a biblioteca openpyxl para carregar o arquivo e iterar sobre as planilhas, extraindo o ID (que pode ser obtido de diferentes atributos dependendo da versão do openpyxl) e o título de cada aba, retornando\sphinxhyphen{}os em uma lista estruturada.
:param path: Caminho do arquivo Excel a ser lido.
:type path: \DUrole{sphinx_autodoc_typehints-type}{\sphinxcode{\sphinxupquote{str}}}
:returns: Lista de tuplas, onde cada tupla contém o ID da planilha (int) e o nome da aba (str) presente no arquivo Excel.
\begin{quote}\begin{description}
\sphinxlineitem{Return type}
\sphinxAtStartPar
\DUrole{sphinx_autodoc_typehints-type}{\sphinxcode{\sphinxupquote{List}}{[}\sphinxcode{\sphinxupquote{Tuple}}{[}\sphinxcode{\sphinxupquote{int}}, \sphinxcode{\sphinxupquote{str}}{]}{]}}

\end{description}\end{quote}

\end{fulllineitems}

\index{list\_sheets() (in module excel\_io)@\spxentry{list\_sheets()}\spxextra{in module excel\_io}}

\begin{fulllineitems}
\phantomsection\label{\detokenize{api/excel_io:excel_io.list_sheets}}
\pysigstartsignatures
\pysiglinewithargsret
{\sphinxcode{\sphinxupquote{excel\_io.}}\sphinxbfcode{\sphinxupquote{list\_sheets}}}
{\sphinxparam{\DUrole{n}{path}}}
{}
\pysigstopsignatures\begin{quote}\begin{description}
\sphinxlineitem{Return type}
\sphinxAtStartPar
\DUrole{sphinx_autodoc_typehints-type}{\sphinxcode{\sphinxupquote{list}}{[}\sphinxcode{\sphinxupquote{str}}{]}}

\end{description}\end{quote}

\end{fulllineitems}

\index{ensure\_sheet\_exists() (in module excel\_io)@\spxentry{ensure\_sheet\_exists()}\spxextra{in module excel\_io}}

\begin{fulllineitems}
\phantomsection\label{\detokenize{api/excel_io:excel_io.ensure_sheet_exists}}
\pysigstartsignatures
\pysiglinewithargsret
{\sphinxcode{\sphinxupquote{excel\_io.}}\sphinxbfcode{\sphinxupquote{ensure\_sheet\_exists}}}
{\sphinxparam{\DUrole{n}{path}}\sphinxparamcomma \sphinxparam{\DUrole{n}{sheet\_name}}}
{}
\pysigstopsignatures\begin{quote}\begin{description}
\sphinxlineitem{Return type}
\sphinxAtStartPar
\DUrole{sphinx_autodoc_typehints-type}{\sphinxcode{\sphinxupquote{None}}}

\end{description}\end{quote}

\end{fulllineitems}

\index{is\_excel\_lock\_present() (in module excel\_io)@\spxentry{is\_excel\_lock\_present()}\spxextra{in module excel\_io}}

\begin{fulllineitems}
\phantomsection\label{\detokenize{api/excel_io:excel_io.is_excel_lock_present}}
\pysigstartsignatures
\pysiglinewithargsret
{\sphinxcode{\sphinxupquote{excel\_io.}}\sphinxbfcode{\sphinxupquote{is\_excel\_lock\_present}}}
{\sphinxparam{\DUrole{n}{xlsx\_path}}}
{}
\pysigstopsignatures\begin{quote}\begin{description}
\sphinxlineitem{Return type}
\sphinxAtStartPar
\DUrole{sphinx_autodoc_typehints-type}{\sphinxcode{\sphinxupquote{bool}}}

\end{description}\end{quote}

\end{fulllineitems}

\index{get\_sheet\_headers() (in module excel\_io)@\spxentry{get\_sheet\_headers()}\spxextra{in module excel\_io}}

\begin{fulllineitems}
\phantomsection\label{\detokenize{api/excel_io:excel_io.get_sheet_headers}}
\pysigstartsignatures
\pysiglinewithargsret
{\sphinxcode{\sphinxupquote{excel\_io.}}\sphinxbfcode{\sphinxupquote{get\_sheet\_headers}}}
{\sphinxparam{\DUrole{n}{xlsx\_path}}\sphinxparamcomma \sphinxparam{\DUrole{n}{sheet\_name}}}
{}
\pysigstopsignatures\begin{quote}\begin{description}
\sphinxlineitem{Return type}
\sphinxAtStartPar
\DUrole{sphinx_autodoc_typehints-type}{\sphinxcode{\sphinxupquote{list}}{[}\sphinxcode{\sphinxupquote{str}}{]}}

\end{description}\end{quote}

\end{fulllineitems}

\index{read\_headers\_from\_excel() (in module excel\_io)@\spxentry{read\_headers\_from\_excel()}\spxextra{in module excel\_io}}

\begin{fulllineitems}
\phantomsection\label{\detokenize{api/excel_io:excel_io.read_headers_from_excel}}
\pysigstartsignatures
\pysiglinewithargsret
{\sphinxcode{\sphinxupquote{excel\_io.}}\sphinxbfcode{\sphinxupquote{read\_headers\_from\_excel}}}
{\sphinxparam{\DUrole{n}{path}}\sphinxparamcomma \sphinxparam{\DUrole{n}{preferred\_sheet}}}
{}
\pysigstopsignatures
\sphinxAtStartPar
Lê a linha 1 e retorna (sheet\_name\_usado, headers, has\_header).
has\_header = True se encontrou pelo menos 1 valor não vazio na linha 1.
\sphinxhyphen{} Se preferred\_sheet não existir, usa a aba ativa.
\begin{quote}\begin{description}
\sphinxlineitem{Return type}
\sphinxAtStartPar
\DUrole{sphinx_autodoc_typehints-type}{\sphinxcode{\sphinxupquote{Tuple}}{[}\sphinxcode{\sphinxupquote{str}}, \sphinxcode{\sphinxupquote{List}}{[}\sphinxcode{\sphinxupquote{str}}{]}, \sphinxcode{\sphinxupquote{bool}}{]}}

\end{description}\end{quote}

\end{fulllineitems}

\index{write\_headers\_from\_campos() (in module excel\_io)@\spxentry{write\_headers\_from\_campos()}\spxextra{in module excel\_io}}

\begin{fulllineitems}
\phantomsection\label{\detokenize{api/excel_io:excel_io.write_headers_from_campos}}
\pysigstartsignatures
\pysiglinewithargsret
{\sphinxcode{\sphinxupquote{excel\_io.}}\sphinxbfcode{\sphinxupquote{write\_headers\_from\_campos}}}
{\sphinxparam{\DUrole{n}{path}}\sphinxparamcomma \sphinxparam{\DUrole{n}{sheet\_name}}\sphinxparamcomma \sphinxparam{\DUrole{n}{campos}}}
{}
\pysigstopsignatures
\sphinxAtStartPar
Garante que a linha 1 tenha cabeçalhos (títulos) baseados em campos do app.
Se a aba não existir, cria.
\begin{quote}\begin{description}
\sphinxlineitem{Return type}
\sphinxAtStartPar
\DUrole{sphinx_autodoc_typehints-type}{\sphinxcode{\sphinxupquote{None}}}

\end{description}\end{quote}

\end{fulllineitems}

\index{ensure\_workbook() (in module excel\_io)@\spxentry{ensure\_workbook()}\spxextra{in module excel\_io}}

\begin{fulllineitems}
\phantomsection\label{\detokenize{api/excel_io:excel_io.ensure_workbook}}
\pysigstartsignatures
\pysiglinewithargsret
{\sphinxcode{\sphinxupquote{excel\_io.}}\sphinxbfcode{\sphinxupquote{ensure\_workbook}}}
{\sphinxparam{\DUrole{n}{path}}\sphinxparamcomma \sphinxparam{\DUrole{n}{sheet\_name}}\sphinxparamcomma \sphinxparam{\DUrole{n}{headers}}}
{}
\pysigstopsignatures
\sphinxAtStartPar
Garante que o arquivo Excel exista, que a aba exista e que a linha 1 tenha os cabeçalhos (títulos) especificados.
:param path: Caminho do arquivo Excel.
:type path: \DUrole{sphinx_autodoc_typehints-type}{\sphinxcode{\sphinxupquote{str}}}
:param sheet\_name: Nome da aba onde os dados serão escritos.
:type sheet\_name: \DUrole{sphinx_autodoc_typehints-type}{\sphinxcode{\sphinxupquote{str}}}
:param headers: Lista de cabeçalhos (títulos) a serem garantidos na linha 1 da aba. Se a aba já existir, os cabeçalhos serão mesclados com os existentes, mantendo a ordem dos novos e sem duplicatas.
:type headers: \DUrole{sphinx_autodoc_typehints-type}{\sphinxcode{\sphinxupquote{List}}{[}\sphinxcode{\sphinxupquote{str}}{]}}
\begin{quote}\begin{description}
\sphinxlineitem{Return type}
\sphinxAtStartPar
\DUrole{sphinx_autodoc_typehints-type}{\sphinxcode{\sphinxupquote{None}}}

\end{description}\end{quote}

\end{fulllineitems}

\index{apply\_column\_type\_rules() (in module excel\_io)@\spxentry{apply\_column\_type\_rules()}\spxextra{in module excel\_io}}

\begin{fulllineitems}
\phantomsection\label{\detokenize{api/excel_io:excel_io.apply_column_type_rules}}
\pysigstartsignatures
\pysiglinewithargsret
{\sphinxcode{\sphinxupquote{excel\_io.}}\sphinxbfcode{\sphinxupquote{apply\_column\_type\_rules}}}
{\sphinxparam{\DUrole{n}{path}}\sphinxparamcomma \sphinxparam{\DUrole{n}{sheet\_name}}\sphinxparamcomma \sphinxparam{\DUrole{n}{campos}}}
{}
\pysigstopsignatures
\sphinxAtStartPar
Aplica regras de formatação e validação de dados nas colunas do Excel com base nos tipos definidos nos campos do app.
\begin{quote}\begin{description}
\sphinxlineitem{Return type}
\sphinxAtStartPar
\DUrole{sphinx_autodoc_typehints-type}{\sphinxcode{\sphinxupquote{None}}}

\end{description}\end{quote}

\end{fulllineitems}

\index{parse\_br\_number() (in module excel\_io)@\spxentry{parse\_br\_number()}\spxextra{in module excel\_io}}

\begin{fulllineitems}
\phantomsection\label{\detokenize{api/excel_io:excel_io.parse_br_number}}
\pysigstartsignatures
\pysiglinewithargsret
{\sphinxcode{\sphinxupquote{excel\_io.}}\sphinxbfcode{\sphinxupquote{parse\_br\_number}}}
{\sphinxparam{\DUrole{n}{value}}}
{}
\pysigstopsignatures
\sphinxAtStartPar
Converte texto numérico para int ou float, aceitando padrão brasileiro.

\sphinxAtStartPar
Exemplos aceitos:
\sphinxhyphen{} 10
\sphinxhyphen{} 10,5
\sphinxhyphen{} 1.234,56

\sphinxAtStartPar
Retorna:
\sphinxhyphen{} int, se for inteiro
\sphinxhyphen{} float, se tiver parte decimal
\sphinxhyphen{} None, se não for número válido

\end{fulllineitems}

\index{parse\_br\_date\_or\_text() (in module excel\_io)@\spxentry{parse\_br\_date\_or\_text()}\spxextra{in module excel\_io}}

\begin{fulllineitems}
\phantomsection\label{\detokenize{api/excel_io:excel_io.parse_br_date_or_text}}
\pysigstartsignatures
\pysiglinewithargsret
{\sphinxcode{\sphinxupquote{excel\_io.}}\sphinxbfcode{\sphinxupquote{parse\_br\_date\_or\_text}}}
{\sphinxparam{\DUrole{n}{value}}}
{}
\pysigstopsignatures
\sphinxAtStartPar
Retorna datetime se a data for válida e suportada pelo Excel.
Se for anterior a 1900, retorna o texto original.
Se não for uma data válida, retorna o texto original.

\end{fulllineitems}

\index{append\_row\_typed() (in module excel\_io)@\spxentry{append\_row\_typed()}\spxextra{in module excel\_io}}

\begin{fulllineitems}
\phantomsection\label{\detokenize{api/excel_io:excel_io.append_row_typed}}
\pysigstartsignatures
\pysiglinewithargsret
{\sphinxcode{\sphinxupquote{excel\_io.}}\sphinxbfcode{\sphinxupquote{append\_row\_typed}}}
{\sphinxparam{\DUrole{n}{file\_path}}\sphinxparamcomma \sphinxparam{\DUrole{n}{sheet\_name}}\sphinxparamcomma \sphinxparam{\DUrole{n}{campos}}\sphinxparamcomma \sphinxparam{\DUrole{n}{row\_by\_title}}}
{}
\pysigstopsignatures
\sphinxAtStartPar
Adiciona uma nova linha na planilha, respeitando o tipo de cada campo.

\sphinxAtStartPar
Regras principais:
\sphinxhyphen{} texto/email/telefone/… \sphinxhyphen{}\textgreater{} grava como texto
\sphinxhyphen{} booleano \sphinxhyphen{}\textgreater{} grava como texto
\sphinxhyphen{} numero:
\begin{itemize}
\item {} 
\sphinxAtStartPar
grava como int/float

\item {} 
\sphinxAtStartPar
sem forçar casas decimais

\end{itemize}
\begin{itemize}
\item {} \begin{description}
\sphinxlineitem{data:}\begin{itemize}
\item {} 
\sphinxAtStartPar
se ano \textgreater{}= 1900, grava como datetime com formato DD/MM/YYYY

\item {} 
\sphinxAtStartPar
se ano \textless{} 1900, grava como texto

\end{itemize}

\end{description}

\item {} 
\sphinxAtStartPar
campos vazios \sphinxhyphen{}\textgreater{} célula vazia

\end{itemize}
\begin{quote}\begin{description}
\sphinxlineitem{Return type}
\sphinxAtStartPar
\DUrole{sphinx_autodoc_typehints-type}{\sphinxcode{\sphinxupquote{None}}}

\end{description}\end{quote}

\end{fulllineitems}

\index{delete\_column\_by\_header() (in module excel\_io)@\spxentry{delete\_column\_by\_header()}\spxextra{in module excel\_io}}

\begin{fulllineitems}
\phantomsection\label{\detokenize{api/excel_io:excel_io.delete_column_by_header}}
\pysigstartsignatures
\pysiglinewithargsret
{\sphinxcode{\sphinxupquote{excel\_io.}}\sphinxbfcode{\sphinxupquote{delete\_column\_by\_header}}}
{\sphinxparam{\DUrole{n}{path}}\sphinxparamcomma \sphinxparam{\DUrole{n}{sheet\_name}}\sphinxparamcomma \sphinxparam{\DUrole{n}{header\_name}}}
{}
\pysigstopsignatures
\sphinxAtStartPar
Exclui uma coluna do arquivo Excel com base no nome do cabeçalho. O método verifica se o arquivo e a aba existem, se a linha 1 contém o cabeçalho especificado e, se todas as condições forem atendidas, exclui a coluna correspondente ao cabeçalho fornecido.
:param path: Caminho do arquivo Excel.
:type path: \DUrole{sphinx_autodoc_typehints-type}{\sphinxcode{\sphinxupquote{str}}}
:param sheet\_name: Nome da aba onde a coluna será excluída.
:type sheet\_name: \DUrole{sphinx_autodoc_typehints-type}{\sphinxcode{\sphinxupquote{str}}}
:param header\_name: Nome do cabeçalho da coluna a ser excluída.
:type header\_name: \DUrole{sphinx_autodoc_typehints-type}{\sphinxcode{\sphinxupquote{str}}}
\begin{quote}\begin{description}
\sphinxlineitem{Return type}
\sphinxAtStartPar
\DUrole{sphinx_autodoc_typehints-type}{\sphinxcode{\sphinxupquote{None}}}

\end{description}\end{quote}

\end{fulllineitems}

\index{rename\_column\_header() (in module excel\_io)@\spxentry{rename\_column\_header()}\spxextra{in module excel\_io}}

\begin{fulllineitems}
\phantomsection\label{\detokenize{api/excel_io:excel_io.rename_column_header}}
\pysigstartsignatures
\pysiglinewithargsret
{\sphinxcode{\sphinxupquote{excel\_io.}}\sphinxbfcode{\sphinxupquote{rename\_column\_header}}}
{\sphinxparam{\DUrole{n}{path}}\sphinxparamcomma \sphinxparam{\DUrole{n}{sheet\_name}}\sphinxparamcomma \sphinxparam{\DUrole{n}{old\_header}}\sphinxparamcomma \sphinxparam{\DUrole{n}{new\_header}}}
{}
\pysigstopsignatures
\sphinxAtStartPar
Renomeia o cabeçalho de uma coluna no arquivo Excel. O método verifica se o arquivo e a aba existem, se a linha 1 contém o cabeçalho antigo especificado e, se todas as condições forem atendidas, atualiza o valor do cabeçalho para o novo nome fornecido.
:param path: Caminho do arquivo Excel.
:type path: \DUrole{sphinx_autodoc_typehints-type}{\sphinxcode{\sphinxupquote{str}}}
:param sheet\_name: Nome da aba onde a coluna será renomeada.
:type sheet\_name: \DUrole{sphinx_autodoc_typehints-type}{\sphinxcode{\sphinxupquote{str}}}
:param old\_header: Nome do cabeçalho atual da coluna a ser renomeada.
:type old\_header: \DUrole{sphinx_autodoc_typehints-type}{\sphinxcode{\sphinxupquote{str}}}
:param new\_header: Novo nome do cabeçalho da coluna.
:type new\_header: \DUrole{sphinx_autodoc_typehints-type}{\sphinxcode{\sphinxupquote{str}}}
\begin{quote}\begin{description}
\sphinxlineitem{Return type}
\sphinxAtStartPar
\DUrole{sphinx_autodoc_typehints-type}{\sphinxcode{\sphinxupquote{None}}}

\end{description}\end{quote}

\end{fulllineitems}


\sphinxstepscope


\section{formatters module}
\label{\detokenize{api/formatters:module-formatters}}\label{\detokenize{api/formatters:formatters-module}}\label{\detokenize{api/formatters::doc}}\index{module@\spxentry{module}!formatters@\spxentry{formatters}}\index{formatters@\spxentry{formatters}!module@\spxentry{module}}\index{normalize\_masked\_text() (in module formatters)@\spxentry{normalize\_masked\_text()}\spxextra{in module formatters}}

\begin{fulllineitems}
\phantomsection\label{\detokenize{api/formatters:formatters.normalize_masked_text}}
\pysigstartsignatures
\pysiglinewithargsret
{\sphinxcode{\sphinxupquote{formatters.}}\sphinxbfcode{\sphinxupquote{normalize\_masked\_text}}}
{\sphinxparam{\DUrole{n}{value}}}
{}
\pysigstopsignatures
\sphinxAtStartPar
Normaliza um texto que pode conter caracteres de máscara (como em campos de formulário) e retorna uma versão limpa e formatada do texto. O método remove caracteres de máscara comuns, como sublinhados, espaços, hífens, pontos, barras e parênteses, e verifica se o texto resultante contém apenas caracteres de separação. Se o texto contiver apenas caracteres de separação sem letras ou dígitos, ele retorna uma string vazia. Além disso, se o texto contiver uma barra sem dígitos, ou se tiver dígitos mas menos de 8 quando houver uma barra, ele também retorna uma string vazia. Caso contrário, retorna o texto limpo e formatado.
:param value: O texto a ser normalizado.
:type value: \DUrole{sphinx_autodoc_typehints-type}{\sphinxcode{\sphinxupquote{str}}}
:returns: O texto normalizado, sem caracteres de máscara e formatado de acordo com as regras descritas.
\begin{quote}\begin{description}
\sphinxlineitem{Return type}
\sphinxAtStartPar
\DUrole{sphinx_autodoc_typehints-type}{\sphinxcode{\sphinxupquote{str}}}

\end{description}\end{quote}

\end{fulllineitems}

\index{digits\_only() (in module formatters)@\spxentry{digits\_only()}\spxextra{in module formatters}}

\begin{fulllineitems}
\phantomsection\label{\detokenize{api/formatters:formatters.digits_only}}
\pysigstartsignatures
\pysiglinewithargsret
{\sphinxcode{\sphinxupquote{formatters.}}\sphinxbfcode{\sphinxupquote{digits\_only}}}
{\sphinxparam{\DUrole{n}{value}}}
{}
\pysigstopsignatures
\sphinxAtStartPar
Remove todos os caracteres não numéricos de uma string, retornando apenas os dígitos. O método utiliza uma expressão regular para substituir qualquer caractere que não seja um dígito por uma string vazia, resultando em uma string composta apenas por números. Se a entrada for None ou vazia, o método retorna uma string vazia.
:param value: A string de entrada a ser processada.
:type value: \DUrole{sphinx_autodoc_typehints-type}{\sphinxcode{\sphinxupquote{str}}}
:returns: A string contendo apenas os dígitos da entrada original.
\begin{quote}\begin{description}
\sphinxlineitem{Return type}
\sphinxAtStartPar
\DUrole{sphinx_autodoc_typehints-type}{\sphinxcode{\sphinxupquote{str}}}

\end{description}\end{quote}

\end{fulllineitems}

\index{strip\_mask\_chars() (in module formatters)@\spxentry{strip\_mask\_chars()}\spxextra{in module formatters}}

\begin{fulllineitems}
\phantomsection\label{\detokenize{api/formatters:formatters.strip_mask_chars}}
\pysigstartsignatures
\pysiglinewithargsret
{\sphinxcode{\sphinxupquote{formatters.}}\sphinxbfcode{\sphinxupquote{strip\_mask\_chars}}}
{\sphinxparam{\DUrole{n}{value}}}
{}
\pysigstopsignatures
\sphinxAtStartPar
Remove caracteres de máscara comuns de uma string, como sublinhados, espaços, hífens, pontos, barras e parênteses, e retorna a string limpa. O método utiliza uma expressão regular para substituir esses caracteres por uma string vazia, resultando em uma versão da string sem os caracteres de máscara. Se a entrada for None ou vazia, o método retorna uma string vazia.
:param value: A string de entrada a ser processada, que pode conter caracteres de máscara.
:type value: \DUrole{sphinx_autodoc_typehints-type}{\sphinxcode{\sphinxupquote{str}}}
:returns: A string limpa, sem caracteres de máscara.
\begin{quote}\begin{description}
\sphinxlineitem{Return type}
\sphinxAtStartPar
\DUrole{sphinx_autodoc_typehints-type}{\sphinxcode{\sphinxupquote{str}}}

\end{description}\end{quote}

\end{fulllineitems}

\index{parse\_decimal\_br() (in module formatters)@\spxentry{parse\_decimal\_br()}\spxextra{in module formatters}}

\begin{fulllineitems}
\phantomsection\label{\detokenize{api/formatters:formatters.parse_decimal_br}}
\pysigstartsignatures
\pysiglinewithargsret
{\sphinxcode{\sphinxupquote{formatters.}}\sphinxbfcode{\sphinxupquote{parse\_decimal\_br}}}
{\sphinxparam{\DUrole{n}{s}}}
{}
\pysigstopsignatures
\sphinxAtStartPar
Converte uma string que representa um número no formato brasileiro (com vírgula como separador decimal) para um valor float. O método remove quaisquer caracteres que não sejam dígitos, vírgulas, pontos ou sinais de menos, e então processa a string para garantir que o formato seja interpretado corretamente. Se a string contiver tanto vírgula quanto ponto, o método assume que o ponto é um separador de milhar e a vírgula é o separador decimal. Se a string contiver apenas vírgula, ela é tratada como separador decimal. O método tenta converter a string resultante para float e retorna o valor convertido, ou None se a conversão falhar.
\begin{quote}\begin{description}
\sphinxlineitem{Parameters}
\sphinxAtStartPar
\sphinxstyleliteralstrong{\sphinxupquote{s}} (\DUrole{sphinx_autodoc_typehints-type}{\sphinxcode{\sphinxupquote{str}}}) \textendash{} A string contendo um valor numérico no formato brasileiro (com vírgula como separador decimal).

\sphinxlineitem{Returns}
\sphinxAtStartPar
O valor numérico convertido para float, ou None se a conversão falhar.

\sphinxlineitem{Return type}
\sphinxAtStartPar
\DUrole{sphinx_autodoc_typehints-type}{\sphinxcode{\sphinxupquote{Optional}}{[}\sphinxcode{\sphinxupquote{float}}{]}}

\end{description}\end{quote}

\end{fulllineitems}

\index{parse\_date\_any() (in module formatters)@\spxentry{parse\_date\_any()}\spxextra{in module formatters}}

\begin{fulllineitems}
\phantomsection\label{\detokenize{api/formatters:formatters.parse_date_any}}
\pysigstartsignatures
\pysiglinewithargsret
{\sphinxcode{\sphinxupquote{formatters.}}\sphinxbfcode{\sphinxupquote{parse\_date\_any}}}
{\sphinxparam{\DUrole{n}{s}}}
{}
\pysigstopsignatures
\sphinxAtStartPar
Converte uma string que representa uma data em um objeto date do Python, tentando vários formatos comuns de data. O método aceita strings em formatos como dd/mm/yyyy, dd\sphinxhyphen{}mm\sphinxhyphen{}yyyy, yyyy\sphinxhyphen{}mm\sphinxhyphen{}dd e yyyy/mm/dd. Ele tenta converter a string para um objeto date usando cada um desses formatos, e retorna o primeiro resultado bem\sphinxhyphen{}sucedido. Se a conversão falhar para todos os formatos, o método retorna None. Se a entrada for None ou vazia, ele também retorna None.
\begin{quote}\begin{description}
\sphinxlineitem{Parameters}
\sphinxAtStartPar
\sphinxstyleliteralstrong{\sphinxupquote{s}} (\DUrole{sphinx_autodoc_typehints-type}{\sphinxcode{\sphinxupquote{str}}}) \textendash{} A string contendo uma data no formato dd/mm/yyyy, dd\sphinxhyphen{}mm\sphinxhyphen{}yyyy, yyyy\sphinxhyphen{}mm\sphinxhyphen{}dd ou yyyy/mm/dd.

\sphinxlineitem{Returns}
\sphinxAtStartPar
A data convertida, ou None se a conversão falhar.

\sphinxlineitem{Return type}
\sphinxAtStartPar
\DUrole{sphinx_autodoc_typehints-type}{\sphinxcode{\sphinxupquote{Optional}}{[}\sphinxcode{\sphinxupquote{date}}{]}}

\end{description}\end{quote}

\end{fulllineitems}

\index{normalize\_bool\_ptbr() (in module formatters)@\spxentry{normalize\_bool\_ptbr()}\spxextra{in module formatters}}

\begin{fulllineitems}
\phantomsection\label{\detokenize{api/formatters:formatters.normalize_bool_ptbr}}
\pysigstartsignatures
\pysiglinewithargsret
{\sphinxcode{\sphinxupquote{formatters.}}\sphinxbfcode{\sphinxupquote{normalize\_bool\_ptbr}}}
{\sphinxparam{\DUrole{n}{s}}}
{}
\pysigstopsignatures
\sphinxAtStartPar
Normaliza uma string que representa um valor booleano em português, retornando “Sim” para valores que indicam verdadeiro e “Não” para valores que indicam falso. O método aceita várias formas de entrada, como “sim”, “s”, “yes”, “y”, “true”, “1” para verdadeiro, e “não”, “nao”, “n”, “no”, “false”, “0” para falso. A comparação é feita de forma case\sphinxhyphen{}insensitive e a string de entrada é limpa de espaços antes da comparação. Se a string não corresponder a nenhum dos valores reconhecidos para verdadeiro ou falso, o método retorna a string original limpa.
:param s: A string de entrada que representa um valor booleano em português.
:type s: \DUrole{sphinx_autodoc_typehints-type}{\sphinxcode{\sphinxupquote{str}}}
:returns: “Sim” se a string indicar verdadeiro, “Não” se indicar falso, ou a string original limpa se não corresponder a nenhum dos casos anteriores. Retorna None se a entrada for None ou vazia.
\begin{quote}\begin{description}
\sphinxlineitem{Return type}
\sphinxAtStartPar
\DUrole{sphinx_autodoc_typehints-type}{\sphinxcode{\sphinxupquote{Optional}}{[}\sphinxcode{\sphinxupquote{str}}{]}}

\end{description}\end{quote}

\end{fulllineitems}

\index{cast\_value\_by\_type() (in module formatters)@\spxentry{cast\_value\_by\_type()}\spxextra{in module formatters}}

\begin{fulllineitems}
\phantomsection\label{\detokenize{api/formatters:formatters.cast_value_by_type}}
\pysigstartsignatures
\pysiglinewithargsret
{\sphinxcode{\sphinxupquote{formatters.}}\sphinxbfcode{\sphinxupquote{cast\_value\_by\_type}}}
{\sphinxparam{\DUrole{n}{tipo}}\sphinxparamcomma \sphinxparam{\DUrole{n}{raw}}}
{}
\pysigstopsignatures
\sphinxAtStartPar
Converte uma string bruta para um valor do tipo especificado, com suporte para tipos como “numero”, “moeda”, “data” e “booleano”. O método processa a string de entrada com base no tipo solicitado, utilizando funções de formatação específicas para cada tipo. Para “numero” e “moeda”, ele tenta converter a string para um float usando o formato brasileiro. Para “data”, ele tenta converter a string para um objeto date. Para “booleano”, ele normaliza a string para “Sim” ou “Não”. Se a conversão falhar ou se o tipo não for reconhecido, o método retorna a string original limpa.
:param tipo: O tipo para o qual a string deve ser convertida, como “numero”, “moeda”, “data” ou “booleano”.
:type tipo: \DUrole{sphinx_autodoc_typehints-type}{\sphinxcode{\sphinxupquote{str}}}
:param raw: A string bruta a ser convertida para o tipo especificado. O método irá processar essa string de acordo com as regras de formatação para o tipo solicitado e retornar o valor convertido, ou a string original limpa se a conversão falhar ou se o tipo não for reconhecido.
:type raw: \DUrole{sphinx_autodoc_typehints-type}{\sphinxcode{\sphinxupquote{str}}}
:returns: O valor convertido para o tipo especificado, ou a string original limpa se a conversão falhar ou se o tipo não for reconhecido. Para “numero” e “moeda”, retorna um float; para “data”, retorna um objeto date; para “booleano”, retorna “Sim” ou “Não”.
\begin{quote}\begin{description}
\sphinxlineitem{Return type}
\sphinxAtStartPar
\DUrole{sphinx_autodoc_typehints-type}{\sphinxcode{\sphinxupquote{Any}}}

\end{description}\end{quote}

\end{fulllineitems}


\sphinxstepscope


\section{icon\_manager module}
\label{\detokenize{api/icon_manager:module-icon_manager}}\label{\detokenize{api/icon_manager:icon-manager-module}}\label{\detokenize{api/icon_manager::doc}}\index{module@\spxentry{module}!icon\_manager@\spxentry{icon\_manager}}\index{icon\_manager@\spxentry{icon\_manager}!module@\spxentry{module}}\index{IconManager (class in icon\_manager)@\spxentry{IconManager}\spxextra{class in icon\_manager}}

\begin{fulllineitems}
\phantomsection\label{\detokenize{api/icon_manager:icon_manager.IconManager}}
\pysigstartsignatures
\pysiglinewithargsret
{\sphinxbfcode{\sphinxupquote{\DUrole{k}{class}\DUrole{w}{ }}}\sphinxcode{\sphinxupquote{icon\_manager.}}\sphinxbfcode{\sphinxupquote{IconManager}}}
{\sphinxparam{\DUrole{n}{ui\_cfg}}\sphinxparamcomma \sphinxparam{\DUrole{n}{abs\_path}}\sphinxparamcomma \sphinxparam{\DUrole{n}{file\_exists}}}
{}
\pysigstopsignatures
\sphinxAtStartPar
Bases: \sphinxcode{\sphinxupquote{object}}

\sphinxAtStartPar
Centraliza:
\sphinxhyphen{} Resolução de caminho (ex.: PyInstaller \_MEIPASS via resolver externo)
\sphinxhyphen{} Cache de QPixmap
\sphinxhyphen{} Tint dos ícones (SourceIn)
\sphinxhyphen{} Aplicação em botões por key (lido do ui\_cfg{[}“button\_icons”{]})
\index{icon\_for\_key() (icon\_manager.IconManager method)@\spxentry{icon\_for\_key()}\spxextra{icon\_manager.IconManager method}}

\begin{fulllineitems}
\phantomsection\label{\detokenize{api/icon_manager:icon_manager.IconManager.icon_for_key}}
\pysigstartsignatures
\pysiglinewithargsret
{\sphinxbfcode{\sphinxupquote{icon\_for\_key}}}
{\sphinxparam{\DUrole{n}{key}}\sphinxparamcomma \sphinxparam{\DUrole{n}{tint\_hex}}}
{}
\pysigstopsignatures
\sphinxAtStartPar
Retorna um QIcon para uma chave de ícone específica, aplicando uma tintura se necessário. O método busca o caminho do ícone correspondente à chave fornecida no dicionário de configuração ui\_cfg{[}“button\_icons”{]}, verifica se o arquivo existe e, se for encontrado, carrega o QPixmap usando o método \_load\_pixmap. Se o QPixmap for carregado com sucesso, ele aplica a tintura usando o método \_tint\_pixmap e retorna um QIcon criado a partir do pixmap tinturado. Se a chave não for encontrada, se o arquivo não existir ou se o QPixmap for inválido, o método retorna None.
:param key: A chave do ícone a ser buscada no dicionário de configuração ui\_cfg{[}“button\_icons”{]} para obter o caminho do arquivo de ícone correspondente.
:type key: \DUrole{sphinx_autodoc_typehints-type}{\sphinxcode{\sphinxupquote{str}}}
:param tint\_hex: A cor hexadecimal da tintura a ser aplicada ao ícone, caso o arquivo de ícone seja encontrado e carregado com sucesso.
:type tint\_hex: \DUrole{sphinx_autodoc_typehints-type}{\sphinxcode{\sphinxupquote{str}}}
:returns: Um QIcon criado a partir do QPixmap tinturado correspondente à chave fornecida, ou None se a chave não for encontrada, se o arquivo de ícone não existir ou se o QPixmap for inválido.
\begin{quote}\begin{description}
\sphinxlineitem{Return type}
\sphinxAtStartPar
\DUrole{sphinx_autodoc_typehints-type}{\sphinxcode{\sphinxupquote{Optional}}{[}\sphinxcode{\sphinxupquote{QIcon}}{]}}

\end{description}\end{quote}

\end{fulllineitems}

\index{apply\_button\_icon() (icon\_manager.IconManager method)@\spxentry{apply\_button\_icon()}\spxextra{icon\_manager.IconManager method}}

\begin{fulllineitems}
\phantomsection\label{\detokenize{api/icon_manager:icon_manager.IconManager.apply_button_icon}}
\pysigstartsignatures
\pysiglinewithargsret
{\sphinxbfcode{\sphinxupquote{apply\_button\_icon}}}
{\sphinxparam{\DUrole{n}{btn}}\sphinxparamcomma \sphinxparam{\DUrole{n}{key}}\sphinxparamcomma \sphinxparam{\DUrole{n}{tint\_hex}}}
{}
\pysigstopsignatures
\sphinxAtStartPar
Aplica um ícone a um QPushButton com base em uma chave de ícone específica, utilizando o método icon\_for\_key para obter o QIcon correspondente. O método busca o QIcon usando a chave fornecida e, se um ícone válido for retornado, ele é aplicado ao botão usando o método setIcon do QPushButton. Se o ícone não for encontrado ou for inválido, o método não faz nenhuma alteração no botão.
:param btn: O botão ao qual o ícone será aplicado.
:type btn: \DUrole{sphinx_autodoc_typehints-type}{\sphinxcode{\sphinxupquote{QPushButton}}}
:param key: A chave do ícone a ser buscada no dicionário de configuração ui\_cfg{[}“button\_icons”{]} para obter o caminho do arquivo de ícone correspondente.
:type key: \DUrole{sphinx_autodoc_typehints-type}{\sphinxcode{\sphinxupquote{str}}}
:param tint\_hex: A cor hexadecimal da tintura a ser aplicada ao ícone, caso o arquivo de ícone seja encontrado e carregado com sucesso.
:type tint\_hex: \DUrole{sphinx_autodoc_typehints-type}{\sphinxcode{\sphinxupquote{str}}}
\begin{quote}\begin{description}
\sphinxlineitem{Return type}
\sphinxAtStartPar
\DUrole{sphinx_autodoc_typehints-type}{\sphinxcode{\sphinxupquote{None}}}

\end{description}\end{quote}

\end{fulllineitems}

\index{clear\_cache() (icon\_manager.IconManager method)@\spxentry{clear\_cache()}\spxextra{icon\_manager.IconManager method}}

\begin{fulllineitems}
\phantomsection\label{\detokenize{api/icon_manager:icon_manager.IconManager.clear_cache}}
\pysigstartsignatures
\pysiglinewithargsret
{\sphinxbfcode{\sphinxupquote{clear\_cache}}}
{}
{}
\pysigstopsignatures
\sphinxAtStartPar
Limpa o cache de QPixmap armazenados no IconManager. O método remove todas as entradas do dicionário \_pixmap\_cache, liberando a memória ocupada pelos QPixmaps armazenados. Isso pode ser útil para liberar recursos quando os ícones não são mais necessários ou quando se deseja forçar o recarregamento dos ícones a partir dos arquivos.
\begin{quote}\begin{description}
\sphinxlineitem{Return type}
\sphinxAtStartPar
\DUrole{sphinx_autodoc_typehints-type}{\sphinxcode{\sphinxupquote{None}}}

\end{description}\end{quote}

\end{fulllineitems}


\end{fulllineitems}


\sphinxstepscope


\section{license\_dialog module}
\label{\detokenize{api/license_dialog:module-license_dialog}}\label{\detokenize{api/license_dialog:license-dialog-module}}\label{\detokenize{api/license_dialog::doc}}\index{module@\spxentry{module}!license\_dialog@\spxentry{license\_dialog}}\index{license\_dialog@\spxentry{license\_dialog}!module@\spxentry{module}}\index{abs\_path() (in module license\_dialog)@\spxentry{abs\_path()}\spxextra{in module license\_dialog}}

\begin{fulllineitems}
\phantomsection\label{\detokenize{api/license_dialog:license_dialog.abs_path}}
\pysigstartsignatures
\pysiglinewithargsret
{\sphinxcode{\sphinxupquote{license\_dialog.}}\sphinxbfcode{\sphinxupquote{abs\_path}}}
{\sphinxparam{\DUrole{n}{p}}}
{}
\pysigstopsignatures
\sphinxAtStartPar
Resolve um caminho de arquivo relativo para um caminho absoluto, considerando o diretório base do aplicativo ou o diretório de execução PyInstaller. O método verifica se o caminho fornecido é absoluto e, se for, retorna\sphinxhyphen{}o diretamente. Caso contrário, ele determina o diretório base usando a variável sys.\_MEIPASS (que é definida pelo PyInstaller durante a execução) ou o diretório do arquivo atual, e então combina esse diretório base com o caminho relativo fornecido para obter o caminho absoluto. Se a entrada for None ou vazia, o método retorna uma string vazia.
:param p: O caminho relativo do arquivo.
:type p: \DUrole{sphinx_autodoc_typehints-type}{\sphinxcode{\sphinxupquote{str}}}
:returns: O caminho absoluto do arquivo, considerando o diretório base do aplicativo ou o diretório de execução PyInstaller.
\begin{quote}\begin{description}
\sphinxlineitem{Return type}
\sphinxAtStartPar
\DUrole{sphinx_autodoc_typehints-type}{\sphinxcode{\sphinxupquote{str}}}

\end{description}\end{quote}

\end{fulllineitems}

\index{read\_license\_text() (in module license\_dialog)@\spxentry{read\_license\_text()}\spxextra{in module license\_dialog}}

\begin{fulllineitems}
\phantomsection\label{\detokenize{api/license_dialog:license_dialog.read_license_text}}
\pysigstartsignatures
\pysiglinewithargsret
{\sphinxcode{\sphinxupquote{license\_dialog.}}\sphinxbfcode{\sphinxupquote{read\_license\_text}}}
{\sphinxparam{\DUrole{n}{preferred\_names}\DUrole{o}{=}\DUrole{default_value}{None}}\sphinxparamcomma \sphinxparam{\DUrole{n}{fallback\_message}\DUrole{o}{=}\DUrole{default_value}{\textquotesingle{}Licença não encontrada.\textbackslash{}\textbackslash{}n\textbackslash{}\textbackslash{}nInclua um arquivo LICENSE.pt\sphinxhyphen{}BR.txt (ou LICENSE.txt / LICENSE) junto ao instalador.\textquotesingle{}}}}
{}
\pysigstopsignatures\begin{quote}\begin{description}
\sphinxlineitem{Return type}
\sphinxAtStartPar

\sphinxAtStartPar
\DUrole{sphinx_autodoc_typehints-type}{\sphinxcode{\sphinxupquote{str}}}

\sphinxAtStartPar
Lê o conteúdo de um arquivo de licença a partir de uma lista de nomes de arquivos preferenciais, retornando o conteúdo do primeiro arquivo encontrado. O método itera sobre a lista de nomes de arquivos fornecida (ou uma lista padrão se None for fornecida) e tenta ler o conteúdo de cada arquivo usando o caminho absoluto resolvido. Se um arquivo for encontrado e lido com sucesso, seu conteúdo é retornado. Se nenhum dos arquivos na lista for encontrado ou se ocorrer um erro ao ler os arquivos, o método retorna a mensagem de fallback especificada.
Args:
\begin{quote}

\sphinxAtStartPar
preferred\_names (Optional{[}list{[}str{]}{]}, optional): Lista de nomes de arquivos de licença a serem verificados em ordem de preferência. O método tentará ler o conteúdo do arquivo de licença usando cada nome na lista até encontrar um arquivo existente e legível. Se nenhum dos arquivos na lista for encontrado ou se ocorrer um erro ao ler os arquivos, o método retornará a mensagem de fallback. Se a lista for None ou vazia, o método usará uma lista padrão de nomes de arquivos de licença para verificar.
fallback\_message (str, optional):  Defaults to ( “Licença não encontrada.
\end{quote}


\end{description}\end{quote}

\sphinxAtStartPar
“ “Inclua um arquivo LICENSE.pt\sphinxhyphen{}BR.txt (ou LICENSE.txt / LICENSE) junto ao instalador.” ).
\begin{quote}
\begin{description}
\sphinxlineitem{Returns:}
\sphinxAtStartPar
str: O conteúdo do arquivo de licença encontrado, ou a mensagem de fallback se nenhum arquivo for encontrado.

\end{description}
\end{quote}

\end{fulllineitems}

\index{LicenseDialog (class in license\_dialog)@\spxentry{LicenseDialog}\spxextra{class in license\_dialog}}

\begin{fulllineitems}
\phantomsection\label{\detokenize{api/license_dialog:license_dialog.LicenseDialog}}
\pysigstartsignatures
\pysiglinewithargsret
{\sphinxbfcode{\sphinxupquote{\DUrole{k}{class}\DUrole{w}{ }}}\sphinxcode{\sphinxupquote{license\_dialog.}}\sphinxbfcode{\sphinxupquote{LicenseDialog}}}
{\sphinxparam{\DUrole{n}{parent}\DUrole{o}{=}\DUrole{default_value}{None}}\sphinxparamcomma \sphinxparam{\DUrole{n}{theme}\DUrole{o}{=}\DUrole{default_value}{None}}\sphinxparamcomma \sphinxparam{\DUrole{n}{license\_text}\DUrole{o}{=}\DUrole{default_value}{\textquotesingle{}\textquotesingle{}}}}
{}
\pysigstopsignatures
\sphinxAtStartPar
Bases: \sphinxcode{\sphinxupquote{QDialog}}
\begin{quote}

\sphinxAtStartPar
Diálogo para exibir o texto da licença do aplicativo, com suporte para personalização de tema e leitura de arquivos de licença a partir de uma lista de nomes preferenciais. O diálogo é modal e inclui um QTextEdit para exibir o conteúdo da licença, bem como um QDialogButtonBox com um botão de fechar. O método apply\_theme permite personalizar as cores do diálogo com base em um dicionário de tema fornecido.
\end{quote}
\begin{quote}\begin{description}
\sphinxlineitem{Parameters}
\sphinxAtStartPar
\sphinxstyleliteralstrong{\sphinxupquote{QDialog}} (\sphinxstyleliteralemphasis{\sphinxupquote{\_type\_}}) \textendash{} Diálogo para exibir o texto da licença do aplicativo, com suporte para personalização de tema e leitura de arquivos de licença a partir de uma lista de nomes preferenciais. O diálogo é modal e inclui um QTextEdit para exibir o conteúdo da licença, bem como um QDialogButtonBox com um botão de fechar. O método apply\_theme permite personalizar as cores do diálogo com base em um dicionário de tema fornecido.

\end{description}\end{quote}
\index{apply\_theme() (license\_dialog.LicenseDialog method)@\spxentry{apply\_theme()}\spxextra{license\_dialog.LicenseDialog method}}

\begin{fulllineitems}
\phantomsection\label{\detokenize{api/license_dialog:license_dialog.LicenseDialog.apply_theme}}
\pysigstartsignatures
\pysiglinewithargsret
{\sphinxbfcode{\sphinxupquote{apply\_theme}}}
{\sphinxparam{\DUrole{n}{theme}}}
{}
\pysigstopsignatures
\sphinxAtStartPar
Aplica um tema personalizado ao diálogo de licença, definindo as cores de fundo, superfície, texto e borda com base nas especificações do dicionário de tema fornecido. O método extrai as cores do dicionário de tema usando as chaves “background”, “surface”, “text” e “border”, e então constrói uma string de estilo CSS para aplicar essas cores aos elementos do diálogo, como o QDialog, QTextEdit e QPushButton. O estilo é aplicado ao diálogo usando o método setStyleSheet.
:param theme: Dicionário de tema para personalizar as cores do diálogo. O dicionário deve conter as chaves “background”, “surface”, “text” e “border” com valores de cor hexadecimal para personalizar a aparência do diálogo. O método usará essas cores para construir um estilo CSS que será aplicado ao diálogo, definindo as cores de fundo, superfície, texto e borda dos elementos do diálogo.
:type theme: \DUrole{sphinx_autodoc_typehints-type}{\sphinxcode{\sphinxupquote{dict}}}
\begin{quote}\begin{description}
\sphinxlineitem{Return type}
\sphinxAtStartPar
\DUrole{sphinx_autodoc_typehints-type}{\sphinxcode{\sphinxupquote{None}}}

\end{description}\end{quote}

\end{fulllineitems}


\end{fulllineitems}


\sphinxstepscope


\section{loading\_screen module}
\label{\detokenize{api/loading_screen:module-loading_screen}}\label{\detokenize{api/loading_screen:loading-screen-module}}\label{\detokenize{api/loading_screen::doc}}\index{module@\spxentry{module}!loading\_screen@\spxentry{loading\_screen}}\index{loading\_screen@\spxentry{loading\_screen}!module@\spxentry{module}}\index{LoadingScreen (class in loading\_screen)@\spxentry{LoadingScreen}\spxextra{class in loading\_screen}}

\begin{fulllineitems}
\phantomsection\label{\detokenize{api/loading_screen:loading_screen.LoadingScreen}}
\pysigstartsignatures
\pysiglinewithargsret
{\sphinxbfcode{\sphinxupquote{\DUrole{k}{class}\DUrole{w}{ }}}\sphinxcode{\sphinxupquote{loading\_screen.}}\sphinxbfcode{\sphinxupquote{LoadingScreen}}}
{\sphinxparam{\DUrole{n}{app}}\sphinxparamcomma \sphinxparam{\DUrole{n}{title}\DUrole{o}{=}\DUrole{default_value}{\textquotesingle{}Iniciando\textquotesingle{}}}\sphinxparamcomma \sphinxparam{\DUrole{n}{subtitle}\DUrole{o}{=}\DUrole{default_value}{\textquotesingle{}Aguarde enquanto validamos o acesso...\textquotesingle{}}}\sphinxparamcomma \sphinxparam{\DUrole{n}{width}\DUrole{o}{=}\DUrole{default_value}{520}}\sphinxparamcomma \sphinxparam{\DUrole{n}{height}\DUrole{o}{=}\DUrole{default_value}{220}}\sphinxparamcomma \sphinxparam{\DUrole{n}{bg}\DUrole{o}{=}\DUrole{default_value}{\textquotesingle{}\#0B1220\textquotesingle{}}}\sphinxparamcomma \sphinxparam{\DUrole{n}{accent}\DUrole{o}{=}\DUrole{default_value}{\textquotesingle{}\#3B82F6\textquotesingle{}}}\sphinxparamcomma \sphinxparam{\DUrole{n}{text}\DUrole{o}{=}\DUrole{default_value}{\textquotesingle{}\#E6EDF7\textquotesingle{}}}\sphinxparamcomma \sphinxparam{\DUrole{n}{muted}\DUrole{o}{=}\DUrole{default_value}{\textquotesingle{}\#A7B3C6\textquotesingle{}}}}
{}
\pysigstopsignatures
\sphinxAtStartPar
Bases: \sphinxcode{\sphinxupquote{object}}

\sphinxAtStartPar
Classe para exibir uma tela de carregamento personalizada durante a inicialização do aplicativo, com suporte para mensagens dinâmicas e personalização de cores. O método construtor recebe parâmetros para configurar o título, subtítulo, dimensões e cores da tela de carregamento, e cria um QPixmap personalizado usando o método \_make\_pixmap. A classe inclui métodos para mostrar a tela de carregamento, atualizar a mensagem exibida, ocultar a tela e finalizar o splash de forma confiável em relação à janela principal do aplicativo.
\index{show() (loading\_screen.LoadingScreen method)@\spxentry{show()}\spxextra{loading\_screen.LoadingScreen method}}

\begin{fulllineitems}
\phantomsection\label{\detokenize{api/loading_screen:loading_screen.LoadingScreen.show}}
\pysigstartsignatures
\pysiglinewithargsret
{\sphinxbfcode{\sphinxupquote{show}}}
{\sphinxparam{\DUrole{n}{message}\DUrole{o}{=}\DUrole{default_value}{\textquotesingle{}Preparando...\textquotesingle{}}}}
{}
\pysigstopsignatures
\sphinxAtStartPar
Exibe a tela de carregamento e atualiza a mensagem exibida. O método mostra o splash usando o método show do QSplashScreen, e então chama o método update para definir a mensagem exibida na tela de carregamento. A mensagem é alinhada à parte inferior esquerda da tela usando as flags Qt.AlignBottom | Qt.AlignLeft, e a cor do texto é definida como a cor de texto secundário configurada na classe. O método processa os eventos do aplicativo para garantir que a interface seja atualizada corretamente.
:param message: A mensagem a ser exibida na tela de carregamento. Defaults to “Preparando…”.
:type message: \DUrole{sphinx_autodoc_typehints-type}{\sphinxcode{\sphinxupquote{str}}}
\begin{quote}\begin{description}
\sphinxlineitem{Return type}
\sphinxAtStartPar
\DUrole{sphinx_autodoc_typehints-type}{\sphinxcode{\sphinxupquote{None}}}

\end{description}\end{quote}

\end{fulllineitems}

\index{update() (loading\_screen.LoadingScreen method)@\spxentry{update()}\spxextra{loading\_screen.LoadingScreen method}}

\begin{fulllineitems}
\phantomsection\label{\detokenize{api/loading_screen:loading_screen.LoadingScreen.update}}
\pysigstartsignatures
\pysiglinewithargsret
{\sphinxbfcode{\sphinxupquote{update}}}
{\sphinxparam{\DUrole{n}{message}}}
{}
\pysigstopsignatures
\sphinxAtStartPar
Atualiza a mensagem exibida na tela de carregamento. O método utiliza o método showMessage do QSplashScreen para definir a nova mensagem, alinhando\sphinxhyphen{}a à parte inferior esquerda da tela e usando a cor de texto secundário configurada na classe. Após atualizar a mensagem, o método processa os eventos do aplicativo para garantir que a interface seja atualizada e a nova mensagem seja exibida corretamente.
:param message: A nova mensagem a ser exibida na tela de carregamento.
:type message: \DUrole{sphinx_autodoc_typehints-type}{\sphinxcode{\sphinxupquote{str}}}
\begin{quote}\begin{description}
\sphinxlineitem{Return type}
\sphinxAtStartPar
\DUrole{sphinx_autodoc_typehints-type}{\sphinxcode{\sphinxupquote{None}}}

\end{description}\end{quote}

\end{fulllineitems}

\index{hide() (loading\_screen.LoadingScreen method)@\spxentry{hide()}\spxextra{loading\_screen.LoadingScreen method}}

\begin{fulllineitems}
\phantomsection\label{\detokenize{api/loading_screen:loading_screen.LoadingScreen.hide}}
\pysigstartsignatures
\pysiglinewithargsret
{\sphinxbfcode{\sphinxupquote{hide}}}
{}
{}
\pysigstopsignatures
\sphinxAtStartPar
Oculta a tela de carregamento. O método chama o método hide do QSplashScreen para ocultar a tela de carregamento, e então processa os eventos do aplicativo para garantir que a interface seja atualizada corretamente após ocultar o splash.
\begin{quote}\begin{description}
\sphinxlineitem{Return type}
\sphinxAtStartPar
\DUrole{sphinx_autodoc_typehints-type}{\sphinxcode{\sphinxupquote{None}}}

\end{description}\end{quote}

\end{fulllineitems}

\index{close() (loading\_screen.LoadingScreen method)@\spxentry{close()}\spxextra{loading\_screen.LoadingScreen method}}

\begin{fulllineitems}
\phantomsection\label{\detokenize{api/loading_screen:loading_screen.LoadingScreen.close}}
\pysigstartsignatures
\pysiglinewithargsret
{\sphinxbfcode{\sphinxupquote{close}}}
{}
{}
\pysigstopsignatures
\sphinxAtStartPar
Fecha a tela de carregamento de forma confiável. O método chama o método close do QSplashScreen para fechar a tela de carregamento, e processa os eventos do aplicativo antes e depois de fechar o splash para garantir que a interface seja atualizada corretamente e que o fechamento ocorra de forma suave.
\begin{quote}\begin{description}
\sphinxlineitem{Return type}
\sphinxAtStartPar
\DUrole{sphinx_autodoc_typehints-type}{\sphinxcode{\sphinxupquote{None}}}

\end{description}\end{quote}

\end{fulllineitems}

\index{finish\_and\_close() (loading\_screen.LoadingScreen method)@\spxentry{finish\_and\_close()}\spxextra{loading\_screen.LoadingScreen method}}

\begin{fulllineitems}
\phantomsection\label{\detokenize{api/loading_screen:loading_screen.LoadingScreen.finish_and_close}}
\pysigstartsignatures
\pysiglinewithargsret
{\sphinxbfcode{\sphinxupquote{finish\_and\_close}}}
{\sphinxparam{\DUrole{n}{main\_window}}}
{}
\pysigstopsignatures
\sphinxAtStartPar
Fecha o splash de forma confiável.
\sphinxhyphen{} finish() tenta sincronizar com a janela principal
\sphinxhyphen{} close() garante o fechamento mesmo se a janela ainda não foi “exposta”
\begin{quote}\begin{description}
\sphinxlineitem{Return type}
\sphinxAtStartPar
\DUrole{sphinx_autodoc_typehints-type}{\sphinxcode{\sphinxupquote{None}}}

\end{description}\end{quote}

\end{fulllineitems}


\end{fulllineitems}


\sphinxstepscope


\section{models module}
\label{\detokenize{api/models:module-models}}\label{\detokenize{api/models:models-module}}\label{\detokenize{api/models::doc}}\index{module@\spxentry{module}!models@\spxentry{models}}\index{models@\spxentry{models}!module@\spxentry{module}}\index{Campo (class in models)@\spxentry{Campo}\spxextra{class in models}}

\begin{fulllineitems}
\phantomsection\label{\detokenize{api/models:models.Campo}}
\pysigstartsignatures
\pysiglinewithargsret
{\sphinxbfcode{\sphinxupquote{\DUrole{k}{class}\DUrole{w}{ }}}\sphinxcode{\sphinxupquote{models.}}\sphinxbfcode{\sphinxupquote{Campo}}}
{\sphinxparam{\DUrole{n}{id}}\sphinxparamcomma \sphinxparam{\DUrole{n}{titulo}}\sphinxparamcomma \sphinxparam{\DUrole{n}{tipo}\DUrole{o}{=}\DUrole{default_value}{\textquotesingle{}texto\textquotesingle{}}}\sphinxparamcomma \sphinxparam{\DUrole{n}{fixo}\DUrole{o}{=}\DUrole{default_value}{False}}\sphinxparamcomma \sphinxparam{\DUrole{n}{locked}\DUrole{o}{=}\DUrole{default_value}{False}}}
{}
\pysigstopsignatures
\sphinxAtStartPar
Bases: \sphinxcode{\sphinxupquote{object}}

\sphinxAtStartPar
Representa um campo de extração, com propriedades como id, título, tipo, e flags para indicar se o campo é fixo ou bloqueado. O campo é identificado por um id único gerado a partir do título, e possui um tipo que pode ser inferido a partir do título ou definido explicitamente. A flag “fixo” indica se o campo é obrigatório e não pode ser removido, enquanto a flag “locked” indica se o campo está bloqueado para edição. A classe Campo é usada para definir os campos que serão extraídos dos documentos, e suas propriedades são utilizadas para configurar a extração e a exportação dos dados.
\index{id (models.Campo attribute)@\spxentry{id}\spxextra{models.Campo attribute}}

\begin{fulllineitems}
\phantomsection\label{\detokenize{api/models:models.Campo.id}}
\pysigstartsignatures
\pysigline
{\sphinxbfcode{\sphinxupquote{id}}\sphinxbfcode{\sphinxupquote{\DUrole{p}{:}\DUrole{w}{ }str}}}
\pysigstopsignatures
\end{fulllineitems}

\index{titulo (models.Campo attribute)@\spxentry{titulo}\spxextra{models.Campo attribute}}

\begin{fulllineitems}
\phantomsection\label{\detokenize{api/models:models.Campo.titulo}}
\pysigstartsignatures
\pysigline
{\sphinxbfcode{\sphinxupquote{titulo}}\sphinxbfcode{\sphinxupquote{\DUrole{p}{:}\DUrole{w}{ }str}}}
\pysigstopsignatures
\end{fulllineitems}

\index{tipo (models.Campo attribute)@\spxentry{tipo}\spxextra{models.Campo attribute}}

\begin{fulllineitems}
\phantomsection\label{\detokenize{api/models:models.Campo.tipo}}
\pysigstartsignatures
\pysigline
{\sphinxbfcode{\sphinxupquote{tipo}}\sphinxbfcode{\sphinxupquote{\DUrole{p}{:}\DUrole{w}{ }str}}\sphinxbfcode{\sphinxupquote{\DUrole{w}{ }\DUrole{p}{=}\DUrole{w}{ }\textquotesingle{}texto\textquotesingle{}}}}
\pysigstopsignatures
\end{fulllineitems}

\index{fixo (models.Campo attribute)@\spxentry{fixo}\spxextra{models.Campo attribute}}

\begin{fulllineitems}
\phantomsection\label{\detokenize{api/models:models.Campo.fixo}}
\pysigstartsignatures
\pysigline
{\sphinxbfcode{\sphinxupquote{fixo}}\sphinxbfcode{\sphinxupquote{\DUrole{p}{:}\DUrole{w}{ }bool}}\sphinxbfcode{\sphinxupquote{\DUrole{w}{ }\DUrole{p}{=}\DUrole{w}{ }False}}}
\pysigstopsignatures
\end{fulllineitems}

\index{locked (models.Campo attribute)@\spxentry{locked}\spxextra{models.Campo attribute}}

\begin{fulllineitems}
\phantomsection\label{\detokenize{api/models:models.Campo.locked}}
\pysigstartsignatures
\pysigline
{\sphinxbfcode{\sphinxupquote{locked}}\sphinxbfcode{\sphinxupquote{\DUrole{p}{:}\DUrole{w}{ }bool}}\sphinxbfcode{\sphinxupquote{\DUrole{w}{ }\DUrole{p}{=}\DUrole{w}{ }False}}}
\pysigstopsignatures
\end{fulllineitems}


\end{fulllineitems}

\index{sanitize\_id() (in module models)@\spxentry{sanitize\_id()}\spxextra{in module models}}

\begin{fulllineitems}
\phantomsection\label{\detokenize{api/models:models.sanitize_id}}
\pysigstartsignatures
\pysiglinewithargsret
{\sphinxcode{\sphinxupquote{models.}}\sphinxbfcode{\sphinxupquote{sanitize\_id}}}
{\sphinxparam{\DUrole{n}{title}}}
{}
\pysigstopsignatures
\sphinxAtStartPar
Gera um ID sanitizado a partir de um título, convertendo\sphinxhyphen{}o para minúsculas, substituindo espaços por underscores, e removendo caracteres especiais. O método processa o título de entrada para criar um ID que seja seguro para uso em código e arquivos, garantindo que ele contenha apenas caracteres alfanuméricos e underscores. O resultado é um ID único e legível que pode ser usado para identificar campos de extração de forma consistente.
:param title: O título a partir do qual o ID será gerado. O método irá processar essa string para criar um ID sanitizado, convertendo\sphinxhyphen{}a para minúsculas, substituindo espaços por underscores, e removendo caracteres especiais, resultando em um ID seguro para uso em código e arquivos.
:type title: \DUrole{sphinx_autodoc_typehints-type}{\sphinxcode{\sphinxupquote{str}}}
:returns: O ID sanitizado gerado a partir do título, contendo apenas caracteres alfanuméricos e underscores, e convertido para minúsculas. O método garante que o ID seja legível e consistente para identificar campos de extração.
\begin{quote}\begin{description}
\sphinxlineitem{Return type}
\sphinxAtStartPar
\DUrole{sphinx_autodoc_typehints-type}{\sphinxcode{\sphinxupquote{str}}}

\end{description}\end{quote}

\end{fulllineitems}

\index{make\_unique\_id() (in module models)@\spxentry{make\_unique\_id()}\spxextra{in module models}}

\begin{fulllineitems}
\phantomsection\label{\detokenize{api/models:models.make_unique_id}}
\pysigstartsignatures
\pysiglinewithargsret
{\sphinxcode{\sphinxupquote{models.}}\sphinxbfcode{\sphinxupquote{make\_unique\_id}}}
{\sphinxparam{\DUrole{n}{base\_id}}\sphinxparamcomma \sphinxparam{\DUrole{n}{used}}}
{}
\pysigstopsignatures
\sphinxAtStartPar
Gera um ID único a partir de um ID base, garantindo que ele não esteja presente em um conjunto de IDs já utilizados. O método verifica se o ID base já está no conjunto de IDs utilizados, e se estiver, ele adiciona um sufixo numérico incremental (começando com “\_2”) até encontrar um ID que seja único. O ID gerado é adicionado ao conjunto de IDs utilizados para garantir que futuras chamadas ao método também considerem esse ID como utilizado. O resultado é um ID único que pode ser usado para identificar campos de extração ou outros elementos de forma consistente e sem conflitos.
:param base\_id: O ID base a ser usado como base para gerar um ID único.
:type base\_id: \DUrole{sphinx_autodoc_typehints-type}{\sphinxcode{\sphinxupquote{str}}}
:param used: Um conjunto de IDs já utilizados, para garantir unicidade.
:type used: \DUrole{sphinx_autodoc_typehints-type}{\sphinxcode{\sphinxupquote{set}}}
:returns: Um ID único gerado a partir do base\_id e do conjunto de IDs utilizados.
\begin{quote}\begin{description}
\sphinxlineitem{Return type}
\sphinxAtStartPar
\DUrole{sphinx_autodoc_typehints-type}{\sphinxcode{\sphinxupquote{str}}}

\end{description}\end{quote}

\end{fulllineitems}

\index{infer\_type\_from\_title() (in module models)@\spxentry{infer\_type\_from\_title()}\spxextra{in module models}}

\begin{fulllineitems}
\phantomsection\label{\detokenize{api/models:models.infer_type_from_title}}
\pysigstartsignatures
\pysiglinewithargsret
{\sphinxcode{\sphinxupquote{models.}}\sphinxbfcode{\sphinxupquote{infer\_type\_from\_title}}}
{\sphinxparam{\DUrole{n}{header}}}
{}
\pysigstopsignatures
\sphinxAtStartPar
Infere o tipo de um campo a partir do título do campo, utilizando palavras\sphinxhyphen{}chave comuns para identificar tipos como “email”, “telefone”, “data”, “moeda”, “numero” e “booleano”. O método processa o título de entrada, convertendo\sphinxhyphen{}o para minúsculas e verificando a presença de palavras\sphinxhyphen{}chave específicas que indicam cada tipo. Se uma palavra\sphinxhyphen{}chave correspondente for encontrada, o método retorna o tipo inferido. Se nenhuma palavra\sphinxhyphen{}chave for encontrada, ele retorna “texto” como tipo padrão.
:param header: O título do campo a partir do qual o tipo será inferido.
:type header: \DUrole{sphinx_autodoc_typehints-type}{\sphinxcode{\sphinxupquote{str}}}
:returns: O tipo inferido a partir do título do campo.
\begin{quote}\begin{description}
\sphinxlineitem{Return type}
\sphinxAtStartPar
\DUrole{sphinx_autodoc_typehints-type}{\sphinxcode{\sphinxupquote{str}}}

\end{description}\end{quote}

\end{fulllineitems}


\sphinxstepscope


\section{online\_license module}
\label{\detokenize{api/online_license:module-online_license}}\label{\detokenize{api/online_license:online-license-module}}\label{\detokenize{api/online_license::doc}}\index{module@\spxentry{module}!online\_license@\spxentry{online\_license}}\index{online\_license@\spxentry{online\_license}!module@\spxentry{module}}\index{app\_state\_dir() (in module online\_license)@\spxentry{app\_state\_dir()}\spxextra{in module online\_license}}

\begin{fulllineitems}
\phantomsection\label{\detokenize{api/online_license:online_license.app_state_dir}}
\pysigstartsignatures
\pysiglinewithargsret
{\sphinxcode{\sphinxupquote{online\_license.}}\sphinxbfcode{\sphinxupquote{app\_state\_dir}}}
{}
{}
\pysigstopsignatures
\sphinxAtStartPar
Retorna o caminho do diretório de estado do aplicativo, onde dados como cache de licença podem ser armazenados. O método tenta criar um subdiretório com o nome do aplicativo dentro de vários caminhos base comuns (APPDATA, LOCALAPPDATA, PROGRAMDATA) e verifica se é possível escrever nesse diretório. Se nenhum dos caminhos base permitir a criação de um diretório gravável, o método cria um diretório com o nome do aplicativo dentro do diretório home do usuário. O resultado é um caminho de diretório garantido onde o aplicativo pode armazenar seus dados de estado.
:returns: O caminho do diretório de estado do aplicativo, onde dados como cache de licença podem ser armazenados. O método tenta criar um subdiretório com o nome do aplicativo dentro de vários caminhos base comuns (APPDATA, LOCALAPPDATA, PROGRAMDATA) e verifica se é possível escrever nesse diretório. Se nenhum dos caminhos base permitir a criação de um diretório gravável, o método cria um diretório com o nome do aplicativo dentro do diretório home do usuário. O resultado é um caminho de diretório garantido onde o aplicativo pode armazenar seus dados de estado.
\begin{quote}\begin{description}
\sphinxlineitem{Return type}
\sphinxAtStartPar
\DUrole{sphinx_autodoc_typehints-type}{\sphinxcode{\sphinxupquote{str}}}

\end{description}\end{quote}

\end{fulllineitems}

\index{cache\_path() (in module online\_license)@\spxentry{cache\_path()}\spxextra{in module online\_license}}

\begin{fulllineitems}
\phantomsection\label{\detokenize{api/online_license:online_license.cache_path}}
\pysigstartsignatures
\pysiglinewithargsret
{\sphinxcode{\sphinxupquote{online\_license.}}\sphinxbfcode{\sphinxupquote{cache\_path}}}
{}
{}
\pysigstopsignatures
\sphinxAtStartPar
Retorna o caminho completo do arquivo de cache de licença, que é construído combinando o diretório de estado do aplicativo (obtido pela função app\_state\_dir) com o nome do arquivo de cache definido na constante LICENSE\_CACHE\_FILE. O resultado é um caminho absoluto para o arquivo de cache onde os dados de licença podem ser armazenados e recuperados.
:returns: O caminho completo do arquivo de cache de licença, que é construído combinando o diretório de estado do aplicativo (obtido pela função app\_state\_dir) com o nome do arquivo de cache definido na constante LICENSE\_CACHE\_FILE. O resultado é um caminho absoluto para o arquivo de cache onde os dados de licença podem ser armazenados e recuperados.
\begin{quote}\begin{description}
\sphinxlineitem{Return type}
\sphinxAtStartPar
\DUrole{sphinx_autodoc_typehints-type}{\sphinxcode{\sphinxupquote{str}}}

\end{description}\end{quote}

\end{fulllineitems}

\index{clear\_cache() (in module online\_license)@\spxentry{clear\_cache()}\spxextra{in module online\_license}}

\begin{fulllineitems}
\phantomsection\label{\detokenize{api/online_license:online_license.clear_cache}}
\pysigstartsignatures
\pysiglinewithargsret
{\sphinxcode{\sphinxupquote{online\_license.}}\sphinxbfcode{\sphinxupquote{clear\_cache}}}
{}
{}
\pysigstopsignatures
\sphinxAtStartPar
Limpa o cache de licença removendo o arquivo de cache se ele existir. O método obtém o caminho do arquivo de cache usando a função cache\_path, verifica se o arquivo existe e, se for encontrado, tenta removê\sphinxhyphen{}lo usando a função os.remove. Se ocorrer um erro durante a remoção do arquivo, o método ignora a exceção e continua, garantindo que o cache seja limpo sem causar falhas no aplicativo.
\begin{quote}\begin{description}
\sphinxlineitem{Return type}
\sphinxAtStartPar
\DUrole{sphinx_autodoc_typehints-type}{\sphinxcode{\sphinxupquote{None}}}

\end{description}\end{quote}

\end{fulllineitems}

\index{now\_ts() (in module online\_license)@\spxentry{now\_ts()}\spxextra{in module online\_license}}

\begin{fulllineitems}
\phantomsection\label{\detokenize{api/online_license:online_license.now_ts}}
\pysigstartsignatures
\pysiglinewithargsret
{\sphinxcode{\sphinxupquote{online\_license.}}\sphinxbfcode{\sphinxupquote{now\_ts}}}
{}
{}
\pysigstopsignatures
\sphinxAtStartPar
Retorna o timestamp atual em segundos desde a época (1 de janeiro de 1970). O método utiliza a função time.time para obter o tempo atual em segundos como um número de ponto flutuante, e então converte esse valor para um inteiro antes de retorná\sphinxhyphen{}lo. O resultado é o número de segundos inteiros que se passaram desde a época, representando o momento atual.
:returns: O timestamp atual em segundos desde a época (1 de janeiro de 1970), obtido usando a função time.time e convertido para um inteiro.
\begin{quote}\begin{description}
\sphinxlineitem{Return type}
\sphinxAtStartPar
\DUrole{sphinx_autodoc_typehints-type}{\sphinxcode{\sphinxupquote{int}}}

\end{description}\end{quote}

\end{fulllineitems}

\index{load\_cache() (in module online\_license)@\spxentry{load\_cache()}\spxextra{in module online\_license}}

\begin{fulllineitems}
\phantomsection\label{\detokenize{api/online_license:online_license.load_cache}}
\pysigstartsignatures
\pysiglinewithargsret
{\sphinxcode{\sphinxupquote{online\_license.}}\sphinxbfcode{\sphinxupquote{load\_cache}}}
{}
{}
\pysigstopsignatures
\sphinxAtStartPar
Carrega o cache de licença a partir do arquivo de cache, retornando um dicionário com os dados armazenados. O método tenta abrir o arquivo de cache usando o caminho obtido pela função cache\_path, lê seu conteúdo como JSON e retorna o resultado como um dicionário. Se o arquivo não existir ou se ocorrer um erro durante a leitura ou a conversão do JSON, o método retorna um dicionário vazio, indicando que não há dados de cache disponíveis.
:returns: O conteúdo do cache de licença carregado a partir do arquivo de cache, ou um dicionário vazio se o arquivo não existir ou se ocorrer um erro ao ler o arquivo. O método tenta abrir o arquivo de cache usando o caminho obtido pela função cache\_path, lê seu conteúdo como JSON e retorna o resultado como um dicionário. Se o arquivo não existir ou se ocorrer um erro durante a leitura ou a conversão do JSON, o método retorna um dicionário vazio.
\begin{quote}\begin{description}
\sphinxlineitem{Return type}
\sphinxAtStartPar
\DUrole{sphinx_autodoc_typehints-type}{\sphinxcode{\sphinxupquote{dict}}}

\end{description}\end{quote}

\end{fulllineitems}

\index{save\_cache() (in module online\_license)@\spxentry{save\_cache()}\spxextra{in module online\_license}}

\begin{fulllineitems}
\phantomsection\label{\detokenize{api/online_license:online_license.save_cache}}
\pysigstartsignatures
\pysiglinewithargsret
{\sphinxcode{\sphinxupquote{online\_license.}}\sphinxbfcode{\sphinxupquote{save\_cache}}}
{\sphinxparam{\DUrole{n}{data}}}
{}
\pysigstopsignatures
\sphinxAtStartPar
Salva os dados de licença no cache escrevendo um dicionário como JSON em um arquivo de cache. O método recebe um dicionário de dados, converte\sphinxhyphen{}o para uma string JSON formatada usando a função json.dump, e escreve essa string no arquivo de cache localizado pelo caminho obtido pela função cache\_path. Se ocorrer um erro durante a escrita do arquivo, o método irá lançar uma exceção, indicando que o processo de salvamento do cache falhou.
:param data: O dicionário de dados a ser salvo no cache de licença. O método irá converter esse dicionário para JSON e escrevê\sphinxhyphen{}lo no arquivo de cache usando o caminho obtido pela função cache\_path. Se ocorrer um erro durante a escrita do arquivo, o método irá lançar uma exceção.
:type data: \DUrole{sphinx_autodoc_typehints-type}{\sphinxcode{\sphinxupquote{dict}}}
\begin{quote}\begin{description}
\sphinxlineitem{Return type}
\sphinxAtStartPar
\DUrole{sphinx_autodoc_typehints-type}{\sphinxcode{\sphinxupquote{None}}}

\end{description}\end{quote}

\end{fulllineitems}

\index{get\_machine\_guid() (in module online\_license)@\spxentry{get\_machine\_guid()}\spxextra{in module online\_license}}

\begin{fulllineitems}
\phantomsection\label{\detokenize{api/online_license:online_license.get_machine_guid}}
\pysigstartsignatures
\pysiglinewithargsret
{\sphinxcode{\sphinxupquote{online\_license.}}\sphinxbfcode{\sphinxupquote{get\_machine\_guid}}}
{}
{}
\pysigstopsignatures
\sphinxAtStartPar
Obtém o valor do MachineGuid do sistema a partir do registro do Windows. O método abre a chave de registro HKEY\_LOCAL\_MACHINESOFTWAREMicrosoftCryptography, consulta o valor da chave “MachineGuid” e retorna esse valor como uma string. O MachineGuid é um identificador único para o computador, e é comumente usado para fins de licenciamento e identificação de dispositivos. Se ocorrer um erro ao acessar o registro ou ao obter o valor, o método irá lançar uma exceção, indicando que a obtenção do MachineGuid falhou.
:returns: O valor do MachineGuid do sistema, obtido a partir do registro do Windows. O método abre a chave de registro HKEY\_LOCAL\_MACHINESOFTWAREMicrosoftCryptography, consulta o valor da chave “MachineGuid” e retorna esse valor como uma string. Se ocorrer um erro ao acessar o registro ou ao obter o valor, o método irá lançar uma exceção.
\begin{quote}\begin{description}
\sphinxlineitem{Return type}
\sphinxAtStartPar
\DUrole{sphinx_autodoc_typehints-type}{\sphinxcode{\sphinxupquote{str}}}

\end{description}\end{quote}

\end{fulllineitems}

\index{device\_id() (in module online\_license)@\spxentry{device\_id()}\spxextra{in module online\_license}}

\begin{fulllineitems}
\phantomsection\label{\detokenize{api/online_license:online_license.device_id}}
\pysigstartsignatures
\pysiglinewithargsret
{\sphinxcode{\sphinxupquote{online\_license.}}\sphinxbfcode{\sphinxupquote{device\_id}}}
{}
{}
\pysigstopsignatures
\sphinxAtStartPar
Gera um identificador de dispositivo único a partir do valor do MachineGuid do sistema, utilizando o hash SHA\sphinxhyphen{}256 para criar um identificador seguro e consistente. O método obtém o MachineGuid usando a função get\_machine\_guid, converte esse valor para bytes usando a codificação UTF\sphinxhyphen{}8, e então calcula o hash SHA\sphinxhyphen{}256 desses bytes usando a função hashlib.sha256. O resultado é um hash hexadecimal em letras maiúsculas que representa o identificador único do dispositivo, que pode ser usado para fins de licenciamento e identificação de hardware.
:returns: O identificador do dispositivo, gerado a partir do valor do MachineGuid do sistema. O método obtém o MachineGuid usando a função get\_machine\_guid, converte esse valor para bytes usando a codificação UTF\sphinxhyphen{}8, e então calcula o hash SHA\sphinxhyphen{}256 desses bytes usando a função hashlib.sha256. O resultado é um hash hexadecimal em letras maiúsculas que representa o identificador único do dispositivo, que pode ser usado para fins de licenciamento e identificação de hardware.
\begin{quote}\begin{description}
\sphinxlineitem{Return type}
\sphinxAtStartPar
\DUrole{sphinx_autodoc_typehints-type}{\sphinxcode{\sphinxupquote{str}}}

\end{description}\end{quote}

\end{fulllineitems}

\index{call\_api() (in module online\_license)@\spxentry{call\_api()}\spxextra{in module online\_license}}

\begin{fulllineitems}
\phantomsection\label{\detokenize{api/online_license:online_license.call_api}}
\pysigstartsignatures
\pysiglinewithargsret
{\sphinxcode{\sphinxupquote{online\_license.}}\sphinxbfcode{\sphinxupquote{call\_api}}}
{\sphinxparam{\DUrole{n}{action}}\sphinxparamcomma \sphinxparam{\DUrole{n}{email}}\sphinxparamcomma \sphinxparam{\DUrole{n}{dev\_id}}}
{}
\pysigstopsignatures
\sphinxAtStartPar
Faz uma chamada POST para a API de licenciamento, enviando um payload JSON com a ação, o e\sphinxhyphen{}mail do usuário e o identificador do dispositivo. O método constrói um dicionário de payload contendo os parâmetros fornecidos, e então utiliza a função requests.post para enviar esse payload como JSON para a URL da API definida na constante API\_URL. A chamada tem um timeout de 15 segundos para evitar bloqueios prolongados. Após receber a resposta, o método verifica o status HTTP e o tipo de conteúdo da resposta para garantir que seja uma resposta JSON válida. Se a resposta indicar sucesso (status code 200 e Content\sphinxhyphen{}Type contendo “application/json”), o método retorna o conteúdo da resposta como um dicionário. Caso contrário, ele lança uma exceção RuntimeError com uma mensagem apropriada indicando o erro ocorrido.
:param action: A ação a ser realizada na API, como “activate” ou “renew”, que indica o tipo de operação de licenciamento
:type action: \DUrole{sphinx_autodoc_typehints-type}{\sphinxcode{\sphinxupquote{str}}}
:param email: O e\sphinxhyphen{}mail do usuário para o qual a licença está sendo verificada ou ativada. O método inclui esse e\sphinxhyphen{}mail no payload enviado para a API, e é usado para identificar o usuário associado à licença.
:type email: \DUrole{sphinx_autodoc_typehints-type}{\sphinxcode{\sphinxupquote{str}}}
:param dev\_id: O identificador do dispositivo, gerado a partir do MachineGuid do sistema, que é enviado no payload para a API de licenciamento. Esse identificador é usado para associar a licença a um dispositivo específico e garantir que as verificações de licença sejam feitas corretamente com base no hardware do usuário.
:type dev\_id: \DUrole{sphinx_autodoc_typehints-type}{\sphinxcode{\sphinxupquote{str}}}
:raises RuntimeError: Se a resposta da API indicar um erro, como um status HTTP diferente de 200 ou um Content\sphinxhyphen{}Type que não contenha “application/json”, o método lança uma exceção RuntimeError com uma mensagem que descreve o erro ocorrido, incluindo o status code e o conteúdo da resposta para facilitar a depuração. Se ocorrer um erro durante a chamada HTTP, como um timeout ou uma falha de conexão, a exceção correspondente será propagada para o chamador, indicando que a comunicação com a API falhou.
:raises RuntimeError: Se a resposta da API indicar um erro, como um status HTTP diferente de 200 ou um Content\sphinxhyphen{}Type que não contenha “application/json”, o método lança uma exceção RuntimeError com uma mensagem que descreve o erro ocorrido, incluindo o status code e o conteúdo da resposta para facilitar a depuração. Se ocorrer um erro durante a chamada HTTP, como um timeout ou uma falha de conexão, a exceção correspondente será propagada para o chamador, indicando que a comunicação com a API falhou.
\begin{quote}\begin{description}
\sphinxlineitem{Returns}
\sphinxAtStartPar
O conteúdo da resposta da API de licenciamento, retornado como um dicionário. O método faz uma chamada POST para a API, enviando um payload JSON com a ação, o e\sphinxhyphen{}mail do usuário e o identificador do dispositivo. Após receber a resposta, o método verifica o status HTTP e o tipo de conteúdo para garantir que seja uma resposta JSON válida. Se a resposta indicar sucesso (status code 200 e Content\sphinxhyphen{}Type contendo “application/json”), o método retorna o conteúdo da resposta como um dicionário. Caso contrário, ele lança uma exceção RuntimeError com uma mensagem apropriada indicando o erro ocorrido.

\sphinxlineitem{Return type}
\sphinxAtStartPar
\DUrole{sphinx_autodoc_typehints-type}{\sphinxcode{\sphinxupquote{dict}}}

\end{description}\end{quote}

\end{fulllineitems}

\index{ensure\_online\_license() (in module online\_license)@\spxentry{ensure\_online\_license()}\spxextra{in module online\_license}}

\begin{fulllineitems}
\phantomsection\label{\detokenize{api/online_license:online_license.ensure_online_license}}
\pysigstartsignatures
\pysiglinewithargsret
{\sphinxcode{\sphinxupquote{online\_license.}}\sphinxbfcode{\sphinxupquote{ensure\_online\_license}}}
{\sphinxparam{\DUrole{n}{email}}}
{}
\pysigstopsignatures
\sphinxAtStartPar
Verifica a validade da licença online para um usuário específico, utilizando o e\sphinxhyphen{}mail do usuário para consultar a API de licenciamento e retornando uma tupla com o resultado da verificação. O método processa o e\sphinxhyphen{}mail de entrada, obtém o identificador do dispositivo usando a função device\_id, e então verifica se há um cache de licença válido que corresponda ao e\sphinxhyphen{}mail e ao dispositivo. Se um cache válido for encontrado, ele retorna um resultado positivo indicando que a licença é válida (True) junto com uma mensagem de “OK (cache)” e um código de status “ok\_cache”. Se não houver um cache válido, o método determina a ação apropriada (“activate” ou “renew”) com base no estado do cache, e faz uma chamada para a API usando a função call\_api. Dependendo da resposta da API, o método salva os dados no cache se a licença for válida, ou retorna mensagens de erro específicas para diferentes condições de falha (como “no\_license”, “blocked”, “device\_limit”, etc.). O resultado final é uma tupla que indica se a licença é válida, uma mensagem descritiva do resultado, e um código de status que pode ser usado para identificar o resultado específico da verificação.
\begin{quote}\begin{description}
\sphinxlineitem{Parameters}
\sphinxAtStartPar
\sphinxstyleliteralstrong{\sphinxupquote{email}} (\DUrole{sphinx_autodoc_typehints-type}{\sphinxcode{\sphinxupquote{str}}}) \textendash{} O e\sphinxhyphen{}mail do usuário para o qual a licença está sendo verificada ou ativada. O método inclui esse e\sphinxhyphen{}mail no payload enviado para a API, e é usado para identificar o usuário associado à licença.

\sphinxlineitem{Returns}
\sphinxAtStartPar
Uma tupla contendo três elementos: um booleano indicando se a licença é válida (True) ou não (False), uma string com uma mensagem descritiva do resultado da verificação da licença, e uma string com um código de status que pode ser usado para identificar o resultado específico da verificação (como “ok”, “no\_license”, “blocked”, etc.). O método realiza a verificação da licença online, utilizando o e\sphinxhyphen{}mail fornecido para consultar a API de licenciamento, e retorna os resultados dessa verificação na forma de uma tupla estruturada.

\sphinxlineitem{Return type}
\sphinxAtStartPar
\DUrole{sphinx_autodoc_typehints-type}{\sphinxcode{\sphinxupquote{Tuple}}{[}\sphinxcode{\sphinxupquote{bool}}, \sphinxcode{\sphinxupquote{str}}, \sphinxcode{\sphinxupquote{str}}{]}}

\end{description}\end{quote}

\end{fulllineitems}


\sphinxstepscope


\section{simple\_lock module}
\label{\detokenize{api/simple_lock:module-simple_lock}}\label{\detokenize{api/simple_lock:simple-lock-module}}\label{\detokenize{api/simple_lock::doc}}\index{module@\spxentry{module}!simple\_lock@\spxentry{simple\_lock}}\index{simple\_lock@\spxentry{simple\_lock}!module@\spxentry{module}}\index{exe\_dir\_base() (in module simple\_lock)@\spxentry{exe\_dir\_base()}\spxextra{in module simple\_lock}}

\begin{fulllineitems}
\phantomsection\label{\detokenize{api/simple_lock:simple_lock.exe_dir_base}}
\pysigstartsignatures
\pysiglinewithargsret
{\sphinxcode{\sphinxupquote{simple\_lock.}}\sphinxbfcode{\sphinxupquote{exe\_dir\_base}}}
{}
{}
\pysigstopsignatures
\sphinxAtStartPar
Retorna o diretório base do executável, considerando o ambiente de execução (PyInstaller ou desenvolvimento). O método verifica se o atributo sys.\_MEIPASS está presente, o que indica que o aplicativo está sendo executado a partir de um pacote criado pelo PyInstaller, e retorna esse diretório como base. Caso contrário, ele retorna o diretório do arquivo atual como base. O resultado é o caminho do diretório onde o aplicativo está sendo executado, que pode ser usado para localizar recursos e arquivos relacionados ao aplicativo.
:returns: O diretório base do executável, considerando o ambiente de execução (PyInstaller ou desenvolvimento). O método verifica se o atributo sys.\_MEIPASS está presente, o que indica que o aplicativo está sendo executado a partir de um pacote criado pelo PyInstaller, e retorna esse diretório como base. Caso contrário, ele retorna o diretório do arquivo atual como base. O resultado é o caminho do diretório onde o aplicativo está sendo executado, que pode ser usado para localizar recursos e arquivos relacionados ao aplicativo.
\begin{quote}\begin{description}
\sphinxlineitem{Return type}
\sphinxAtStartPar
\DUrole{sphinx_autodoc_typehints-type}{\sphinxcode{\sphinxupquote{str}}}

\end{description}\end{quote}

\end{fulllineitems}

\index{internal\_dir() (in module simple\_lock)@\spxentry{internal\_dir()}\spxextra{in module simple\_lock}}

\begin{fulllineitems}
\phantomsection\label{\detokenize{api/simple_lock:simple_lock.internal_dir}}
\pysigstartsignatures
\pysiglinewithargsret
{\sphinxcode{\sphinxupquote{simple\_lock.}}\sphinxbfcode{\sphinxupquote{internal\_dir}}}
{}
{}
\pysigstopsignatures
\sphinxAtStartPar
Retorna o caminho do diretório interno do aplicativo, que é uma subpasta chamada “\_internal” localizada dentro do diretório base do executável. O método utiliza a função exe\_dir\_base para obter o diretório base do executável e, em seguida, combina esse diretório com o nome da subpasta “\_internal” para formar o caminho completo do diretório interno. Esse diretório é usado para armazenar arquivos relacionados ao aplicativo que não devem ser acessados ou modificados diretamente pelo usuário, como o marcador de licença e outros arquivos de configuração ou estado.
:returns: O caminho do diretório interno do aplicativo, que é uma subpasta chamada “\_internal” localizada dentro do diretório base do executável. O método utiliza a função exe\_dir\_base para obter o diretório base do executável e, em seguida, combina esse diretório com o nome da subpasta “\_internal” para formar o caminho completo do diretório interno. Esse diretório é usado para armazenar arquivos relacionados ao aplicativo que não devem ser acessados ou modificados diretamente pelo usuário, como o marcador de licença e outros arquivos de configuração ou estado.
\begin{quote}\begin{description}
\sphinxlineitem{Return type}
\sphinxAtStartPar
\DUrole{sphinx_autodoc_typehints-type}{\sphinxcode{\sphinxupquote{str}}}

\end{description}\end{quote}

\end{fulllineitems}

\index{internal\_marker\_path() (in module simple\_lock)@\spxentry{internal\_marker\_path()}\spxextra{in module simple\_lock}}

\begin{fulllineitems}
\phantomsection\label{\detokenize{api/simple_lock:simple_lock.internal_marker_path}}
\pysigstartsignatures
\pysiglinewithargsret
{\sphinxcode{\sphinxupquote{simple\_lock.}}\sphinxbfcode{\sphinxupquote{internal\_marker\_path}}}
{}
{}
\pysigstopsignatures
\sphinxAtStartPar
Retorna o caminho do arquivo de marcador interno, que é um arquivo chamado “\_internal.dat” localizado dentro do diretório interno do aplicativo. O método utiliza a função internal\_dir para obter o caminho do diretório interno e, em seguida, combina esse diretório com o nome do arquivo de marcador para formar o caminho completo do arquivo de marcador. Esse arquivo é criado pelo instalador e serve como um indicador de que o aplicativo foi instalado corretamente, além de ser usado para armazenar informações relacionadas à licença e ao estado do aplicativo.
:returns: O caminho do arquivo de marcador interno, que é um arquivo chamado “\_internal.dat” localizado dentro do diretório interno do aplicativo. O método utiliza a função internal\_dir para obter o caminho do diretório interno e, em seguida, combina esse diretório com o nome do arquivo de marcador para formar o caminho completo do arquivo de marcador. Esse arquivo é criado pelo instalador e serve como um indicador de que o aplicativo foi instalado corretamente, além de ser usado para armazenar informações relacionadas à licença e ao estado do aplicativo.
\begin{quote}\begin{description}
\sphinxlineitem{Return type}
\sphinxAtStartPar
\DUrole{sphinx_autodoc_typehints-type}{\sphinxcode{\sphinxupquote{str}}}

\end{description}\end{quote}

\end{fulllineitems}

\index{programdata\_dir() (in module simple\_lock)@\spxentry{programdata\_dir()}\spxextra{in module simple\_lock}}

\begin{fulllineitems}
\phantomsection\label{\detokenize{api/simple_lock:simple_lock.programdata_dir}}
\pysigstartsignatures
\pysiglinewithargsret
{\sphinxcode{\sphinxupquote{simple\_lock.}}\sphinxbfcode{\sphinxupquote{programdata\_dir}}}
{}
{}
\pysigstopsignatures
\sphinxAtStartPar
Retorna o caminho do diretório ProgramData específico para o aplicativo, que é uma subpasta com o nome do aplicativo localizada dentro do diretório ProgramData do sistema. O método obtém o caminho do diretório ProgramData usando a variável de ambiente “PROGRAMDATA” (ou um valor padrão se a variável não estiver definida), e então combina esse caminho com o nome do aplicativo para formar o caminho completo do diretório específico do aplicativo dentro do ProgramData. O método também garante que esse diretório exista, criando\sphinxhyphen{}o se necessário, para que possa ser usado para armazenar arquivos relacionados ao estado e à configuração do aplicativo.
:returns: O caminho do diretório ProgramData específico para o aplicativo, que é uma subpasta com o nome do aplicativo localizada dentro do diretório ProgramData do sistema. O método obtém o caminho do diretório ProgramData usando a variável de ambiente “PROGRAMDATA” (ou um valor padrão se a variável não estiver definida), e então combina esse caminho com o nome do aplicativo para formar o caminho completo do diretório específico do aplicativo dentro do ProgramData. O método também garante que esse diretório exista, criando\sphinxhyphen{}o se necessário, para que possa ser usado para armazenar arquivos relacionados ao estado e à configuração do aplicativo.
\begin{quote}\begin{description}
\sphinxlineitem{Return type}
\sphinxAtStartPar
\DUrole{sphinx_autodoc_typehints-type}{\sphinxcode{\sphinxupquote{str}}}

\end{description}\end{quote}

\end{fulllineitems}

\index{license\_path\_programdata() (in module simple\_lock)@\spxentry{license\_path\_programdata()}\spxextra{in module simple\_lock}}

\begin{fulllineitems}
\phantomsection\label{\detokenize{api/simple_lock:simple_lock.license_path_programdata}}
\pysigstartsignatures
\pysiglinewithargsret
{\sphinxcode{\sphinxupquote{simple\_lock.}}\sphinxbfcode{\sphinxupquote{license\_path\_programdata}}}
{}
{}
\pysigstopsignatures
\sphinxAtStartPar
Retorna o caminho do arquivo de licença localizado no diretório ProgramData específico para o aplicativo. O método utiliza a função programdata\_dir para obter o caminho do diretório específico do aplicativo dentro do ProgramData, e então combina esse diretório com o nome do arquivo de licença (definido pela constante PROGRAMDATA\_FILE) para formar o caminho completo do arquivo de licença. Esse arquivo é usado para armazenar informações relacionadas à licença do aplicativo, e o método garante que o caminho retornado seja válido para uso na leitura e escrita de dados relacionados à licença.
:returns: O caminho completo do arquivo de licença localizado no diretório ProgramData específico para o aplicativo.
\begin{quote}\begin{description}
\sphinxlineitem{Return type}
\sphinxAtStartPar
\DUrole{sphinx_autodoc_typehints-type}{\sphinxcode{\sphinxupquote{str}}}

\end{description}\end{quote}

\end{fulllineitems}

\index{machine\_id() (in module simple\_lock)@\spxentry{machine\_id()}\spxextra{in module simple\_lock}}

\begin{fulllineitems}
\phantomsection\label{\detokenize{api/simple_lock:simple_lock.machine_id}}
\pysigstartsignatures
\pysiglinewithargsret
{\sphinxcode{\sphinxupquote{simple\_lock.}}\sphinxbfcode{\sphinxupquote{machine\_id}}}
{}
{}
\pysigstopsignatures
\sphinxAtStartPar
Calcula o hash SHA\sphinxhyphen{}256 do MachineGuid do registro do Windows e retorna como uma string em letras maiúsculas. O método chama a função \_machine\_guid para obter o MachineGuid da máquina e, em seguida, calcula o hash SHA\sphinxhyphen{}256 desse valor usando a biblioteca hashlib. O resultado é convertido para uma string hexadecimal e transformado em letras maiúsculas antes de ser retornado. Esse hash é usado como um identificador único para a máquina, permitindo que o aplicativo associe informações ou licenças a um dispositivo específico de forma segura.
:returns: O hash SHA\sphinxhyphen{}256 do MachineGuid do registro do Windows, retornado como uma string em letras maiúsculas. O método chama a função \_machine\_guid para obter o MachineGuid da máquina e, em seguida, calcula o hash SHA\sphinxhyphen{}256 desse valor usando a biblioteca hashlib. O resultado é convertido para uma string hexadecimal e transformado em letras maiúsculas antes de ser retornado. Esse hash é usado como um identificador único para a máquina, permitindo que o aplicativo associe informações ou licenças a um dispositivo específico de forma segura.
\begin{quote}\begin{description}
\sphinxlineitem{Return type}
\sphinxAtStartPar
\DUrole{sphinx_autodoc_typehints-type}{\sphinxcode{\sphinxupquote{str}}}

\end{description}\end{quote}

\end{fulllineitems}

\index{machine\_id\_hash() (in module simple\_lock)@\spxentry{machine\_id\_hash()}\spxextra{in module simple\_lock}}

\begin{fulllineitems}
\phantomsection\label{\detokenize{api/simple_lock:simple_lock.machine_id_hash}}
\pysigstartsignatures
\pysiglinewithargsret
{\sphinxcode{\sphinxupquote{simple\_lock.}}\sphinxbfcode{\sphinxupquote{machine\_id\_hash}}}
{\sphinxparam{\DUrole{n}{mid}}}
{}
\pysigstopsignatures
\sphinxAtStartPar
Calcula o hash SHA\sphinxhyphen{}256 de um identificador de máquina (Machine ID) fornecido como argumento e retorna o resultado como uma string em letras maiúsculas. O método recebe um valor de identificador de máquina (mid) como entrada, codifica esse valor em UTF\sphinxhyphen{}8 e calcula o hash SHA\sphinxhyphen{}256 usando a biblioteca hashlib. O resultado é convertido para uma string hexadecimal e transformado em letras maiúsculas antes de ser retornado. Esse hash é usado para criar um identificador único e seguro para a máquina, permitindo que o aplicativo associe informações ou licenças a um dispositivo específico de forma segura.
:param mid: O valor do identificador da máquina (Machine ID) a ser convertido em hash SHA\sphinxhyphen{}256.
:type mid: \DUrole{sphinx_autodoc_typehints-type}{\sphinxcode{\sphinxupquote{str}}}
:returns: O hash SHA\sphinxhyphen{}256 do valor passado como argumento \sphinxtitleref{mid}, retornado como uma string em letras maiúsculas.
\begin{quote}\begin{description}
\sphinxlineitem{Return type}
\sphinxAtStartPar
\DUrole{sphinx_autodoc_typehints-type}{\sphinxcode{\sphinxupquote{str}}}

\end{description}\end{quote}

\end{fulllineitems}

\index{DATA\_BLOB (class in simple\_lock)@\spxentry{DATA\_BLOB}\spxextra{class in simple\_lock}}

\begin{fulllineitems}
\phantomsection\label{\detokenize{api/simple_lock:simple_lock.DATA_BLOB}}
\pysigstartsignatures
\pysigline
{\sphinxbfcode{\sphinxupquote{\DUrole{k}{class}\DUrole{w}{ }}}\sphinxcode{\sphinxupquote{simple\_lock.}}\sphinxbfcode{\sphinxupquote{DATA\_BLOB}}}
\pysigstopsignatures
\sphinxAtStartPar
Bases: \sphinxcode{\sphinxupquote{Structure}}

\sphinxAtStartPar
Representa a estrutura DATA\_BLOB usada pela API de criptografia do Windows (DPAPI). Essa estrutura é usada para armazenar dados binários que serão protegidos ou desprotegidos usando as funções CryptProtectData e CryptUnprotectData. A estrutura contém dois campos: cbData, que é um valor DWORD que indica o tamanho dos dados em bytes, e pbData, que é um ponteiro para os dados binários a serem protegidos ou desprotegidos. Essa estrutura é essencial para a interação com a API de criptografia do Windows, permitindo que os dados sejam manipulados de forma segura e eficiente.
:param ctypes: A classe DATA\_BLOB herda de ctypes.Structure, o que permite que ela seja usada para definir uma estrutura de dados compatível com a API de criptografia do Windows. A estrutura é definida usando a sintaxe de classes do Python, e os campos da estrutura são especificados na lista \_fields\_, onde cada campo é representado por uma tupla contendo o nome do campo e seu tipo de dados correspondente (por exemplo, wt.DWORD para cbData e ctypes.POINTER(ctypes.c\_byte) para pbData). Essa definição permite que a estrutura seja usada para armazenar e manipular dados binários de forma segura ao interagir com as funções de criptografia do Windows.
:type ctypes: \_type\_
\index{cbData (simple\_lock.DATA\_BLOB attribute)@\spxentry{cbData}\spxextra{simple\_lock.DATA\_BLOB attribute}}

\begin{fulllineitems}
\phantomsection\label{\detokenize{api/simple_lock:simple_lock.DATA_BLOB.cbData}}
\pysigstartsignatures
\pysigline
{\sphinxbfcode{\sphinxupquote{cbData}}}
\pysigstopsignatures
\sphinxAtStartPar
Structure/Union member

\end{fulllineitems}

\index{pbData (simple\_lock.DATA\_BLOB attribute)@\spxentry{pbData}\spxextra{simple\_lock.DATA\_BLOB attribute}}

\begin{fulllineitems}
\phantomsection\label{\detokenize{api/simple_lock:simple_lock.DATA_BLOB.pbData}}
\pysigstartsignatures
\pysigline
{\sphinxbfcode{\sphinxupquote{pbData}}}
\pysigstopsignatures
\sphinxAtStartPar
Structure/Union member

\end{fulllineitems}


\end{fulllineitems}

\index{dpapi\_protect\_localmachine() (in module simple\_lock)@\spxentry{dpapi\_protect\_localmachine()}\spxextra{in module simple\_lock}}

\begin{fulllineitems}
\phantomsection\label{\detokenize{api/simple_lock:simple_lock.dpapi_protect_localmachine}}
\pysigstartsignatures
\pysiglinewithargsret
{\sphinxcode{\sphinxupquote{simple\_lock.}}\sphinxbfcode{\sphinxupquote{dpapi\_protect\_localmachine}}}
{\sphinxparam{\DUrole{n}{plaintext}}}
{}
\pysigstopsignatures
\sphinxAtStartPar
Protege um objeto de bytes usando a API de criptografia do Windows (DPAPI) com a opção CRYPTPROTECT\_LOCAL\_MACHINE, que vincula os dados protegidos à máquina em que foram protegidos. O método recebe um objeto de bytes como entrada, converte esses bytes em uma estrutura DATA\_BLOB usando a função \_bytes\_to\_blob, e então chama a função CryptProtectData para proteger os dados. O resultado é um novo objeto de bytes que representa os dados protegidos, que pode ser armazenado ou transmitido com segurança. Se ocorrer um erro durante o processo de proteção, o método levanta uma exceção OSError com o código de erro correspondente.
:param plaintext:
:type plaintext: \DUrole{sphinx_autodoc_typehints-type}{\sphinxcode{\sphinxupquote{bytes}}}
:raises OSError: Se ocorrer um erro durante o processo de proteção dos dados usando a API de criptografia do Windows, o método levanta uma exceção OSError com o código de erro correspondente, que pode ser obtido usando ctypes.get\_last\_error() para fornecer informações sobre a natureza do erro ocorrido.
\begin{quote}\begin{description}
\sphinxlineitem{Returns}
\sphinxAtStartPar
Um novo objeto de bytes que representa os dados protegidos usando a API de criptografia do Windows (DPAPI) com a opção CRYPTPROTECT\_LOCAL\_MACHINE. O método converte os bytes de entrada em uma estrutura DATA\_BLOB, chama a função CryptProtectData para proteger os dados, e retorna o resultado como um objeto de bytes. Se ocorrer um erro durante o processo de proteção, o método levanta uma exceção OSError com o código de erro correspondente.

\sphinxlineitem{Return type}
\sphinxAtStartPar
\DUrole{sphinx_autodoc_typehints-type}{\sphinxcode{\sphinxupquote{bytes}}}

\end{description}\end{quote}

\end{fulllineitems}

\index{dpapi\_unprotect() (in module simple\_lock)@\spxentry{dpapi\_unprotect()}\spxextra{in module simple\_lock}}

\begin{fulllineitems}
\phantomsection\label{\detokenize{api/simple_lock:simple_lock.dpapi_unprotect}}
\pysigstartsignatures
\pysiglinewithargsret
{\sphinxcode{\sphinxupquote{simple\_lock.}}\sphinxbfcode{\sphinxupquote{dpapi\_unprotect}}}
{\sphinxparam{\DUrole{n}{ciphertext}}}
{}
\pysigstopsignatures
\sphinxAtStartPar
Desprotege um objeto de bytes protegido usando a API de criptografia do Windows (DPAPI). O método recebe um objeto de bytes protegido como entrada, converte esses bytes em uma estrutura DATA\_BLOB usando a função \_bytes\_to\_blob, e então chama a função CryptUnprotectData para desproteger os dados. O resultado é um novo objeto de bytes que representa os dados desprotegidos, que pode ser usado ou processado conforme necessário. Se ocorrer um erro durante o processo de desproteção, o método levanta uma exceção OSError com o código de erro correspondente.
:param ciphertext: Um objeto de bytes protegido usando a API de criptografia do Windows (DPAPI).
:type ciphertext: \DUrole{sphinx_autodoc_typehints-type}{\sphinxcode{\sphinxupquote{bytes}}}
:raises OSError: Se ocorrer um erro durante o processo de desproteção dos dados usando a API de criptografia do Windows, o método levanta uma exceção OSError com o código de erro correspondente, que pode ser obtido usando ctypes.get\_last\_error() para fornecer informações sobre a natureza do erro ocorrido.
\begin{quote}\begin{description}
\sphinxlineitem{Returns}
\sphinxAtStartPar
Um novo objeto de bytes que representa os dados desprotegidos usando a API de criptografia do Windows (DPAPI). O método recebe um objeto de bytes protegido como entrada, converte esses bytes em uma estrutura DATA\_BLOB usando a função \_bytes\_to\_blob, chama a função CryptUnprotectData para desproteger os dados, e retorna o resultado como um objeto de bytes. Se ocorrer um erro durante o processo de desproteção, o método levanta uma exceção OSError com o código de erro correspondente.

\sphinxlineitem{Return type}
\sphinxAtStartPar
\DUrole{sphinx_autodoc_typehints-type}{\sphinxcode{\sphinxupquote{bytes}}}

\end{description}\end{quote}

\end{fulllineitems}

\index{ensure\_or\_mark() (in module simple\_lock)@\spxentry{ensure\_or\_mark()}\spxextra{in module simple\_lock}}

\begin{fulllineitems}
\phantomsection\label{\detokenize{api/simple_lock:simple_lock.ensure_or_mark}}
\pysigstartsignatures
\pysiglinewithargsret
{\sphinxcode{\sphinxupquote{simple\_lock.}}\sphinxbfcode{\sphinxupquote{ensure\_or\_mark}}}
{\sphinxparam{\DUrole{n}{allow\_create}\DUrole{o}{=}\DUrole{default_value}{False}}}
{}
\pysigstopsignatures
\sphinxAtStartPar
Política final (conforme combinado):
\begin{itemize}
\item {} 
\sphinxAtStartPar
O instalador cria: \sphinxcode{\sphinxupquote{\textless{}pasta app\textgreater{}\textbackslash{}\_internal\textbackslash{}\_internal.dat}} (marker somente leitura)

\item {} 
\sphinxAtStartPar
O app NUNCA escreve em \sphinxcode{\sphinxupquote{\_internal}}

\item {} 
\sphinxAtStartPar
O marcador por máquina fica em ProgramData e é DPAPI (LocalMachine):

\end{itemize}

\sphinxAtStartPar
\sphinxcode{\sphinxupquote{C:\textbackslash{}ProgramData\textbackslash{}LeadsApp\textbackslash{}cache\_string.dat}}

\sphinxAtStartPar
Regras
\begin{enumerate}
\sphinxsetlistlabels{\arabic}{enumi}{enumii}{}{.}%
\item {} 
\sphinxAtStartPar
Se marker do instalador não existir =\textgreater{} bloqueia (evita rodar app copiado solto).

\end{enumerate}

\sphinxAtStartPar
2. Se o arquivo de ProgramData NÃO existir:
\sphinxhyphen{} \sphinxcode{\sphinxupquote{allow\_create=False}} =\textgreater{} bloqueia pedindo ativação online primeiro.
\sphinxhyphen{} \sphinxcode{\sphinxupquote{allow\_create=True}}  =\textgreater{} cria o marcador por máquina (use somente após validação online OK).
3. Se o arquivo existir =\textgreater{} valida DPAPI + HMAC + machine\_id\_hash.
\begin{quote}\begin{description}
\sphinxlineitem{Parameters}
\sphinxAtStartPar
\sphinxstyleliteralstrong{\sphinxupquote{allow\_create}} (\DUrole{sphinx_autodoc_typehints-type}{\sphinxcode{\sphinxupquote{bool}}}) \textendash{} Indica se o método deve criar um novo marcador por máquina em ProgramData caso o arquivo
de marcador não exista. Se \sphinxcode{\sphinxupquote{False}}, bloqueia e solicita ativação online. Se \sphinxcode{\sphinxupquote{True}},
cria o marcador por máquina usando DPAPI (use somente após validação online bem\sphinxhyphen{}sucedida).

\sphinxlineitem{Return type}
\sphinxAtStartPar
\DUrole{sphinx_autodoc_typehints-type}{\sphinxcode{\sphinxupquote{Tuple}}{[}\sphinxcode{\sphinxupquote{bool}}, \sphinxcode{\sphinxupquote{str}}{]}}

\end{description}\end{quote}

\end{fulllineitems}


\sphinxstepscope


\section{user\_login module}
\label{\detokenize{api/user_login:module-user_login}}\label{\detokenize{api/user_login:user-login-module}}\label{\detokenize{api/user_login::doc}}\index{module@\spxentry{module}!user\_login@\spxentry{user\_login}}\index{user\_login@\spxentry{user\_login}!module@\spxentry{module}}\index{app\_state\_dir() (in module user\_login)@\spxentry{app\_state\_dir()}\spxextra{in module user\_login}}

\begin{fulllineitems}
\phantomsection\label{\detokenize{api/user_login:user_login.app_state_dir}}
\pysigstartsignatures
\pysiglinewithargsret
{\sphinxcode{\sphinxupquote{user\_login.}}\sphinxbfcode{\sphinxupquote{app\_state\_dir}}}
{}
{}
\pysigstopsignatures
\sphinxAtStartPar
Diretório de estado com fallback real:
1) \%APPDATA\% (Roaming)
2) \%LOCALAPPDATA\% (Local)
3) \%PROGRAMDATA\%LeadsAppuser\_rootLeadsApp
\begin{quote}\begin{description}
\sphinxlineitem{Returns}
\sphinxAtStartPar
O caminho do diretório de estado do aplicativo, determinado por meio de uma série de tentativas. O método verifica primeiro o diretório de roaming (\%APPDATA\%), seguido pelo diretório local (\%LOCALAPPDATA\%), e finalmente tenta criar um diretório específico para o aplicativo dentro do diretório de dados do programa (\%PROGRAMDATA\%). Se todas as tentativas falharem, o método retorna um caminho de fallback dentro do diretório do usuário.

\sphinxlineitem{Return type}
\sphinxAtStartPar
\DUrole{sphinx_autodoc_typehints-type}{\sphinxcode{\sphinxupquote{str}}}

\end{description}\end{quote}

\end{fulllineitems}

\index{login\_path() (in module user\_login)@\spxentry{login\_path()}\spxextra{in module user\_login}}

\begin{fulllineitems}
\phantomsection\label{\detokenize{api/user_login:user_login.login_path}}
\pysigstartsignatures
\pysiglinewithargsret
{\sphinxcode{\sphinxupquote{user\_login.}}\sphinxbfcode{\sphinxupquote{login\_path}}}
{}
{}
\pysigstopsignatures
\sphinxAtStartPar
Retorna o caminho completo do arquivo de login do aplicativo. O método constrói o caminho do arquivo de login combinando o diretório de estado do aplicativo, obtido por meio da função app\_state\_dir(), com o nome do arquivo de login definido pela constante LOGIN\_FILE. O resultado é o caminho completo onde as informações de login do usuário serão armazenadas ou recuperadas.
:returns: O caminho completo do arquivo de login do aplicativo.
\begin{quote}\begin{description}
\sphinxlineitem{Return type}
\sphinxAtStartPar
\DUrole{sphinx_autodoc_typehints-type}{\sphinxcode{\sphinxupquote{str}}}

\end{description}\end{quote}

\end{fulllineitems}

\index{load\_login() (in module user\_login)@\spxentry{load\_login()}\spxextra{in module user\_login}}

\begin{fulllineitems}
\phantomsection\label{\detokenize{api/user_login:user_login.load_login}}
\pysigstartsignatures
\pysiglinewithargsret
{\sphinxcode{\sphinxupquote{user\_login.}}\sphinxbfcode{\sphinxupquote{load\_login}}}
{}
{}
\pysigstopsignatures
\sphinxAtStartPar
Carrega o e\sphinxhyphen{}mail e o timestamp do login salvo.
\begin{quote}\begin{description}
\sphinxlineitem{Returns}
\sphinxAtStartPar
Uma tupla contendo o e\sphinxhyphen{}mail salvo (ou None se não houver) e o timestamp do login (ou None se não houver).

\sphinxlineitem{Return type}
\sphinxAtStartPar
\DUrole{sphinx_autodoc_typehints-type}{\sphinxcode{\sphinxupquote{Tuple}}{[}\sphinxcode{\sphinxupquote{Optional}}{[}\sphinxcode{\sphinxupquote{str}}{]}, \sphinxcode{\sphinxupquote{Optional}}{[}\sphinxcode{\sphinxupquote{int}}{]}{]}}

\end{description}\end{quote}

\end{fulllineitems}

\index{save\_login() (in module user\_login)@\spxentry{save\_login()}\spxextra{in module user\_login}}

\begin{fulllineitems}
\phantomsection\label{\detokenize{api/user_login:user_login.save_login}}
\pysigstartsignatures
\pysiglinewithargsret
{\sphinxcode{\sphinxupquote{user\_login.}}\sphinxbfcode{\sphinxupquote{save\_login}}}
{\sphinxparam{\DUrole{n}{email}}}
{}
\pysigstopsignatures
\sphinxAtStartPar
Salva o e\sphinxhyphen{}mail do login com um timestamp atual.
\begin{quote}\begin{description}
\sphinxlineitem{Parameters}
\sphinxAtStartPar
\sphinxstyleliteralstrong{\sphinxupquote{email}} (\DUrole{sphinx_autodoc_typehints-type}{\sphinxcode{\sphinxupquote{str}}}) \textendash{} O e\sphinxhyphen{}mail a ser salvo. O método irá limpar espaços em branco e converter o e\sphinxhyphen{}mail para letras minúsculas antes de salvar. O e\sphinxhyphen{}mail é armazenado junto com um timestamp que indica quando o login foi salvo, permitindo que o aplicativo determine se o login é válido ou expirado com base no tempo decorrido desde o último login.

\sphinxlineitem{Return type}
\sphinxAtStartPar
\DUrole{sphinx_autodoc_typehints-type}{\sphinxcode{\sphinxupquote{None}}}

\end{description}\end{quote}

\end{fulllineitems}

\index{clear\_login() (in module user\_login)@\spxentry{clear\_login()}\spxextra{in module user\_login}}

\begin{fulllineitems}
\phantomsection\label{\detokenize{api/user_login:user_login.clear_login}}
\pysigstartsignatures
\pysiglinewithargsret
{\sphinxcode{\sphinxupquote{user\_login.}}\sphinxbfcode{\sphinxupquote{clear\_login}}}
{}
{}
\pysigstopsignatures
\sphinxAtStartPar
Limpa o login salvo removendo o arquivo de login. O método tenta excluir o arquivo de login do sistema, que contém as informações de e\sphinxhyphen{}mail e timestamp do login. Se o arquivo existir, ele é removido, efetivamente limpando o login salvo. Se ocorrer qualquer erro durante a tentativa de remoção do arquivo (por exemplo, se o arquivo não existir ou se houver problemas de permissão), o método captura a exceção e continua sem bloquear o fluxo do programa, garantindo que a limpeza do login seja realizada de forma segura sem causar interrupções.
\begin{quote}\begin{description}
\sphinxlineitem{Return type}
\sphinxAtStartPar
\DUrole{sphinx_autodoc_typehints-type}{\sphinxcode{\sphinxupquote{None}}}

\end{description}\end{quote}

\end{fulllineitems}

\index{is\_login\_expired() (in module user\_login)@\spxentry{is\_login\_expired()}\spxextra{in module user\_login}}

\begin{fulllineitems}
\phantomsection\label{\detokenize{api/user_login:user_login.is_login_expired}}
\pysigstartsignatures
\pysiglinewithargsret
{\sphinxcode{\sphinxupquote{user\_login.}}\sphinxbfcode{\sphinxupquote{is\_login\_expired}}}
{\sphinxparam{\DUrole{n}{ts}}}
{}
\pysigstopsignatures
\sphinxAtStartPar
Verifica se o login é expirado com base no timestamp fornecido. O método compara o timestamp do login com o timestamp atual para determinar se o tempo decorrido desde o último login excede o limite definido por LOGIN\_MAX\_AGE\_SECONDS. Se o timestamp fornecido for None, o método considera o login como expirado. Caso contrário, ele calcula a diferença entre o timestamp atual e o timestamp do login e verifica se essa diferença é maior que o limite de idade máxima do login.
:param ts: O timestamp do login a ser verificado. Este valor representa o momento em que o login foi salvo e é usado para determinar se o login é válido ou expirado com base no tempo decorrido desde então.
:type ts: \DUrole{sphinx_autodoc_typehints-type}{\sphinxcode{\sphinxupquote{Optional}}{[}\sphinxcode{\sphinxupquote{int}}{]}}
:returns: Retorna True se o login for considerado expirado (ou seja, se o timestamp for None ou se o tempo decorrido desde o login exceder o limite definido por LOGIN\_MAX\_AGE\_SECONDS), e False caso contrário (ou seja, se o login ainda for válido).
\begin{quote}\begin{description}
\sphinxlineitem{Return type}
\sphinxAtStartPar
\DUrole{sphinx_autodoc_typehints-type}{\sphinxcode{\sphinxupquote{bool}}}

\end{description}\end{quote}

\end{fulllineitems}

\index{prompt\_email() (in module user\_login)@\spxentry{prompt\_email()}\spxextra{in module user\_login}}

\begin{fulllineitems}
\phantomsection\label{\detokenize{api/user_login:user_login.prompt_email}}
\pysigstartsignatures
\pysiglinewithargsret
{\sphinxcode{\sphinxupquote{user\_login.}}\sphinxbfcode{\sphinxupquote{prompt\_email}}}
{\sphinxparam{\DUrole{n}{force}\DUrole{o}{=}\DUrole{default_value}{False}}}
{}
\pysigstopsignatures
\sphinxAtStartPar
Exibe um diálogo para o usuário inserir seu e\sphinxhyphen{}mail de login. O método verifica se há um e\sphinxhyphen{}mail salvo anteriormente e se ele ainda é válido (não expirado). Se houver um e\sphinxhyphen{}mail válido salvo, ele é retornado imediatamente, a menos que o parâmetro force seja definido como True, o que força o usuário a inserir um novo e\sphinxhyphen{}mail. Se não houver um e\sphinxhyphen{}mail válido salvo ou se force for True, o método exibe um diálogo de entrada para o usuário digitar seu e\sphinxhyphen{}mail. O método valida o formato do e\sphinxhyphen{}mail inserido usando uma expressão regular e continua solicitando até que um e\sphinxhyphen{}mail válido seja fornecido ou até que o usuário cancele a operação. Se um e\sphinxhyphen{}mail válido for inserido, ele é salvo para uso futuro e retornado pelo método. Se o usuário cancelar a operação, o método retorna None.
:param force: Se True, força o usuário a inserir um novo e\sphinxhyphen{}mail mesmo que haja um e\sphinxhyphen{}mail salvo válido. Defaults to False.
:type force: \DUrole{sphinx_autodoc_typehints-type}{\sphinxcode{\sphinxupquote{bool}}}
:returns: O e\sphinxhyphen{}mail inserido pelo usuário se for válido, ou o e\sphinxhyphen{}mail salvo anteriormente se ainda for válido, ou None se o usuário cancelar a operação.
\begin{quote}\begin{description}
\sphinxlineitem{Return type}
\sphinxAtStartPar
\DUrole{sphinx_autodoc_typehints-type}{\sphinxcode{\sphinxupquote{Optional}}{[}\sphinxcode{\sphinxupquote{str}}{]}}

\end{description}\end{quote}

\end{fulllineitems}


\sphinxstepscope


\section{widgets module}
\label{\detokenize{api/widgets:module-widgets}}\label{\detokenize{api/widgets:widgets-module}}\label{\detokenize{api/widgets::doc}}\index{module@\spxentry{module}!widgets@\spxentry{widgets}}\index{widgets@\spxentry{widgets}!module@\spxentry{module}}\index{EmailInputWidget (class in widgets)@\spxentry{EmailInputWidget}\spxextra{class in widgets}}

\begin{fulllineitems}
\phantomsection\label{\detokenize{api/widgets:widgets.EmailInputWidget}}
\pysigstartsignatures
\pysiglinewithargsret
{\sphinxbfcode{\sphinxupquote{\DUrole{k}{class}\DUrole{w}{ }}}\sphinxcode{\sphinxupquote{widgets.}}\sphinxbfcode{\sphinxupquote{EmailInputWidget}}}
{\sphinxparam{\DUrole{n}{domains}}\sphinxparamcomma \sphinxparam{\DUrole{n}{parent}\DUrole{o}{=}\DUrole{default_value}{None}}}
{}
\pysigstopsignatures
\sphinxAtStartPar
Bases: \sphinxcode{\sphinxupquote{QWidget}}

\sphinxAtStartPar
Encapsula:
\sphinxhyphen{} QLineEdit (parte local) + QComboBox (domínio)
\sphinxhyphen{} Monta e lê email final
\sphinxhyphen{} Atualiza lista de domínios preservando seleção
\begin{itemize}
\item {} 
\sphinxAtStartPar
Se o local contém ‘@’, considera colado completo e retorna ok=True

\item {} 
\sphinxAtStartPar
Se local não vazio e domínio vazio =\textgreater{} ok=False

\item {} 
\sphinxAtStartPar
Se local vazio =\textgreater{} retorna “” e ok=True (campo em branco)

\end{itemize}
\index{set\_domains() (widgets.EmailInputWidget method)@\spxentry{set\_domains()}\spxextra{widgets.EmailInputWidget method}}

\begin{fulllineitems}
\phantomsection\label{\detokenize{api/widgets:widgets.EmailInputWidget.set_domains}}
\pysigstartsignatures
\pysiglinewithargsret
{\sphinxbfcode{\sphinxupquote{set\_domains}}}
{\sphinxparam{\DUrole{n}{domains}}}
{}
\pysigstopsignatures\begin{quote}\begin{description}
\sphinxlineitem{Return type}
\sphinxAtStartPar
\DUrole{sphinx_autodoc_typehints-type}{\sphinxcode{\sphinxupquote{None}}}

\end{description}\end{quote}

\end{fulllineitems}

\index{local() (widgets.EmailInputWidget method)@\spxentry{local()}\spxextra{widgets.EmailInputWidget method}}

\begin{fulllineitems}
\phantomsection\label{\detokenize{api/widgets:widgets.EmailInputWidget.local}}
\pysigstartsignatures
\pysiglinewithargsret
{\sphinxbfcode{\sphinxupquote{local}}}
{}
{}
\pysigstopsignatures\begin{quote}\begin{description}
\sphinxlineitem{Return type}
\sphinxAtStartPar
\DUrole{sphinx_autodoc_typehints-type}{\sphinxcode{\sphinxupquote{str}}}

\end{description}\end{quote}

\end{fulllineitems}

\index{domain() (widgets.EmailInputWidget method)@\spxentry{domain()}\spxextra{widgets.EmailInputWidget method}}

\begin{fulllineitems}
\phantomsection\label{\detokenize{api/widgets:widgets.EmailInputWidget.domain}}
\pysigstartsignatures
\pysiglinewithargsret
{\sphinxbfcode{\sphinxupquote{domain}}}
{}
{}
\pysigstopsignatures\begin{quote}\begin{description}
\sphinxlineitem{Return type}
\sphinxAtStartPar
\DUrole{sphinx_autodoc_typehints-type}{\sphinxcode{\sphinxupquote{str}}}

\end{description}\end{quote}

\end{fulllineitems}

\index{set\_parts() (widgets.EmailInputWidget method)@\spxentry{set\_parts()}\spxextra{widgets.EmailInputWidget method}}

\begin{fulllineitems}
\phantomsection\label{\detokenize{api/widgets:widgets.EmailInputWidget.set_parts}}
\pysigstartsignatures
\pysiglinewithargsret
{\sphinxbfcode{\sphinxupquote{set\_parts}}}
{\sphinxparam{\DUrole{n}{local}}\sphinxparamcomma \sphinxparam{\DUrole{n}{domain}}}
{}
\pysigstopsignatures\begin{quote}\begin{description}
\sphinxlineitem{Return type}
\sphinxAtStartPar
\DUrole{sphinx_autodoc_typehints-type}{\sphinxcode{\sphinxupquote{None}}}

\end{description}\end{quote}

\end{fulllineitems}

\index{set\_email() (widgets.EmailInputWidget method)@\spxentry{set\_email()}\spxextra{widgets.EmailInputWidget method}}

\begin{fulllineitems}
\phantomsection\label{\detokenize{api/widgets:widgets.EmailInputWidget.set_email}}
\pysigstartsignatures
\pysiglinewithargsret
{\sphinxbfcode{\sphinxupquote{set\_email}}}
{\sphinxparam{\DUrole{n}{email}}}
{}
\pysigstopsignatures\begin{quote}\begin{description}
\sphinxlineitem{Return type}
\sphinxAtStartPar
\DUrole{sphinx_autodoc_typehints-type}{\sphinxcode{\sphinxupquote{None}}}

\end{description}\end{quote}

\end{fulllineitems}

\index{get\_email() (widgets.EmailInputWidget method)@\spxentry{get\_email()}\spxextra{widgets.EmailInputWidget method}}

\begin{fulllineitems}
\phantomsection\label{\detokenize{api/widgets:widgets.EmailInputWidget.get_email}}
\pysigstartsignatures
\pysiglinewithargsret
{\sphinxbfcode{\sphinxupquote{get\_email}}}
{}
{}
\pysigstopsignatures\begin{quote}\begin{description}
\sphinxlineitem{Return type}
\sphinxAtStartPar
\DUrole{sphinx_autodoc_typehints-type}{\sphinxcode{\sphinxupquote{Tuple}}{[}\sphinxcode{\sphinxupquote{str}}, \sphinxcode{\sphinxupquote{bool}}{]}}

\end{description}\end{quote}

\end{fulllineitems}


\end{fulllineitems}

\index{FieldRowWidget (class in widgets)@\spxentry{FieldRowWidget}\spxextra{class in widgets}}

\begin{fulllineitems}
\phantomsection\label{\detokenize{api/widgets:widgets.FieldRowWidget}}
\pysigstartsignatures
\pysiglinewithargsret
{\sphinxbfcode{\sphinxupquote{\DUrole{k}{class}\DUrole{w}{ }}}\sphinxcode{\sphinxupquote{widgets.}}\sphinxbfcode{\sphinxupquote{FieldRowWidget}}}
{\sphinxparam{\DUrole{n}{field\_id}}\sphinxparamcomma \sphinxparam{\DUrole{n}{label\_text}}\sphinxparamcomma \sphinxparam{\DUrole{n}{input\_widget}}\sphinxparamcomma \sphinxparam{\DUrole{n}{parent}\DUrole{o}{=}\DUrole{default_value}{None}}}
{}
\pysigstopsignatures
\sphinxAtStartPar
Bases: \sphinxcode{\sphinxupquote{QWidget}}
\index{editRequested (widgets.FieldRowWidget attribute)@\spxentry{editRequested}\spxextra{widgets.FieldRowWidget attribute}}

\begin{fulllineitems}
\phantomsection\label{\detokenize{api/widgets:widgets.FieldRowWidget.editRequested}}
\pysigstartsignatures
\pysigline
{\sphinxbfcode{\sphinxupquote{editRequested}}}
\pysigstopsignatures
\sphinxAtStartPar
alias of \sphinxcode{\sphinxupquote{str}}

\end{fulllineitems}

\index{deleteRequested (widgets.FieldRowWidget attribute)@\spxentry{deleteRequested}\spxextra{widgets.FieldRowWidget attribute}}

\begin{fulllineitems}
\phantomsection\label{\detokenize{api/widgets:widgets.FieldRowWidget.deleteRequested}}
\pysigstartsignatures
\pysigline
{\sphinxbfcode{\sphinxupquote{deleteRequested}}}
\pysigstopsignatures
\sphinxAtStartPar
alias of \sphinxcode{\sphinxupquote{str}}

\end{fulllineitems}

\index{lockToggled (widgets.FieldRowWidget attribute)@\spxentry{lockToggled}\spxextra{widgets.FieldRowWidget attribute}}

\begin{fulllineitems}
\phantomsection\label{\detokenize{api/widgets:widgets.FieldRowWidget.lockToggled}}
\pysigstartsignatures
\pysigline
{\sphinxbfcode{\sphinxupquote{lockToggled}}}
\pysigstopsignatures
\sphinxAtStartPar
alias of \sphinxcode{\sphinxupquote{str}}

\end{fulllineitems}

\index{set\_locked() (widgets.FieldRowWidget method)@\spxentry{set\_locked()}\spxextra{widgets.FieldRowWidget method}}

\begin{fulllineitems}
\phantomsection\label{\detokenize{api/widgets:widgets.FieldRowWidget.set_locked}}
\pysigstartsignatures
\pysiglinewithargsret
{\sphinxbfcode{\sphinxupquote{set\_locked}}}
{\sphinxparam{\DUrole{n}{locked}}}
{}
\pysigstopsignatures\begin{quote}\begin{description}
\sphinxlineitem{Return type}
\sphinxAtStartPar
\DUrole{sphinx_autodoc_typehints-type}{\sphinxcode{\sphinxupquote{None}}}

\end{description}\end{quote}

\end{fulllineitems}


\end{fulllineitems}

\index{BoolInputWidget (class in widgets)@\spxentry{BoolInputWidget}\spxextra{class in widgets}}

\begin{fulllineitems}
\phantomsection\label{\detokenize{api/widgets:widgets.BoolInputWidget}}
\pysigstartsignatures
\pysiglinewithargsret
{\sphinxbfcode{\sphinxupquote{\DUrole{k}{class}\DUrole{w}{ }}}\sphinxcode{\sphinxupquote{widgets.}}\sphinxbfcode{\sphinxupquote{BoolInputWidget}}}
{\sphinxparam{\DUrole{n}{parent}\DUrole{o}{=}\DUrole{default_value}{None}}}
{}
\pysigstopsignatures
\sphinxAtStartPar
Bases: \sphinxcode{\sphinxupquote{QWidget}}

\sphinxAtStartPar
Widget de entrada booleana baseado em QComboBox.
\index{value() (widgets.BoolInputWidget method)@\spxentry{value()}\spxextra{widgets.BoolInputWidget method}}

\begin{fulllineitems}
\phantomsection\label{\detokenize{api/widgets:widgets.BoolInputWidget.value}}
\pysigstartsignatures
\pysiglinewithargsret
{\sphinxbfcode{\sphinxupquote{value}}}
{}
{}
\pysigstopsignatures\begin{quote}\begin{description}
\sphinxlineitem{Return type}
\sphinxAtStartPar
\DUrole{sphinx_autodoc_typehints-type}{\sphinxcode{\sphinxupquote{str}}}

\end{description}\end{quote}

\end{fulllineitems}

\index{set\_value() (widgets.BoolInputWidget method)@\spxentry{set\_value()}\spxextra{widgets.BoolInputWidget method}}

\begin{fulllineitems}
\phantomsection\label{\detokenize{api/widgets:widgets.BoolInputWidget.set_value}}
\pysigstartsignatures
\pysiglinewithargsret
{\sphinxbfcode{\sphinxupquote{set\_value}}}
{\sphinxparam{\DUrole{n}{v}}}
{}
\pysigstopsignatures\begin{quote}\begin{description}
\sphinxlineitem{Return type}
\sphinxAtStartPar
\DUrole{sphinx_autodoc_typehints-type}{\sphinxcode{\sphinxupquote{None}}}

\end{description}\end{quote}

\end{fulllineitems}


\end{fulllineitems}



\renewcommand{\indexname}{Python Module Index}
\begin{sphinxtheindex}
\let\bigletter\sphinxstyleindexlettergroup
\bigletter{a}
\item\relax\sphinxstyleindexentry{app}\sphinxstyleindexpageref{api/app:\detokenize{module-app}}
\indexspace
\bigletter{c}
\item\relax\sphinxstyleindexentry{colors}\sphinxstyleindexpageref{api/colors:\detokenize{module-colors}}
\item\relax\sphinxstyleindexentry{config}\sphinxstyleindexpageref{api/config:\detokenize{module-config}}
\indexspace
\bigletter{d}
\item\relax\sphinxstyleindexentry{dialogs}\sphinxstyleindexpageref{api/dialogs:\detokenize{module-dialogs}}
\indexspace
\bigletter{e}
\item\relax\sphinxstyleindexentry{excel\_io}\sphinxstyleindexpageref{api/excel_io:\detokenize{module-excel_io}}
\indexspace
\bigletter{f}
\item\relax\sphinxstyleindexentry{formatters}\sphinxstyleindexpageref{api/formatters:\detokenize{module-formatters}}
\indexspace
\bigletter{i}
\item\relax\sphinxstyleindexentry{icon\_manager}\sphinxstyleindexpageref{api/icon_manager:\detokenize{module-icon_manager}}
\indexspace
\bigletter{l}
\item\relax\sphinxstyleindexentry{license\_dialog}\sphinxstyleindexpageref{api/license_dialog:\detokenize{module-license_dialog}}
\item\relax\sphinxstyleindexentry{loading\_screen}\sphinxstyleindexpageref{api/loading_screen:\detokenize{module-loading_screen}}
\indexspace
\bigletter{m}
\item\relax\sphinxstyleindexentry{models}\sphinxstyleindexpageref{api/models:\detokenize{module-models}}
\indexspace
\bigletter{o}
\item\relax\sphinxstyleindexentry{online\_license}\sphinxstyleindexpageref{api/online_license:\detokenize{module-online_license}}
\indexspace
\bigletter{s}
\item\relax\sphinxstyleindexentry{simple\_lock}\sphinxstyleindexpageref{api/simple_lock:\detokenize{module-simple_lock}}
\indexspace
\bigletter{u}
\item\relax\sphinxstyleindexentry{user\_login}\sphinxstyleindexpageref{api/user_login:\detokenize{module-user_login}}
\indexspace
\bigletter{w}
\item\relax\sphinxstyleindexentry{widgets}\sphinxstyleindexpageref{api/widgets:\detokenize{module-widgets}}
\end{sphinxtheindex}

\renewcommand{\indexname}{Index}
\printindex
\end{document}